% =====================================================================
%  Harald Høffding, TYPES OF RELIGIOUS THOUGHT.
%  English translation of: Religiøse Tanketyper.  København: Gyldendalske Boghandel, Nordisk Forlag, 1927.
%
%  GENERATED; DO NOT EDIT BY HAND.  Built by ebook-method/religioese-tanketyper/tr_parts.py build --write from
%  tr_template.tex and the parts in ebook-method/religioese-tanketyper/tr/ (translated 2026-10-17 from
%  texts/hoeffding/religioese-tanketyper/transcription.tex, following TRANSLATOR-BRIEF.md and GLOSSARY.md there).
%  Page numbers in the margin [N] are those of the Danish print.  Footnote marks 1) 2) as printed.
%  Title: Types of Religious Thought; Højsind = magnanimity (Hans, 2026-10-17).  Oplevelse = experience; Erfaring =
%  experience, marked [Erfaring] where the two are contrasted (Hans, 2026-10-14).
%  Pp. 27-28 are missing from the scan the transcription was made from; their text was set from the e-book (see the
%  transcription's header) and is translated from that text.
%  ODD SPOTS: the seven translators listed 61 places in the Danish that looked wrong or unclear (tr/odd-*.tsv).  31 were
%  read at the image (400-dpi line crops): 1 transcription fault (psykologisk-historisk, p. 34), 3 divisions whose hyphen
%  the scan's OCR lost, the verse on p. 125 restructured; 11 misprints of the print kept and marked in the Danish
%  (Han sidste Ord, sammme, Kirkegaard, nøvrig, Krafte ...), translated here as the sense requires.  The rest were
%  questions of grammar, reference or period spelling, settled in translating.
% =====================================================================
\documentclass[12pt,a4paper,openany]{book}
\usepackage[utf8]{inputenc}
\usepackage[T1]{fontenc}
\usepackage{libertinus}
\usepackage{textalpha}
\usepackage[english]{babel}
\usepackage{setspace}
\usepackage[margin=1.2in]{geometry}
\usepackage{fancyhdr}
\setlength{\headheight}{28pt}
\addtolength{\topmargin}{-13.5pt}
\usepackage{hyperref}
\pagestyle{fancy}
\fancyhf{}
\fancyhead[LE,RO]{\thepage}
\fancyhead[RE]{\textit{Harald Høffding: Types of Religious Thought}}
\fancyhead[LO]{\textit{\leftmark}}
\fancypagestyle{plain}{\fancyhf{}\fancyhead[LE,RO]{\thepage}}
\setstretch{1.3}
\usepackage{marginnote}
\usepackage{graphicx}
\DeclareUnicodeCharacter{0254}{\reflectbox{c}}  % ɔ: the old Danish id-est mark (a reversed c)
\newcommand{\starmark}[1]{\arabic{footnote})}   % the print numbers its notes 1) 2) per page
\renewcommand{\thefootnote}{\starmark{\value{footnote}}}
\newcommand{\opage}[1]{\marginnote{\footnotesize\textit{[#1]}}}
\newcommand{\apage}[1]{\marginnote{\footnotesize\textit{[#1]}}}
\newcommand{\whiteline}{\par\vspace{\baselineskip}\par}
% the printed Indhold: an essay or chapter line (title, page), a section line
\newcommand{\ientry}[2]{\par\smallskip\noindent #1\dotfill\,#2\par}
\newcommand{\isub}[2]{\par\noindent\hspace*{2em}#1\dotfill\,#2\par}
% the headings (no \addcontentsline: the print has its own Indhold, and a footnote in a heading breaks a toc entry)
\pdfstringdefDisableCommands{\renewcommand{\footnote}[2][]{}}
\newcommand{\prefacehead}[1]{\chapter*{#1}\markboth{Preface}{Preface}}
\newcommand{\sectionhead}[1]{\par\medskip\begin{center}#1\end{center}\smallskip}
\newcommand{\essayhead}[2][]{\clearpage\begin{center}{\Large #2}\end{center}\if\relax\detokenize{#1}\relax\markboth{#2}{#2}\else\markboth{#1}{#1}\fi\medskip}   % [running head] when the title carries a note
\newcommand{\asterism}{\par\medskip\begin{center}*\hspace{2.5em}*\\[-2pt]*\end{center}\medskip}   % the print's three stars, two above one
\newcommand{\piecerule}{\par\medskip\begin{center}\rule{0.2\textwidth}{0.4pt}\end{center}\medskip}   % the short rule at the end of a piece
\newcommand{\essaysource}[1]{\begin{center}{\small #1}\end{center}\medskip}
\newcommand{\subhead}[1]{\par\medskip\begin{center}#1\end{center}\smallskip}
\newcommand{\romanhead}[1]{\par\medskip\begin{center}#1.\end{center}\smallskip}

\begin{document}

\frontmatter
\begin{titlepage}
\centering
\vspace*{2cm}
{\large HARALD HØFFDING}\par
\vspace{1.5cm}
{\Huge TYPES OF\\[6pt] RELIGIOUS THOUGHT}\par
\vspace{1cm}
{[\textit{fleuron}]}\par
\vfill
{\small GYLDENDALSKE BOGHANDEL · NORDISK FORLAG\\ COPENHAGEN · MCMXXVII}\par
\vspace{1cm}
{\small\textit{Translated from the Danish}}\par
\end{titlepage}
\clearpage

\chapter*{CONTENTS}
\markboth{Contents}{Contents}
\opage{5}%
\noindent\hfill{\small Page}\par
\ientry{Preface (1927)}{7}
\ientry{Philosophical Points of View (1923)}{9}
\ientry{On My Work as an Author (1923)}{17}
\medskip
\ientry{Plato on Art and Religion (1924)}{41}
\ientry{Giordano Bruno (1922)}{50}
\ientry{Spinoza (1927)}{59}
\ientry{Pascal and Kierkegaard (1923)}{70}
\ientry{Sibbern and Kierkegaard (1923)}{98}
\ientry{Philosophy and Ecclesiastical Rationalism (1922)}{112}
\ientry{A Modern Buddhist (1926)}{117}
\medskip
\ientry{``Bakkehus og Solbjerg'' (1924)}{126}
\ientry{``De Vises Sten'' (1923)}{136}
\ientry{``For aaben Dør'' (1926)}{147}
\ientry{``Den grønne Mark'' (1926)}{154}
\medskip
\ientry{Shakespeare: ``Stormen'' (1926)}{158}
\clearpage

% --- p. 7 ---
\setcounter{footnote}{0}%

\opage{7}\prefacehead{PREFACE}

This book contains a series of lectures and essays
from recent years. I have felt a need to try once more to
shed light on the question of the relation between thought and faith,
which has occupied me since my youth.
In my ``Religionsfilosofi'' I have given the answer
to the question that I have been able to give. There are,
in my conviction, only three ways by which the religious
problem can be illuminated, an epistemological, a psychological,
and an ethical way, all three ways each in turn illuminated
by what the history of religion can teach us. In the book
that I present here I have, as the title indicates, brought forward
some types of thoughts about religion, as they have emerged at
different times in thinkers and poets.

As an introduction to the historical presentations I give,
in the first two pieces, a brief account of my own philosophical
conception, which naturally underlies both the
choice of types and the characterization of them. The first
piece is a lecture that I, at the invitation of the University,
gave in its ceremonial hall some days after my eightieth birthday,
and in which I pointed out that my philosophizing was essentially
determined by the wish to find out how
one should humanly take one's stand, on the one hand toward the struggle for
existence, on the other toward the spiritual division of labor, which
among other things seems to force science and religion onto entirely different
paths. The second piece gives a survey of my
work as an author, which can show in what relation my work in the philosophy of religion
stands to my other works.

Next come characterizations of thinkers who have asserted
the right of scientific inquiry in the understanding of the
% --- p. 8 ---
\setcounter{footnote}{0}%
\opage{8}religious life, as well as of thinkers who have held that there was
a sharp, an insuperable opposition between intellectual
and religious interest. The characterizations of Spinoza and Pascal
constitute the most important sections of the book. It is owing to a request
from a French philosophical journal that I came to draw
a comparison between Pascal and Kierkegaard, the latter
of whom may with a certain justice be described as the Danish Pascal.
The comparative method has here proved beneficial for illuminating
our great religious thinker. I had already earlier
employed this method in a short essay on Kierkegaard
and Nietzsche (in the third volume of my Mindre Arbejder).
The essay on Sibbern and Kierkegaard, too, is an example
of the application of this method. And the article on Arthur
Pfungst's attempt to practice Buddhism in the midst of
European cultural life can be seen from the same point of view, in that
its subject forms an analogy to what Kierkegaard called
``modern Christianity''.

It has interested me greatly that some of our poets
have sought to give an account to themselves of the inward union of
poetry and religion, which will ever assert itself whenever
a faith develops along life's own way.

When at the end I let Shakespeare's Prospero have
the word, it is because I see in him the poetic representative
of the magnanimity which in the time of the Renaissance naturally
came to stand as the foremost trait of character, both
in relation to the great horizons of nature and changing fates,
and in relation to the people we meet in the struggle for
existence. Here is a form of faith that will always stand
its test.

Carlsberg, 18 April 1927.

\hfill\textit{Harald Høffding.}

\mainmatter
% --- p. 9 ---
\setcounter{footnote}{0}%

\opage{9}\essayhead{PHILOSOPHICAL POINTS OF VIEW}

\essaysource{Lecture given at the University on the occasion of my \\ eightieth birthday.}

It is a great honor that the University shows me by inviting
me to give a lecture on the occasion of
my eightieth birthday. In no better way could one
honor and gladden me.

It is now 64 years since I first walked up the University steps.
Already in my school days I heard Carsten
Hauch's lectures on the folk ballads, Ewald, Oehlenschläger
and Shakespeare. Then came my student years, in
which the conflicting conceptions asserted
by philosophers and theologians gave me plenty to think about. And
52 years ago I gave my first lecture as a Privatdocent.
Since that time there has been a constant interaction
between my lecturing, my personal
development and my work as an author. There is a natural
need to communicate what one has been occupied with in the inner chamber,
and through attempts at communication thought can
become clearer and more definite. I have often found true what
Zeuthen once told me as his experience, that both
the significance and the difficulty of a thought often first really
come out for us when we try to express it. But
when I think back on my work at the University,
I remember first and foremost with gratitude
the receptiveness that was shown toward what I could offer,
within the various circles of listeners --- the narrow
circle of those who wished to orient themselves more thoroughly in
philosophy, the larger circle of young students in their first
years, and the circle, sometimes quite large, that heard my
% --- p. 10 ---
\setcounter{footnote}{0}%
\opage{10}evening lectures. And then my work at the University brought
me good collegial relations with other university teachers,
out of which not infrequently a warm personal relationship developed.
By far the greatest part of my life has been tied
to the University.

\asterism

As the subject of this lecture I have chosen a look
at the problems that especially set my reflection in motion,
at the motives of thought that have especially been decisive
for my standpoint.

\asterism

When I began my more independent philosophizing,
Darwinism exerted an overwhelming influence on minds.
For my part my reflection was especially set
in motion by Darwin's strong emphasis that even the
most splendid forms of life must ``struggle for life'', that is to say,
do work, satisfy certain conditions in order to
persist and develop. Darwin has wrongly seen in the very struggle
for life a cause of the origin of new species. But
that does not diminish the significance of the picture he has given
us of nature. Nature no longer appears as an idyll,
but as a world full of changes, clashes and
struggles.

The question then had naturally to arise: how
does it go with the spiritual life, which after all is to persist and
develop within nature, and which is subject to its
laws? How can then, under this alternation, which
is determined merely by the relentless demands of life and circumstances,
the spiritual goods, work in the service of truth and goodness
retain their significance?

The answer I thought I found, and which I sought in a series of
writings to substantiate, amounted in brief to this, that the
spiritual life in the world is itself an uninterrupted struggle to
% --- p. 11 ---
\setcounter{footnote}{0}%
\opage{11}persist --- that it is an activity, a work, whether the work
is carried on with consciousness or not.

I found as the characteristic marks of spiritual life two
traits that were precisely of significance when it came to
asserting oneself against an environment. The one trait is openness,
receptiveness toward the new, toward all changes.
The other is that the new is brought into connection
with our earlier experiences, so that there arises
a constant interaction, through which now earlier
experiences are corrected through the new, now the new are
interpreted by means of the earlier. This calls for an inner
work, a struggle to overcome the chaotic and manifold,
which can scatter, as well as the habitual, which can
dull and blind. As the fundamental character of spiritual life
there then stands a unifying energy. Energy is the capacity to
perform work; work presupposes resistance; and in the
spiritual domain the resistance consists in dispersion and
dullness and in the often so sharp opposition between the
new and the habitual. The fuller and more comprehensive a
content of observations and memories, of impulses and
of inclinations that appears, the greater the energy that is required
to hold it together and form it into a whole.

At this point a study of the opposition between
diseased and healthy mental life became of great importance to me. Insanity
expresses itself precisely as incoherence or as
the one-sided preponderance of individual elements. When geniuses are not
infrequently insane, it is connected with the fact that in
them such a great abundance of drives, thoughts and
images wells up that a very great energy is required to order
them and direct them. It is no art to be master of
oneself when there is not much to be master of.

\asterism

Whereas psychology has the task of showing how
the spiritual life actually unfolds as a work for
% --- p. 12 ---
\setcounter{footnote}{0}%
\opage{12}persisting, it is the task of epistemology to investigate
which forms of thought underlie it, when the question is
asked not merely about our actual thoughts but about our valid
thoughts, that is to say, the thoughts that can lead to a secure,
again and again confirmed orientation in the existence outside
us and in ourselves. Science too is a struggle against chaos
and against the habitual, a work to get beyond the
scattered and the blind. It is the task of philosophy to point out
the fundamental thoughts (``categories'') which assert themselves
in this work, both when questions are asked and when answers are given.
Often for philosophy the questions are more interesting than
the answers. Such fundamental thoughts are born and die like other
thoughts. The task then is, supported by the history of the sciences,
to show how both their birth and their death are conditioned
by the demands that the constantly continued striving after truth
makes. In the life of thought, as everywhere, it is not the determinate
forms that persist, but the forces that lead to the arising of the forms
and can themselves continue to work under new forms.

The analogy that thus appears between science
and personality, in that on both domains the matter is a struggle
against chaos, has been of great importance to me. It is not
alien powers to which we surrender ourselves when we dedicate our life
to scientific work. Science is the common work of the race
to attain orientation in existence, just as it
is the task of the single personality to orient itself in its
small circle. And the very work in the small circle cannot be
carried on without participation in the great common work, in appropriation
or in cooperation. A conclusion to this work cannot
be thought of as long as existence presents new sides
to us. And it might indeed also be that existence itself is not
ever finished once and for all!

\asterism

Most sharply the doctrine of life as a struggle for existence had to
intervene in the ethical domain. It has indeed also
% --- p. 13 ---
\setcounter{footnote}{0}%
\opage{13}been said that it caused ``the greatest crisis that ethics has at
any time had to pass through''. Is life at all about anything
other than procuring means of self-preservation?
--- Against such a question I have especially
maintained that self-preservation is always conditioned by the stage of development
at which one happens to stand. There is always a
definite level that one will maintain and below which one will not sink
down. That this holds in economic respects, shows
daily experience. But it holds also in ethical respects.
The spiritual goods to which one has access one will not be
excluded from; precisely in our days many a hard struggle is fought
for this reason in needy regions. The stage on which one
stands will condition the way in which the struggle is
conducted. Conversely, the way in which the struggle is conducted
can show that one in reality did not stand on so high a stage
as one thought. --- But when there are very many different
stages, and not only actions but also judgments
about actions are conditioned by them, how then is a universally valid
ethics possible? Here I saw no other possibility
than to take one's stand on one of these stages and to investigate
which ethical judgments must be passed from it. On bare
ground one cannot philosophize, no more in ethics than
in other domains. For my part I found in
the feeling of humanity that developed at the close of classical
antiquity, was deepened by Christianity and in the eighteenth
century was freed from dogmatic barriers, a point of departure
that it would be fruitful to use, in order then
to compare the results one arrived at with
the results one would have to arrive at from other points of view.
A presuppositionless ethics there is not.

\asterism

Besides the natural-historical motive that asserted itself
for me in the questions hitherto indicated, there was also
a cultural-historical motive that set my reflection in
motion. This was especially the case with the religious
% --- p. 14 ---
\setcounter{footnote}{0}%
\opage{14}problem. The sting of the religious problem had for me arisen
from the division of labor in the domain of spiritual life.
At the most primitive stages we know, religion stands
as the only form of spiritual culture. It is in a similar
way when a great religion stands at its first source, thus in
Buddhism and Christianity in their original forms. Then religion is One
and All, determines all thoughts and all tasks. Later on it is attempted
perhaps --- on the grandest scale by Catholicism in the thirteenth
century --- to unite all the culture of the age into one great whole,
so that religious ideas give the final points of view
and solve the last riddles. But on none of these stages
do we stand now. Religion has become a single element in the life of the spirit,
neither the only one nor the all-governing. Science,
art and morality stand on their own legs, and social life
develops according to its own laws. Religion is
no longer the pillar of fire that leads humanity's wandering
through the desert; it is rather the ambulance that
follows behind and picks up the weary and the suffering. It
is the consoler, not the guide. This too is great. And there
is no other spiritual power at hand that can take over
this task. But that does not take the sting from the problem.

If anywhere, it holds here that what persists
in the world, the spiritual as well as the material, is not the forms,
under which life appears at a given time, but the
working forces that have created the old forms and can
create new ones. For my part I have thought to find as the kernel
of religion at all stages a need and will
to hold fast the persistence of the values of life beyond the limits
within which human will can work for them.
One can read off through the history of religions what men
at various times have set highest. And the need
that has been at work here can continue to work, even after the
spiritual conditions have changed to a considerable degree.

There is, as I once stated, a standpoint at
which, through life's mixture of tragedy and comedy,
of victory and defeat, a fundamental mood has been won which
% --- p. 15 ---
\setcounter{footnote}{0}%
\opage{15}at once has the character of seriousness --- by reason of the
great things that have been experienced, and the pain it has often cost
--- and of irony toward every attempt to express in images or
concepts what in its fullness and by its oppositions bursts
every form one tries to fix. The view of life or life-feeling
that thus arises holds fast the conviction of a great coherence
in existence, as regards values as well as other
subjects and elements, and can therefore turn jestingly
against the dispensations of life, because they stand powerless
before the highest and the best, as well as against all attempts
to utter the unutterable. --- Whether one will give
such a faith the name of religion or not is immaterial.
We perhaps lack an apt word for what may be present
here. If we have in many cases too many words, then
here we have too few. As regards the old religions, their
lasting value, from a standpoint like the one indicated,
will essentially consist in the great poetry through which they have
been able to express the oppositions of life, its high points
and its depths. It is also through their
cult and the art that has entered into the service of this cult
that they still today have so great an influence on
the minds of people.

\piecerule

I have here touched on some of the motives of thought that have
asserted themselves in my philosophizing. Of the significance of
philosophy I have never doubted. I have been in doubt about
means and ways, still more often about my abilities to work
in the service of philosophy and satisfy the intellectual and
personal demands that philosophy makes on its true devotees.
But that there are here tasks and works that must continually be
taken up, of that I have never doubted, however high it
was from my standpoint up to the great figures of philosophy.

On my work desk stands a photograph of a bust
of Plato, which one of our art historians has given
me. When I sit down in the morning to my work
% --- p. 16 ---
\setcounter{footnote}{0}%
\opage{16}in front of this picture, I often feel unworthy to count
myself in the same guild as Plato. The age of Platonism is past,
and all ``isms'' will vanish. But the drive of thought that led
Plato, all his long life through until he died 80
years old (a fitting time to die at), in ever renewed
striving to seek back to the fundamental thoughts of knowledge and of life,
it will not die. And it will be able to stir itself in
the most modest activity, when this is joined with
honesty and with intellectual probity.

\asterism

When 53 years ago I had defended my dissertation
for the doctorate, I concluded, according to custom,
but not merely because it was custom, by expressing
the best wishes for the University. With that I will also
conclude this evening. May this University, hitherto Denmark's only
scientific institution of higher learning, continue to work through the ages
and bear witness to our people's ability and will to
contribute to the intellectual culture of mankind!

% --- p. 17 ---
\setcounter{footnote}{0}%

\opage{17}\essayhead{ON MY WORK AS AN AUTHOR}

In my earliest youth I lived myself into two lines of thought
that were later to set me the gravest problems.
Already in my school years I had pictured theology as
my future study. But reading Plato's Symposium
in the last year of school filled me with deep admiration
and brought me a testimony that on purely human
ground too high and significant
thoughts could arise. In the course of my years as a student I continued
the study of Plato, and to this was then added the study of German
philosophy from Kant to Hegel. The opposition between theology
and philosophy was stressed very sharply by the philosophical professors who
at that time
taught at the University, Rasmus Nielsen and Brøchner, each
in his own way. I then felt myself more and
more divided. I wandered from the theological lecture hall
to the philosophical, and back again; they lay each at its
end of the long university corridor. It was a time full
of torment, and I was often on the point of despair. But the author
in whom I immersed myself most was Søren Kierkegaard, and with
his help I at last caught sight of a new possibility. His mighty
struggle against the National Church, and against modern
Christianity generally, still resounded. But it was nevertheless
above all his philosophical writings that captivated me. I learned
from him that religious conviction was in the end a matter of personal experience
and rested on personal need, and thereby
the task was set me of testing whether the content of Christian faith
really corresponded to my own innermost need. The answer
came only after a long struggle. For a time I tried
to keep my personal faith and my philosophical interest
entirely apart from each other. Kierkegaard regarded all
% --- p. 18 ---
\setcounter{footnote}{0}%
\opage{18}scientific knowledge as ``inessential'' in relation to the
religious need. But this ``inessential'' knowledge came
now ever more to have ``essential'' significance for me. Yet
Kierkegaard also prepared my
later development in a more positive way. He maintained that the formal
peculiarity of the spiritual life consisted in the unity and concentration
with which the content of life is held together, and that the standard
of a spiritual life's value consists in the relation
between the fullness of the content and the energy with which the concentration
is exercised. This thought, which in the history of philosophy is represented
especially by Leibniz and Kant, forms the foundation of
the characterization Kierkegaard has given of the different
ways of taking life, his so-called ``stages''. Its significance
was naturally not clear to me from the first.
Only through a critical study of Kant and of positivism,
and especially during the preparatory work for my presentation of psychology,
did its great theoretical and practical significance
dawn on me. I now felt myself ever more a philosopher, and
after I had taken the doctorate in 1870 with a treatise
on ``Den antike Opfattelse af Menneskets Villie'', I began
as a Privatdocent to lecture on philosophy.

After I had lectured on various parts of the history of
philosophy (especially on ``post-Hegelian philosophy in Germany''
and on ``English philosophy in our time'') I gave in the autumn of
1874 a course of lectures on ``Introduction to Philosophy'',
in which I set out in all brevity the conception of the
philosophical problems and their mutual relations that had
taken firm hold in me. And then my work as an author proper began,
as I set myself the task of treating
each problem for itself. The short program I had
given in the lectures mentioned had naturally to be changed in many ways in the
working out of detail and through the new studies I
undertook. There was a great deal of work
to be done in the various fields. When I now look
back, I can divide my work in the following 40 years
into three periods. From 1875 to 1887 psychology and ethics stood in
the first rank; from 1887 to 1895 I was chiefly occupied
% --- p. 19 ---
\setcounter{footnote}{0}%
\opage{19}with the history of modern philosophy, and from 1895 on I have
been engaged with philosophy of religion and epistemology. Epistemology
became more and more for me the central discipline
in philosophy. Philosophy is for me an endeavor
to comprehend the work of human thought under certain
general points of view, an endeavor that is prepared within
the individual sciences when these become conscious of their guiding
thoughts.

My ``Psykologi paa Grundlag af Erfaring'' appeared in the autumn of
1882. The main thought of the book, or rather the hypothesis
it seeks to ground and confirm, is that the
fundamental law of mental life is the law of synthesis, of comprehensive
activity. I find the grounding of this hypothesis partly in the way
in which consciousness passes over into unconsciousness,
and conversely awakens from it, partly in an analysis of some
important mental phenomena and processes (memory, comparison,
grief, doubt, decision). And the confirmation I find
in an investigation of the conditions under which
sensation arises, as these are illuminated by
Fechner and the experimental psychology founded by him,
further in an investigation of the association of ideas,
imagination, and the main forms of the life of feeling and the life of will. Synthetic
energy, will in the widest sense of the word, is for me
the central expression of mental life. As long as
mental life exists, work is done to hold together a
more or less rich content. Feelings are symptoms of
this willing, and cognition is from the first its servant,
later perhaps its master. In the later editions an attempt has been made
to carry through this point of view with greater consistency.

The basic thought of the book developed in me successively,
from the Kierkegaard studies of my earliest youth through my
Kant studies and through studies of the diseased mental life in
comparison with the healthy. --- For the French translation
of the book Pierre Janet wrote a preface, and
since, as said, a comparison between diseased and healthy mental
% --- p. 20 ---
\setcounter{footnote}{0}%
\opage{20}life had been of great importance for the working out of
my hypothesis, it was a particular joy to me that the book could
appear in a major language under the aegis of an eminent
psychopathologist.

I see the shortcomings of the book very well (perhaps I do not know
them all). But I believe that its main thoughts are justified
and will constantly emerge anew, perhaps in forms
that we cannot now imagine. The sources of psychological insight
are indeed so manifold (self-observation, experiments,
physiology, psychopathology, the history of literature and religion,
etc.) that it is already now more difficult to use them
all than it was when in the years 1876---1882 I worked out
my psychology.

For my activity as a university teacher, especially after
I became a professor in 1883, the psychological way of regarding things
that I had practiced was of great importance.
Toward the spiritual movements of the time there was here a standpoint
that could meet the most diverse tendencies with
intellectual sympathy, without anticipating the final
epistemological and religio-philosophical thoughts. My
old teacher Sibbern (whom I had heard in my first year as a student,
but whom I came to understand properly only later)
said that just as for the botanist there is no such thing as
a weed, so for the psychologist there is no spiritual
phenomenon that cannot be an object of observation and
interest. Before we pass judgment on a spiritual phenomenon's
intellectual or ethico-religious value, we must recognize and
understand it. All too often one goes the opposite way and
stamps as stupidity or sin what yet has its
natural root in the connection of life and, rightly treated,
perhaps can reach up into the world of truth and goodness.

The next problem that occupied me was the ethical one. Already in
1876 I had published a little work on it (Om
Grundlaget for den humane Etik). But it now seemed to me
that the presuppositions of an ethics had not been sufficiently
% --- p. 21 ---
\setcounter{footnote}{0}%
\opage{21}critically investigated here. In my psychology I had made a
sharp distinction between fact and value (validity).
The ethical problem now presented itself to me in the following
way. Just as logic and mathematics rest on presuppositions common to all
human beings, and just as the empirical science,
e.g. psychology, rests on the law of causality, which is
universally recognized, so too there are universally recognized
presuppositions on which ethics can build, and which are
they? It cost me no little work before I came to the
conviction that, when one makes sufficiently strict demands
on the grounding, there is no ethics that in a
rational way can be forced upon all, as is the case
with logical, mathematical and natural-scientific propositions. Psychological
observation and historical experience show
that human actions are judged on the basis of highly different
presuppositions. Even if ethical standpoints may be
in agreement in their results, they will yet often go different ways
in their grounding. By logical analysis we can get back
to the presuppositions peculiar to a single definite ethical standpoint,
but these presuppositions themselves
are connected with psychological and historical conditions
that are not present in quite the same way at
other standpoints. I found certain main types by which
the presuppositions of ethical judgments can be determined --- from the standpoint of the moment,
the standpoint of individuality, the standpoint of family,
class and nationality to the universal standpoint of humanity.
It was impossible for me to derive these
different standpoints from ``the Idea of Man'' and to ascribe
to any of them a unique necessity.

The position of ethics can to a certain degree be compared with
the position of astronomy before Newton. Then too different systems were possible
in astronomy. I do not know whether ethics
will get its Newton; so far in any case he has not come.
Every attempt at an ethics must therefore be clear
about its psychological and historical presuppositions. As for
myself, I found a significant foundation for
an ethics in the universal sympathy, the general humanity
% --- p. 22 ---
\setcounter{footnote}{0}%
\opage{22}feeling that developed in the Greco-Roman world,
especially after the time of Alexander, and then, deepened and widened by
Christian love of one's neighbor, emerged again in the
ideas of the 18th century. My task was now to analyze and
deepen this foundation and from it to draw conclusions with
regard to the value of human modes of action and institutions.
What could be grounded from this standpoint had
then to be compared with the conclusions that could be
drawn from other standpoints (the self-preservation of the individual,
family, nation, etc.). ``The Idea of Man'' expresses
itself for us, as it happens, in different types, some simultaneous, some appearing
successively. Here we have what Fichte from
his idealistic point of view has called a projectio irrationalis. The special evaluations are determined in the end by a
fundamental motive, a fundamental value. One can set up a civil
or social morality without asking about the motives, but the
ethical concerns, as it happens, not only action but also
disposition. And as regards the motives, a distinction must be made between
the motive of evaluation (fundamental value) and the motive of action. They
do not always coincide. Thus humanity
as a motive of evaluation can give rise to special motives of action
such as sense of justice, courage, helpfulness, etc. The motive of evaluation
grounds in the end the value that is attributed to the
different motives of action. What in one ethical system
is a motive of action is perhaps in another system a motive of evaluation,
and conversely.

The special working out of my Ethics was undertaken on the
basis of a psychological, historical and economic
material that I had collected in the course of the years. As far
as possible I went into the special fields of human life.

In my conviction neither a supernatural
nor a metaphysical sanction of ethics is needed.
Human nature, as we know it from
experience, offers sufficient motives of evaluation and of action.
My Ethics, which appeared in 1887, was attacked partly
from the theological side, partly from certain antitheological sides.
There are antitheologians who meet in brotherly union with
% --- p. 23 ---
\setcounter{footnote}{0}%
\opage{23}the theologians in the assertion that an ethics without supernatural
sanction and without metaphysical foundation is impossible. Against
both parties I maintain that even when one has given up the
old faith, serious and positive work for ethical conviction
and conduct of life is not only possible but necessary.

--- It is now 35\footnote[1]{In 1927 it is 40 years since,} years since my Ethics appeared, and its most essential
thoughts I still hold fast. The book also continues
to find readers both abroad and at home. --- Only on
a few points have I made changes
and additions in later editions.

After the completion of my psychology and my ethics I set
myself the task of giving a presentation of the history of modern
philosophy. This task had a natural connection
with the plan of undertaking an investigation of
the main philosophical problems in their mutual relations,
which I had already conceived in 1874, and which occupied
me greatly.

As I stated above, Greek thinking and
view of life played a large role in my personal development. By it I
had won a deep conviction of the independent significance
of science and of the human life. And in my
Ethics I had maintained that our ethics is the Greek ethics,
deepened and widened by the spiritual life of the following ages. Greek thinking
lies, in general, at the root of all European
philosophy, and every study that goes deeper must again and again
go back to it. I have conducted laboratory exercises
on Plato's ``State'' and on Aristotle's ``Ethics''. But I
had to limit myself, and for the posing of the problems in our time
the philosophy of the last three centuries was naturally of the greatest
significance.

In my treatment of the history of modern philosophy I laid
greater weight on the problems than on the systems, and
I sought rather to characterize than to report. Through
characterization one passes in a natural way into
criticism. Characterization includes also an investigation
of the relation between the philosopher's personality and his doctrine,
% --- p. 24 ---
\setcounter{footnote}{0}%
\opage{24}and several of the small biographies in my book I have quite
written con amore. Moreover I have laid weight on philosophy's relation
to other sciences, and therefore in my presentation
men such as Copernicus, Galileo, Newton, Darwin and
Robert Mayer occupy a prominent place, in that their researches
occasion significant changes in the philosophical
formulations of problems. --- The book appeared in 1894---1895. On a few
points I had earlier done monographic preparatory work
(on Kant's philosophical course of development and on Kierkegaard);
other points I have later taken up for monographic
treatment (Rousseau, Spinoza). In comparison
with most German presentations I have laid greater weight
on French and English philosophy in the 19th century. I
had here a preparatory work in my book ``Den
engelske Filosofi i vor Tid'' published in 1874. --- Later I have written on
``Moderne Filosofer'' (1903) and on ``Henri Bergson's Filosofi''
(1914)\footnote[1]{Later I have (see below) written on Plato and Spinoza.} ---

After the publication of my historical work my studies concentrated
on epistemology and philosophy of religion.
One of the preparatory works for my epistemology lies in a
treatise on ``det psykologiske Grundlag for logiske
Domme'' (1898). But for the time being I was especially occupied with the
religious problem, after it had for some time not interested
me. I had not thought of undertaking a special
investigation of it, beyond the hints that had been given
in my psychology and in my ethics. But now the reading of biographies
had become my favorite reading, and I thereby came
also to occupy myself with Augustine's ``Confessiones''.
They interested me extraordinarily, both psychologically and
from the standpoint of the history of religion. And when I now immersed myself in the tenth book,
where the great teacher of the church raises the question: Qvid
est quod amo, quum deum amo? (What is it that I love,
when I love God?) --- it became clear to me that there lay an
important source for the understanding of religious phenomena in the self-testimony
of personalities who, while they live wholly
% --- p. 25 ---
\setcounter{footnote}{0}%
\opage{25}and fully in their religion, have need and capacity to clarify to themselves, in the form
of reflection, what they experience in the times of religious
rapture. I then collected a series of similar
testimonies --- from Indian literature, from the mysticism of the Middle Ages, from Luther and from Pascal --- and analyzed them psychologically.
I was naturally aware that by this
route one cannot solve the historical problem of the origin of religion.
Here I have often been misunderstood, especially by those
who think that when they have ``refuted'' the religion of the Hottentots,
they have ``refuted'' all religion. In general the
question of the first origin of religions can no more
be answered directly than the question of the first
origin of language. But by the route I took one could
perhaps discover what sustains religion, --- what
gives it, once it has come into being, its power and
significance. And this is for the philosopher, who wishes to investigate the present
state of the problem, more important than the question
of the first origin of forms of cult and of dogmas. I saw
naturally also that the study of the subjective and
individual side of religions had to be combined with a study of the history of religion.
And here I sought especially to trace what
standard historians of religion apply when they call
one religion higher or lower than another. By both routes
I sought to form a grounded conjecture, a hypothesis,
about the foundation of religion in its different forms.
And with that I thought I could also find the basis for
an understanding of the position of religion in our days and for
a conjecture about its future. --- In a speech at a University festival
in 1897 I gave a provisional characterization of my
standpoint, and in a course of lectures which I in two semesters
(1899---1900) gave for students of all faculties, I gave
a detailed presentation of my conception. In the following
year (1901) my Religionsfilosofi then appeared.

Every ecclesiastical theology has a sum total of religious
conceptions to which it ascribes literal validity. Within
one and the same religion there can yet be different
opinions as to where the boundary is to be drawn between the
% --- p. 26 ---
\setcounter{footnote}{0}%
\opage{26}figurative and the literal. But philosophy must subject
this question to a principled investigation, in that it compares
religious conceptions with the conceptions that
are formed on a scientific basis and by scientific method.
The first investigation within the philosophy of religion
will therefore be epistemological. An analysis of the main concepts
with which religion operates led me to
the result that these concepts are formed by analogy, and
an analogy that cannot --- like the analogies that in science
serve as working directions --- find confirmation
by close observation. The language of religion is a pictorial language
that stems from a time when no distinction was made
between image and thought. Such a distinction can be made only
when a strictly testing science has developed.
It is not so much the results of science as
the scientific method, the scientific way of
posing questions and forming concepts, that
is of significance here. The questions of religion can neither be posed
nor answered by science, and the answers of religion
are not answers to the questions of science.

These are thus different interests that appear in
religion and science. What, then, is the interest we
call the religious? This question the psychological part of the philosophy
of religion tries to answer. My endeavor
is to regard religion as a psychological
phenomenon. The task is to find out in which region of mental life
religion has its roots, --- what need it can
satisfy, --- and then to investigate the conditions for the continued
existence of this need. It is not enough here to establish
that the same dogmas are taught and the same cult practiced. One must
investigate how these dogmas and this cult work
on the different individualities, and especially on what it rests
that human beings again and again return to them. It
then turns out that both dogmas and forms of cult are connected
very closely with what for human beings at the
given stage is the supremely valuable, the object of their
wish and of their striving. Here the concept of fundamental
% --- p. 27 ---
% pp. 27-28: missing from the scan; translated from text set from the e-book (see the transcription).
\setcounter{footnote}{0}%
\opage{27}value presents itself again, and my hypothesis in the philosophy of religion is then that
the essence of religion has its ground in a need for or a wish to see
the persistence of fundamental value assured, also beyond the limits within
which human work can act for this persistence. When it
now comes to the point that the religious conceptions cannot be held fast in
their literal meaning, then the need to hold fast the persistence of
what is valuable must seek expression in symbols, in poetry of life, i.e. in a
poetry growing out of life itself. The ideal would be that each individual
produced his own symbols. But it will also, from the standpoint of the individual,
be justified, by involuntary interpretation, to find again in the traditional religious
conceptions the thoughts to which his experience of life has led him. Far more than
those who are clearly conscious of it have made the transition from
faith in dogma to poetry of life. One takes from the religious
tradition what one needs and lets the rest lie. This is done by
every church and every religious party. By that route modern
Christianity has become another religion than primitive Christianity, in that it
has in fact let fall essential presuppositions that decisively determine the
ethics of the earliest Christians, in order to accommodate itself to the new historical
conditions.

It is for me ever more a purely psychological interest that leads me
to occupy myself with the religious problem. I find there, where one
can gain an insight into the inner life of deeply religious personalities, a
richness in mood, a grandeur in imagination, a passion in will,
that is not only an attractive and captivating object for the psychologist
as such, but also leads him to ask whether forces like
these can disappear from the spiritual life of human beings without new mental
forces forming a real replacement for them. The idea of psychic
equivalents, which already appears in Plato, leads to the search for new
sources of enthusiasm and inner rapture. And this search is actually
itself religious, since it is motivated by a need to see the values
conserved.

Sometimes I have put the religious problem in the following way. Is
religion an organ or is it a tumor? If it is an organ, then it may
be perilous to life to operate it away, even if it should belong to
the 107 superfluous organs which according to Metchnikoff we still possess.
The question can only be whether it can continue to get nourishment, and whether
its functions can be preserved through exercise, or whether it dies away from
lack of both. If it is a tumor, an attempt to
operate it away may occasion such a loss of blood that life is in
danger; it would be safer to ligate it, so that it can die from
lack of nourishment. --- In both cases it is future experience that
must bring the decision.

This last consideration leads over to the ethical part of the
philosophy of religion, which partly investigates the possibility of a conflict
between ethics and religion (namely if faith in the persistence of values
keeps one from working to produce and conserve values), partly
discusses the question to what extent faith in the persistence of value can yield
strength and power for ethical work. ---

There are no other points of view in the philosophy of religion than the three
indicated, the epistemological, the psychological and the ethical. A
special method of the philosophy of religion there is not. ---

In a treatise that I have called ``Forord og Efterskrift til min
Religionsfilosofi'', and which is found in the second volume of my ``Mindre
Arbejder'', I have tried to illuminate my position toward the religious
problem as philosopher and as human being. Also in ``den store Humor'', which
will be discussed later, this position is illuminated. As a philosopher I have only to
do with the general principles of science and of life; as a human being I have
to do with the special nuancing and application of these principles
to which my life's work and my personal experiences have led me. I
believe that the philosopher and the human being in me have drawn nearer to each other in
the course of life, and that neither of them has been harmed thereby. --- I
refer especially to the section ``Tragik og Humor'' in the book on humor, as the
clearest and deepest expression I have
% --- p. 29 ---
\setcounter{footnote}{0}%
\opage{29}been able to give for the region where my philosophy and my
personal experience of life meet each other.

The problem of knowledge had made itself felt indirectly
in all my writings. In the Psychology I had marked very sharply
the difference between the psychological and the epistemological
point of view. In the Ethics the strict
demand for scientific grounding had led me to doubt
about a rational ethics. In the philosophy of religion epistemological
points of view lay at the base of the denial of the
intellectual significance of religion. In the introduction to ``Den
nyere Filosofis Historie'' I had set up the problem of knowledge
as the central problem over against the psychological,
the ethico-religious and the cosmological (metaphysical)
problem. In a university program (1902) I gave a preliminary
investigation of the main philosophical problems in their
mutual relations. After these and other preparatory works there appeared
in 1910 ``Den menneskelige Tanke, dens Former og
dens Opgaver''.

The book first seeks out the life of thought in its involuntary forms.
Already in the involuntary work of thought can be traced
the forms of thought that in scientific thinking are developed
further, specialized and carried through with greater exactness
and consistency. After the investigations of
the involuntary functions of thought have shown the possibility of
knowledge, the point is to find the definite forms
by which real science can be gained. The realm of possibility
is greater than that of actuality. Here Wallenstein's
words hold: ``Eng ist die Welt, und das Gehirn ist weit''. It is
the history of the sciences that shows us the forms of thought in
which true knowledge can be gained. These forms, the so-called
categories, are not the same at all times. It was
Kant's fateful error that he believed he could set up once
for all a system of categories. The categories
arise sporadically, but can afterwards be collected and ordered,
e.g. in such a way that one proceeds from the most comprehensive
to the more special. I begin my theory of categories with
% --- p. 30 ---
\setcounter{footnote}{0}%
\opage{30}the concept of synthesis and the correlate concept of relation; these
two concepts open the series of the fundamental categories,
which show themselves in all thinking, the pre-scientific
as well as the scientific. From there one passes to the
formal categories (identity; ground and consequence, and others), which
underlie logic and mathematics, and from there to the
real categories (cause, totality, development), which underlie
all empirical science, and from there again to the
ideal categories (the concepts of value), which show themselves
in aesthetics, ethics and philosophy of religion. The transition from
the real categories to the ideal I find in this, that every
evaluation presupposes a relation to a totality (whether
this consists of a single moment or a whole life of personality,
of a single individual or a larger or smaller
society). --- Just as the categories arise, they can
also die; thus I regard e.g. ``substance'' as a
dying category.

The theory of categories can be called a descriptive epistemology
and stands as such between psychology and the history of science
on the one side, the discussion of the philosophical problems
on the other side. The problems arise through the relation
of the categories to the factually given subject matters, which are not
mere products of certain intellectual functions. This shows itself most clearly
in the problem of knowledge. Can our
forms of knowledge be derived from the given subject matters, or have
they, however much they are conditioned and specialized by the subject matters,
their ultimate ground in the structure of the human spirit
(which yet need not be unchangeable)? Here
criticism as founded by Kant will, as I try to show,
always maintain a first place in the epistemological
discussion, because it can apply the analytic
or regressive method for finding out the forms and principles of knowledge. Such an investigation underlies
dogmatism, empiricism, pragmatism and the theory of postulates alike.
Only analysis of our knowledge and its conditions
can teach us what is to be set up as dogma,
as experience, as economy or as postulate.

% --- p. 31 ---
\setcounter{footnote}{0}%
\opage{31}In one of his letters William James refers to me as
``good pragmatist and pluralist''. That rests on a misunderstanding.
A pragmatic (``economic'') view of science
has for me only historical interest, namely as
an illumination of the motives that have led to science, not
as a grounding of the validity of science. And among
these motives I give the purely intellectual interest the
first place, whereas most pragmatists do not express themselves
distinctly enough on this point. And pluralism is for
me a resistance that is to be overcome; it sets a task;
but is no solution and makes (as in our literature
Herbert Iversen clearly saw) all solution, even all striving
for solution, impossible. --- James, whom from my stay
in America I remember with cordial friendship, has inferred too
much from the agreement with which we both stood
against the metaphysics of the Absolute. Yet in America I had on two occasions
the opportunity to develop my philosophical
points of view in James' presence. The two lectures were worked together, at
James' suggestion, into a treatise in an American philosophical
journal. They are printed in Danish in the third volume of
my ``Mindre Arbejder'' under the title ``En filosofisk Bekendelse''.
---

The cosmological (metaphysical) problem arises when
the question is raised whether a scientific
world-view can be won. The difficulty here is that
the whole of existence is never given as a subject matter for our thinking.
It is always a single part or side of existence
that we take as basis; ``the whole of existence'' is then conceived by analogy
with the ``primal phenomenon'' thus taken as basis. The history
of world-views is thus the history of
primal phenomena. Within our experience analogy can be a good
means of knowledge, in that by its help we illuminate one
field of experience by means of another. But existence
as a whole does not lie completed before us, and even if this
could be imagined, we still could not characterize it
through comparison with other subject matters. In the choice of the primal phenomenon,
of the part or side of existence from which
% --- p. 32 ---
\setcounter{footnote}{0}%
\opage{32}one infers to the whole, and in the way in which
the analogy is laid out and carried through, the individual personality
of the thinker involuntarily makes itself known, and every attempt
at a world-view will therefore lie on the border
between science and art. The cosmological problem
is accordingly conceived more and more as a border problem, and
its treatment will have philosophical value only when
it is prepared by thorough epistemological work.
But there will probably always be spirits whose intellectual
urge to expand can be satisfied neither by
religious images nor by scientific concepts.

In the section on the problem of evaluation ``Den menneskelige
Tanke'' gives partly a purely principled, partly a more detailed
presentation of the points of view that underlie
my Ethics and my philosophy of religion. ---

When I am to state in all brevity my place in relation to the
philosophical tendencies, I must first lay weight on the fact that the problem of knowledge
has more and more become the central
problem for me. Actually, in my treatment
of other problems I pursue in the end this problem in its different
ramifications. Kant's great question about the possibility
of knowledge, in particular of experience, has also been
decisive for me. I have learned much from positivism,
both in its French and in its English form, and my conception
could be called a critical positivism, because I lay
great weight on the content of experience, while I
see a problem in the very concept of experience. And I have learned
much from Spinoza, but my doctrine of unity, my assertion of
the unity and lawfulness of the world is not dogmatic like
Spinoza's; it is a critical monism. --- But all such
-isms have very little to say.

I have tried to philosophize in such a way that I held fast
philosophy as an independent scientific field. It
works for clarity and self-reflection in the spiritual world,
with an open sense for all sides of life, in the conviction
that only clear thinking can let the fullness of life come
to its due, although this thinking itself cannot take
% --- p. 33 ---
\setcounter{footnote}{0}%
\opage{33}the place of the fullness of life within which it seeks to
orient itself. I have an inward conviction of the significance
of philosophy as an indispensable and continually active
element in the spiritual life.

In my 72nd year I took my leave from the University and
could now devote myself entirely to the working out of various
ideas and sketches that had not yet assumed definite
forms. Shortly before, the Academy of Sciences had chosen
me as occupant of an ``honorary residence'' that the founder of
the Carlsberg Foundation had placed at the Academy's disposal. In
a beautiful residence with a large garden I found a good place to work.
In the years I have lived here, three kinds of works
have claimed me: psychological, epistemological
and philosophico-historical.

It had long been my wish to carry out a work
in the field of descriptive and analyzing psychology,
where according to my conception there was still much to do.
Experimental psychology had by its nature chiefly
occupied itself with the elementary phenomena of mental life.
But it is of great importance for the spiritual
life that also within the higher, more complicated
and totalized mental states, in which precisely the innermost
tendencies of our mental life appear, a clear and
definite connection should be demonstrated. And it is not only a matter of finding
connection. It is also important to characterize the
differences and oppositions that can assert themselves here. I had for a particularly long time occupied myself with the concept of
humor. In my Psychology I had quite briefly described
humor as a special form of the feeling of the
ludicrous, and I had presented scorn and humor as
the great opposites in this field. Most phenomena of laughter
lie between these two opposites. Now I wished
to investigate humor more closely, and I came
to regard it as one of the feelings in which life can
as it were sum itself up for the individual, --- as a kind of
higher life-feeling, which is determined by the relation between
% --- p. 34 ---
\setcounter{footnote}{0}%
\opage{34}a serious foundation (which can be of differing nature)
and the manifold ills, trifles and absurdities of life.
In everyday speech one understands by humor
all kinds of more or less witty jesting; that is what I
call ``den lille Humor''. But I wished to describe ``den
store Humor'', which presupposes a deep foundation of experience in
the life of feeling, while the foreground finds expression in laughter,
or at least in a smile. It was likewise my task to determine
the relation between humor as life-feeling and the religious
feeling. Here Kierkegaard, in his eager striving
to maintain sharp differences, had placed humor in an all too
low position. The section in which I treat this point,
namely the chapter on ``Tragik og Humor'', I regard myself
as the most important in the book and at the same time as the classic expression
of my own total feeling toward life. --- The method
in this book is, as in my philosophy of religion, psychological-historical.
I found among other things here a welcome opportunity
to express my admiration for two of the greatest figures
I know, the one belonging to the world of thought, the other
to the world of poetry: Socrates and Shakespeare. --- In the book's
last section (``Humor og Filosofi'') I think I have drawn
attention to an important matter, namely the relation
between a researcher's state of mind while he is engaged
in intellectual work and while he moves
in the circumstances of daily life. One cannot without more
infer from a researcher's scientific character to his
state of mind under the small and great events of daily life.

Another psychological work I published in 1918 under the title
of ``Oplevelse og Tydning''. My friend Théodore Flournoy in
Geneva had some years before published some very interesting
notes by a lady who had consulted him
about her mental condition. There made itself known
in her inner life a sharp opposition between experiences
of the most exalted and ideal kind and other experiences
of the most sensual-passionate kind. Often
% --- p. 35 ---
\setcounter{footnote}{0}%
\opage{35}those first experiences took on the character of ecstasies. In her sound
moments the psychologically understanding lady stood freely and
clearly over against these highly different experiences. In particular
she did not have the inclination, frequent among mystics, especially in earlier times,
to regard ecstatic experiences
as revelations of the divine figures of the Bible and the ecclesiastical tradition,
nor did such conceptions weave themselves into the visions themselves. In this respect
she formed an interesting contrast to the Spanish
nun Sta Teresa (16th cent.), with whom I had occupied myself in the preparatory work
for my philosophy of religion,
and in whose writings I now immersed myself anew. My treatise
seeks to show how differently experiences that are purely
psychologically of the same kind are interpreted under different historical
presuppositions and by different personalities. I
found here an example illustrating the basic thought of my
philosophy of religion. And at the same time I found here an analogy of epistemological
interest. For science too is an
interpretation of given qualitative experiences, and here too
experience and observation must not be confounded with interpretation, or
conversely. --- The book contains besides various treatises, especially
in the philosophy of religion, to which the title fits
more or less.

My epistemological works of the later years
consisted in more special investigations of various concepts
set up in ``Den menneskelige Tanke''. First I treated
the concept of totality (1917). This concept plays a role
in all reflection. Either reflection starts from a given
totality, whose elements it seeks to disentangle, or else
it seeks to form from a given manifold a new totality.
Where totalities present themselves of themselves, we dissolve
them; but when a chaos is given, we strive to transform
it into a totality. My motto was an epigram of Goethe,
in which he calls life ``einen wahren Schein'' and ``ein
ernstes Spiel''. Totality (atom, organism, personality,
society) is a semblance and leads to a mere play of thought when one
% --- p. 36 ---
\setcounter{footnote}{0}%
\opage{36}sees in it something completed once for all, or when one
sees in it not a task but an explanation. But there is
truth in the semblance, and earnestness in the play, when we seek the laws
by which manifold elements work together into a whole.

This line of thought was continued in a treatise on the concept
of relation (1921). It is in the two treatises the two
first and most fundamental categories, synthesis and relation,
that I treat. Both when we analyze totalities,
and when we seek to form a totality out of a manifold,
the point is to demonstrate definite relations --- within the
given totality or between the manifold elements. A
totality is understood through the relation between its parts, and a
chaos is overcome by demonstrating definite relations
between the apparently isolated elements. We think
in relations. Our thought is like a pair of compasses, one leg of which
stands fast, while a place is sought for the other leg,
and our knowledge corresponds to the relation between the two positions.
If we cannot find any position for the other leg,
we end in a problem, perhaps in an insoluble problem. In the treatise
this is demonstrated more closely for the different groups
of categories and for the different main problems. Occasionally
it is said that if Einstein's hypothesis is correct, there lies
in it a new example of the significance of the concept of relation,
a continuation of certain thoughts called forth by Copernicanism.
Quite particularly, finally, the significance of the concept of relation
for ethics\footnote[1]{Later I have (in Videnskabernes Selskabs filosofiske Meddelelser) treated yet another concept that played a role in ``Den menneskelige Tanke'', but required further elaboration, ``Begrebet Analogi''. (1922).} is developed.

In my works in the history of philosophy of the later years
I have been able to make use of notes and collections that
stemmed from the preparation for my laboratory exercises.
I had often gone through Spinoza's chief work with
older students. I could say with Lessing that
if I were to call myself after anyone, I would call myself
after Spinoza. Naturally not in the sense that I should
be an adherent of his system. But to several of his
% --- p. 37 ---
\setcounter{footnote}{0}%
\opage{37}basic thoughts I have come by my own way, and the deep
philosophical disposition that animates his thinking makes him
a great model for all thinkers. I now wished to facilitate
the understanding of his chief work by giving a coherent
analysis of its content and thereby a characterization
of his philosophizing. ``Spinoza's Ethica. Analyse
og Karakteristik'' appeared in 1918, like the two epistemological
treatises, in Videnskabernes Selskabs Publikationer.
---

Of other works in the history of philosophy from these years I mention
``Bemærkninger om den platoniske Dialog Parmenides''
(Videnskabernes Selskabs filosofiske Meddelelser.
1920). My result from an investigation of this remarkable
dialogue is that the criticism given in the first part of
the Platonic doctrine of Ideas strikes not only it, but all systematic
metaphysics, and that the second part contains so much
that is unfinished and unclear that perhaps here a posthumous draft has been
published by uncomprehending disciples. What interested
me here was the different position that the concept of ``the
sudden'' occupies in the doctrine of Ideas, where the leap from one
Idea to another stands as incomprehensible, and in the doctrine
of Eros in the ``Symposium'' (Plato's most important contribution to psychology).
Logically considered, there exists in Plato a sharp
difference between the different Ideas, but for Plato as
psychologist a lawfully determined course of development from
sensuous longing to the highest stage of spiritual life stood as
possible, and ``suddenness'' expresses only the opposition of consummation
to preparation. Plato as psychologist has
corrected himself as epistemologist.\footnote[1]{Later I have published in the same place a treatise on ``Platons Bøger om Staten''.}

A third historical work was called forth by the invitation
to write a book on ``Ledende Tanker i det 19.
Aarhundrede'' for the collection of works on the culture of the 19th century
that Gyldendal's publishing house and a Swedish publisher
are issuing in Danish and Swedish. I set up here the
hypothesis that the 19th century was the age of
% --- p. 38 ---
\setcounter{footnote}{0}%
\opage{38}confirmation by experience, of specialization and of technical application.
The great theoretical and practical ideas of the 17th and 18th centuries
were subjected to the test of experience, of experiment and
of life. The moral character of the 19th century
is its fidelity to the great thoughts of the past. True
fidelity consists not in reveling in the ideal world that
opens before us, but in translating it with perseverance
into reality and working it through in detail. As
a motto one could use the exhortation to the sons of the gods in
``Faust'':

\begin{quote}
Und was in schwankender Erscheinung schwebt,\\
befestiget mit dauernden Gedanken!
\end{quote}

I sought to show that Kant's philosophy lies at the base of
the century's two main tendencies, Romanticism and positivism,
and that the criticism of these two main tendencies in the century's
last part led to a resumption of the problem of knowledge
as the central philosophical problem, --- that the natural sciences'
greatest achievements (the demonstration of the persistence of energy
and of the natural development of organic life) were
the carrying through of thoughts that had already been expressed in the 17th and
18th century, --- and that the century's greatest social
phenomena, the great labor organizations, were a carrying through
of the 18th century's urge for freedom, since freedom
consistently must also be freedom of association. In
the religious field an analogous trait appeared in the demand
that religious conceptions, in whatever direction
they may go, shall be an expression of personal experience and
of the single individuals' honest refuge.

By an eminent Danish natural scientist it has been objected
against my conception that it underestimates the significance
of natural science. On this occasion I must point out
that my task is to speak of the leading
thoughts, and these point decidedly back to the preceding
centuries. The last three centuries form a magnificent
whole whose periods supplement one another. Whether the future
will be a continuation of the spiritual work
% --- p. 39 ---
\setcounter{footnote}{0}%
\opage{39}under the guidance of the same basic thoughts, or whether wholly new principles for theory and practice will
emerge, only the future itself can teach. The work of thought can be equally great
in the period of carrying through as in that of founding.

It has been the great good fortune of my life that I have been able to
devote myself to philosophical studies. What, then, has philosophy
been for me? --- Against the doubt and the fog that
plagued me in my earliest youth I have won clarity
as to the direction in which the great tasks for thinking and
acting lie, and how work must be done on their
solution (whether this is possible or not). In the work of thought
itself I have found the best help against despondency
and a heavy heart. Like Spinoza I have experienced that difficulties
and cares vanished when I could only
immerse myself in reflection (modo possem penitus deliberare!).
But philosophy has not been merely a household remedy for me.
It cannot be a household remedy, as it happens, when
it is merely a household remedy. I have won reverence for
the world of thought, wherever it reveals itself amid
the struggle and unrest of the times and points toward the great tasks
for our striving. I believe more in ideals than in
ideas. If we do not get beyond continual striving, this
might perhaps find its explanation in that existence itself
is not finished. Of that part of the world that stirs in human life
and in human thought it holds in any case.

I have the faith that philosophy will also in the future stand
as a necessary element in spiritual culture. It conditions
a spiritual horizon that no special training,
no special subject can give, but to which every subject can
give its contribution. And one thing there is that must particularly be stressed,
and which I can confirm from a long experience of life. Philosophical
reflection has a liberating influence over against habits,
prejudices and traditions. It makes room for a personal
relation to life, a first-hand relation. It is
a basic thought in my philosophy of religion that a genuine faith in life
grows only out of one's own experience. But a work of thought
% --- p. 40 ---
\setcounter{footnote}{0}%
\opage{40}can be the necessary presupposition for our really
getting to build on the very ground of life. A faith in life presupposes
not that all riddles are solved. Just as our courage and
our will must again and again be summoned to battle against fog
and chaos in us and around us, so again and again the courage
of thought will be necessary for work on the great riddles
of existence, which always appear in new forms\footnote[1]{In the treatise ``Erkendelsesteori og Livsopfattelse'' (Vid. Selskabs filosofiske Meddelelser 1926) I have, as in a philosophical testament, expressed my conception of the relation between research and faith in life.}.

What persists in the world of thought as in that of life is
not the forms in which thought and life appear at a given
time, but the forces by which these forms have come
into being, and which for a time have come to rest in them. To these
forces one must go back in order to fix one's hope for
the future. But even if one has the faith that spiritual values
will arise again and again, one has yet no right
to believe that we should be able to know beforehand what these values
will be.

\piecerule

% --- p. 41 ---
\setcounter{footnote}{0}%

\opage{41}\essayhead[Plato's Conception of Art and Religion]{PLATO'S CONCEPTION OF ART \\ AND RELIGION AS IT APPEARS \\ IN THE BOOKS ON THE STATE\footnote[1]{A lecture based on my treatise, published in Videnskabernes Selskabs filosofiske Meddelelser (1924): \textit{Platon's Bøger om Staten.} Analyse og Karakteristik.}}

Plato appears in his writings as thinker, as lawgiver,
and as prophet. The different interests that thus
stir in him not seldom come into conflict with one another.
Not only do the books on the State for this reason stand
in a certain opposition to other works of his, but within
the books on the State itself different points of view are adopted,
especially as regards the conception of art and religion.
As lawgiver he speaks chiefly in the third and fourth books,
and as poet he appears in the artistic shaping
of the work's individual parts. The philosophical foundation
that he believed he had found has made it possible for him to
work into a whole drafts that were evidently composed at different
times and at different stages of his development,
and thereby to create the work, so significant in the history of the spirit,
that is called Plato's ``State''. --- We shall,
as said, confine ourselves here to his conception of art and religion.

1. The books on the State deal with the concept of justice,
and Plato holds that it is easiest to illuminate this concept
by investigating how justice must
appear in a perfect state. In the description he
attempts to give of such a state, there is one question
that engages his whole interest, how the rulers in such
a state ought to be educated. For the time being he stands here at the
% --- p. 42 ---
\setcounter{footnote}{0}%
\opage{42}old Greek stage, where there is required --- besides gymnastic training
--- a spiritual development through absorption
in the national poetry, in which the people's belief in the gods
has found a full expression. An evaluation of the poetry
necessarily becomes an evaluation of the religious conceptions
that appear in it, and Plato then sets up, under
sharp criticism of Homer and Hesiod, certain theological norms
or models. One cannot, he holds, give the poets
free rein, because the young would so easily get false
conceptions of the gods, which the lawgiver cannot permit.
Plato does not say outright whence he knows which
are the right conceptions of the gods, but it is evident
enough from his discussion that the theological norms or
models are determined by his own philosophy. There are three
main norms.

a. The Deity is good and cannot be the cause of anything
evil. The poets must therefore not ascribe evil
deeds to the gods. The gods must not be represented as causes of
misfortunes and sufferings, except insofar as these may become
means of purification and improvement. But now there is ---
Plato teaches both here and in other writings --- more evil
than good in the world, and the Deity can therefore be the
cause only of very little. The real evil in the world must
have another cause than the Deity. --- By pursuing
this line of thought Plato came in his work of old age, the ``Laws'',
to the assumption of an evil world-soul, an assumption
that has been connected with a certain influence of
the doctrine of Zarathustra, which is said to be traceable in Plato in his
last years. --- A particular application of this first norm
is the prohibition against teaching that the gods are at strife
with one another. That would harm the concord that ought to prevail
among the rulers of the state! Not even allegorically may
such mythological conceptions be used, since young people
cannot distinguish between allegory and literal
truth.

b. The Deity ought to be represented as unchangeable. Could
a god change, change would be one with
% --- p. 43 ---
\setcounter{footnote}{0}%
\opage{43}degradation, abasement. A god must therefore not be represented
in human form. And a god will hate lying, never
conceal the truth; truthfulness is a kind of unchangeableness.
And by lying must then be understood not merely the outer
lie, that which expresses itself in words, but the inner lie, which
springs from ignorance, and which only self-deception prevents us from
discovering. It is only a crippled soul that is more ashamed
of voluntary than of involuntary untruth.

c. The third norm is that the gods must not be represented
as subject to strong emotions, such as fear of death, laughter,
grief, erotic desire. What is to be portrayed as the highest is
an equanimity, a self-command, that does not make itself dependent on
outer circumstances but is the right courage in the face of whatever life
may bring. --- And the rulers of the state are to imitate what
is thus asserted as the right conception of the nature
of the gods. The rulers have the one great task of providing for
the freedom and security of the state. Neither in their imagination nor
in their personal conduct may they imitate tragic or
comic poets, or in general appear in any character
other than their own. Plato denies here what he had stated in the ``Symposium'',
that it is the business of the same man to write tragedy
and comedy. It is the necessity of the division of labor that
has gained power over him and now receives an ethical application.
Among the poets the epic poets are preferred to the dramatists, because
the former do not, like the latter, gradually make themselves one with different
characters and thereby split themselves. --- It will not
be possible to bring agreement between Plato's
statements on the point here discussed. He himself, and
not only as a poet, had an eye for the composite character of life.
In a later dialogue (Philebus) he speaks of the tragedy
and comedy of life. But he kept having misgivings
at the thought that a single individual in his art, as poet or actor,
should express this composite nature of life in his imagination
or in his performance.

2. Already from this ethical-pedagogical discussion of art
in the third and fourth books one sees what strict demands
Plato makes of the ruling class in society, although
% --- p. 44 ---
\setcounter{footnote}{0}%
\opage{44}for the time being he merely wishes to reform the old Greek education.
But in his further development of the rulers' task
he comes to sharpen the demands, so that
he not only goes beyond the traditional pedagogy
but comes to make quite new demands, in that he requires
a strict scientific discipline of those who are to govern
the state. The rulers can, he holds, provide for the freedom and security
of the state only if, through all changes and
temptations, they preserve the thought of the state's true nature in
opposition to claims and tasks that might lead to falsifying
the right constitution. The guardians must therefore have the
capacity to defend and hold fast the thought of something that stands
firm and preserves its unity with itself through all changes.
But this capacity is trained only (Plato maintains at the
end of the fifth book) through scientific, especially
mathematical and philosophical, study, and in the sixth and seventh books
he then goes more closely into the arrangement of this study.
When people have often found a great opposition between
the third, fourth, and fifth books on the one hand, and the sixth and seventh
books on the other, it is because they have overlooked
that there is a kinship between practical judgment (e.g.
of a constitution's persisting in unity with itself) and the
strict carrying through of the concept of identity, which is the characteristic
of all science. When Plato sharpens the demands
on the rulers by requiring scientific insight
of them, it is the natural transition from mere likeness to
complete identity that he makes. Here lies the bond of unity
that ties together the various parts of the marvelous work,
even though they were composed at different times.

When Plato then, after having thus sharpened the demands
on the highest culture within society, again (in
the tenth book) returns to discussing the significance of art,
he stresses more strongly than before two points of view that already
made themselves felt in his first discussion of art,
namely the relation of art to truth and reality and
the nature of art as imitation. And he takes into account
not only poetry but also the pictorial arts.

% --- p. 45 ---
\setcounter{footnote}{0}%
\opage{45}Whereas in his first discussion of art he placed it
in relation to theology and morality, it is now its relation to
science that stands in the first line. And he then places
art far below science, because it does indeed give a
picture of something real, but an imperfect picture,
which falls as far short of the real thing or
action as these in turn stand below the general concept
or the general law under which they belong. A
poem about a heroic deed stands far below the heroic deed itself.
One should not write poems about exploits. In this connection
Plato again comes back to the danger of imitating
emotions and passions. He holds that
when the artist is fond of that kind of subject, it is because it is
easier to represent highly passionate natures than energetic
and harmonious natures. He whose thought can move
in the eternal world of the Ideas will in any case not have
an inclination to dwell on the feelings that human
relations can produce, since they lose all significance in the face of
pure truth.

In his first discussion of art (in the third and fourth books)
Plato speaks as lawgiver; in the new discussion that the tenth
book brings, he speaks essentially as philosopher, and the transition
has, as said, taken place through the discussion of science in the
sixth and seventh books. In contrast to the truth that scientific
inquiry can bring, art sinks down to
being an imitation of sensible things. It is Plato's idealism
that makes him a realist or impressionist in his
artistic conception. He thereby comes to fail to recognize
art's capacity to shape individual figures,
to form wholes, where reality gives only fragments
or beginnings. By this capacity to shape given elements
into a whole of individual character, art gives more
than science, which must keep to general concepts,
can give, and to that extent art must be said to stand above science.
Art overcomes chaos better than science
is able to do. --- No one has understood this better
than Plato himself, who in his dialogues has given us so many
% --- p. 46 ---
\setcounter{footnote}{0}%
\opage{46}splendid portrayals of personalities, of scenes, and of
conversations, portrayals that appear with the stamp of plastic wholes.
In speaking against art Plato is then
really speaking against himself. He also declares at the end that
he would be glad if anyone could refute him, since
from his own experience he knows the enchanting power of art.

When Aristotle found the significance of tragedy to lie in its
arousing feelings such as fear and pity in such a
way that we are freed from the torment and bondage that these
feelings bring when it is the real circumstances of life
that call them forth, then people have (Finsler: Platon und die
aristotelische Poetik) perhaps rightly seen in this a direct
objection to Plato's conception of tragedy. Aristotle
has theoretically more sense for the individual than Plato and
has therefore been able to maintain the significance of entering imaginatively into
poetic figures. Such an imaginative entering had for him
a ``purifying'' effect on the mind.

3. The first time Plato in the ``State'' speaks of religion,
it is, as we have seen, on the occasion of poetry, which in
Athens played so great a role in education, since the young
through it received their conceptions of the gods determined.
The criticism of poetry became a criticism of theology. The
weight that he later lays on scientific knowledge leads
to judging from a new side the received religious
conceptions. It is characteristic that the question of
cult, which the history of religion shows to be so decisive
an element in all popular religion, is not discussed further. He
leaves it to ``the interpreter whom we have received from our
fathers, and who, enthroned at the center of the earth, gives answer
to all men.'' It is thus Apollo who is to rule
over cult, and if the lawgivers are wise, they see that they
do not understand how to order cult, but do best to
leave everything pertaining to it to the ancestral god. The main question
becomes how, on the basis of his
philosophy, he must take his stand toward the religious conceptions. He himself
makes use of the conceptions of gods as personal beings,
but there can hardly be any question that his
% --- p. 47 ---
\setcounter{footnote}{0}%
\opage{47}conception of the gods is conditioned by his doctrine of the ``Ideas'',
i.e. of the highest thoughts, which make the things in the
world intelligible to us. The decisive passage is found not in
the ``State'' but in the somewhat later dialogue ``Phaedrus'', where
it is said that whereas we human beings can only in single moments
raise ourselves up to the world of the Ideas, God lives constantly
in this world --- and it is on this that his divinity
rests. The concept ``divine'' thus conditions for Plato
the concept of God. And he knows what the divine is ---
it is those thoughts that in the end make the world intelligible, and
that thereby at the same time raise above everything sensible and perishable.
The Neoplatonist Plotinus alludes to this statement of
Plato and adds that what makes God divine must
be the highest God. In more recent times this would be expressed
by saying that the concept ``divine'' is a concept of value,
and that human beings naturally transfer to their gods the
most valuable qualities they know. We then come to
confirm Goethe's words:

\begin{quote}
Daher der Völker löblicher Gebrauch,\\
dass jeglicher das Beste was er kennt,\\
er Gott, ja seinen Gott benennt.
\end{quote}

If Plato has a theology, then it is in any case
dependent on his philosophy.

But in a certain independence of his conceptions of
the gods stands his belief in immortality. On the great journey
to North Africa and southern Italy that he undertook after the death of Socrates,
and which had so great an influence on his spiritual life,
he learned not only to know the significance of science (especially of mathematics),
but he also entered imaginatively into the
Pythagorean doctrine of the transmigration of souls. Not only in the ``State'',
but both in earlier dialogues (``Gorgias'', ``Meno'',
``Phaedo'') and in later ones (``Phaedrus'') this doctrine comes
forth. It has formed the religious background of his
life of thought. By comparing what he
says in different dialogues on this question, people have even found
that there stood here before Plato's spirit a great imaginative picture,
an ancient Comedia Divina. Plato here appears
% --- p. 48 ---
\setcounter{footnote}{0}%
\opage{48}decidedly as a prophet. At the end of the ``State'' he tells
of a man who fell in battle and came to the realm of souls,
where he became a witness to how the souls, after their judgment
and the reward or punishment that followed, were again sent
down to earth, it then depending on their own choice
into what conditions of life they were placed. The man himself came back to
earth as ``herald of the beyond'' --- just as Orpheus
in an old poem came to the underworld, not, as
in the later myth, to fetch Eurydice, but to proclaim to
men their fate beyond.

The conception of the transmigration of souls and of the souls' return
occurs among some of the most primitive peoples.
Plato could in this respect say, like Myers, the most considerable
representative of modern spiritism, that he acknowledges
his kinship with ``the poor thinkers of the Stone Age''.\footnote[1]{Cf. my characterization of Myers in \textit{Oplevelser og Tydning.} p. 25---33.} But
in the most primitive forms in which it appears,
the conception of the transmigration of souls has really no ethical significance.
Such significance it acquired only in India and
in Greece, when hard fates had broken the original
courage to live.\footnote[2]{Cf. my Religionsfilosofi.\textsuperscript{3} p. 97---100.} In Plato the doctrine of transmigration serves
partly to justify the assumption that the number of souls is
limited (so that the same souls must return if
human life is not to cease), partly to justify the importance
of leading life here in the world in the right way, so
that one may be prepared to make a good choice when
one is to begin life again. This choice is determined, namely,
by the experience of the earlier life, --- and it turned out
that the souls that had fresh memories of their earthly life
made the best choice, whereas those who had been rewarded
with a stay in heaven often chose wrongly! Heaven is thus
no good school. --- Plato's mythology thus has
an ethical core. What matters is to lead life in truth
and harmony and with a manly mind. Plato really
does not get beyond Socrates' undaunted faith that no
% --- p. 49 ---
\setcounter{footnote}{0}%
\opage{49}harm can befall a good man. Just as little as Plato's theology
brings anything new in relation to his philosophy,
since gods as well as men must subordinate themselves to
the supremacy of the Ideas, just as little does the doctrine of transmigration
give him wholly new ethical motives. He stands here much as
Spinoza and Kant later, whose ethics was independent of their
religious conceptions. Plato's position was that justice
can be maintained under outer opposition and without
support from gods and men, but that there is no reason
to believe that under ordinary conditions it would not
find support from both gods and men. Plato approaches
what the Stoics later thought, that the ethical character
in its essence is independent of outer help and of
outer opposition.

\piecerule

% --- p. 50 ---
\setcounter{footnote}{0}%

\opage{50}\essayhead{GIORDANO BRUNO}

The period in Europe's cultural history that is usually called
that of the Renaissance, the time of rebirth, comprises the fourteenth,
fifteenth, and sixteenth centuries. It is opened by the three
great poets, Dante, Boccacio, and Petrarch; after them come
the learned humanists and then in turn the pioneering
natural scientists and philosophers: Pomponazzo and Machiavelli,
Leonardo and Telesio, Bruno, and later Galilei and Campanella.
Here was prepared the world-view and view of life that, in
its main features, can be said to be peculiar to
modern times. Two northern astronomers, Copernicus (about whose
descent Germans and Poles quarrel) and the Dane
Tycho Brahe, gave the foundation, and both were also hailed
with enthusiasm by the man with whom we are to concern ourselves here,
Bruno, the typical representative of the Renaissance in philosophy,
in whom all the threads of that remarkable age are woven together.
Bruno saw in the new astronomy the possibility of a new
world-view, which he presented with enthusiasm,
and whose most important consequences he saw. Supported by one of
science's great turning points, his thought burst
the narrow bounds of what had been handed down from the past, and pointed to new problems,
whose solution was to become the task of the future. It belongs
to Bruno's greatness that he is clear that the work
of winning truth can never be completed.

Bruno was born in 1548 in Nola in southern Italy. His baptismal name
was Filippo, but when at the age of 15 he entered a Dominican monastery
in Naples, he received the name Giordano (after Dominic's
nearest successor). It was a fateful
step he took here, determining the whole
% --- p. 51 ---
\setcounter{footnote}{0}%
\opage{51}course of his life. Probably it was predominantly a desire to
study in peace that led him to the monastery, but it was
the provisional result of these studies that, when he was
28 years old, made him flee from the monastery. A
charge of no fewer than 130 heresies hung over his head.
Doubt about the dogma of the Trinity was the most incriminating;
but he was also reproached for wanting no images of saints
in his cell, only a crucifix. He
fled up to northern Italy, and here he lived for some years
as a tutor of boys and as a teacher of astronomy to young noblemen.
It was probably in these years that his world-view
was founded through a repeated study
of Copernicus' famous work (which had appeared in 1543),
whose ideas he had not been able to accept at first reading.
In his later writings he praises him highly: ``You
magnanimous Copernicus, whose work already made an impression
on me in my very young years, although I thought it conflicted
with both sensation and reason!''

Let us pause here a little. Why did Copernicus' doctrine conflict
of the motion of the earth ``against sensation and reason''?

Well, the senses show us, after all, that it is the sun and the stars
that move, while the earth we live on stands still!
--- Copernicus' answer to this was that we have no right to
regard our standpoint and point of view, the earthly one, as the
only possible one, and he shows by mathematical means that for an
observer on the sun the earth would move, like all the
planets, while the sun would stand still. And at the same time he shows
that this point of view leads to the simplest conception,
free of all the auxiliary hypotheses that in the course of time had
had to be devised in order to carry through the earthly standpoint
as the only right one. We then stand --- as
the German philosopher Cusanus had already said before Copernicus
--- to the stars as seafarers who sail past
coasts. Copernicus already builds on the great law of relativity,
which right up to the very newest time plays so great a
role in science. --- This Copernican conception
Bruno has not only made his own, but he has, with
% --- p. 52 ---
\setcounter{footnote}{0}%
\opage{52}boldness, gone great strides further, especially on two points. Copernicus
regarded the heaven of fixed stars (the ``eighth sphere'') as
immovable and held that all fixed stars were fastened in it, all
thus at an equal distance from the earth. And he still assumed the
old doctrine that the planets sat each in its own rotating crystal sphere,
at different distances from the earth. Bruno shows that we have no
right to set limits to cosmic space and to assume
that all fixed stars are equally far from us. Distant ships out
at sea seem, after all, to be at the same distance from us! It is
a point that needs closer investigation; but when
one blindly believes in the old conception, one never
comes to investigate it! --- One sees that Bruno's bold
thinking expressly demands that confirmation be sought
through experience. --- As regards the planetary spheres, his
doubt found confirmation in Tycho Brahe's investigations of
the motions of comets. According to these investigations the
comets had to pass straight through the planetary spheres and long
since have smashed these crystal heavens to pieces. Bruno therefore
praises in his later writings Tycho Brahe as the
foremost of the astronomers of the time (Ticho Danus, nobilissimus
atqve princeps astronomorum nostri temporis).

So much for the senses. When ``reason'' had
made difficulties, it was because the distances within
the world-system were now becoming so immensely great that thought did not
seem able to get a unified whole out of it. The length of the
earth's diameter was to be vanishing in relation to the earth's
distance from the sun! Bruno shows that reason in reality
has no need of a closed picture of the world.
On the contrary, when our reason is awake, one
thought leads over into the next, just as in the series of numbers we can always
add number to number without limit. The burden of proof rests on the one
who assumes that the world is limited. Should nature be
poorer than our thought? And why should there be absolute
rest anywhere? Why should the earth and the fixed stars
be exceptions from the constant motion that experience
everywhere shows us when we look closely? --- Already as a child
Bruno had learned to see that the conception depends on the
% --- p. 53 ---
\setcounter{footnote}{0}%
\opage{53}standpoint. From the wine- and forest-clad mountain Cikada near
Nola, Vesuvius looked like a naked mass that lay in mist;
but when he once came to Vesuvius, it was the reverse, Cicada
that lay naked and misty. From this experience he understood
Copernicus, indeed understood also that one could not stop
where he stopped. If the earth is a planet, then we are
ourselves in heaven. And when one is perplexed as to where
God's abode is to be, since the firm heaven that hitherto stood
as his abode is no longer there, then Bruno teaches that
the Deity is the infinite force or soul that moves
in All, also in man's soul, indeed even farther in, in
the innermost of our soul, than we ourselves can penetrate! No thought
can encompass the essence of the Deity, since on closer reflection
upon it we are led out into the infinite. The significance of religion
is practical, moral. The Bible does not wish to teach us astronomy
and philosophy, but to work upon the minds and wills of men,
and the religious conceptions have only symbolic significance.
Bruno does not teach that what is true in philosophy can
be false in theology; but he teaches that one and the same
truth is expressed in two ways, in the thoughts of philosophy
and in the images of religion. ---

We left Bruno in northern Italy. For one reason or another
he traveled on northward and came to Geneva. Here he laid aside
the monk's habit and now appeared in doublet and with sword.
He seems to have gone over to Protestantism, though
he denied it later; at least he was excommunicated
by the Protestant clergy and has made a public
apology. From this time on he was unfavorably disposed toward
Protestantism. He found there many popes instead of
one, and he missed the great care for the sick and the poor
that he continued to admire in Catholicism. This made
him homeless. He now roams about in France, England, and
Germany, without our always knowing the reasons why he
changes place. The reason has often been wars and political
disturbances, but just as often his own restless mind, his
demanding behavior, and the remoteness of his train of thought from
the scientific and ecclesiastical conceptions of the time.

% --- p. 54 ---
\setcounter{footnote}{0}%
\opage{54}He lectured as a professor in Toulouse and later in Paris. A
great role in his teaching was played by a theory of the art of memory,
about which great expectations were then entertained,
and which he used as a psychological introduction to
philosophy proper. The French king Henry III had
him summoned and seems to have thought that he was in
possession of a secret wisdom that could solve all riddles. From both
Toulouse and Paris Bruno entered into negotiations
about returning to Italy in reconciliation with the Church.
He thought that the Church could content itself with his
symbolic interpretation of its dogmas. But the answer given
him was that he should return at once to his monastery; then
one could talk about the rest.

The stay in London (from 1583 to 1585) was Bruno's happiest
time. He lived in the house of the French ambassador and moved
in witty and distinguished circles. His most profound
and brilliant works he composed in these years.
Whereas before and afterward he wrote in Latin, the six books
that he composed in London are written in Italian. He defended
the Copernican theory and developed the philosophical
ideas that the great expansion of the picture of the world called forth
in him. In a narrower circle he found interest for his
ideas; Queen Elizabeth also valued conversations
with him, and he praises her in lofty tones as Amphitrite,
the queen of the sea, or as Diana, the royal virgin.
But among the learned and in larger circles he found no understanding,
and he gives drastic descriptions of the learned
stupidity and self-importance he met on his way. England was
on the whole too barbarous for him in manners and way of life. Already in
France he felt himself in a half-barbarous country, but in England
he received a stronger impression in that direction.

When the French ambassador was recalled, Bruno went
with him. In Paris he held a great disputation on his theories.
A single proposition among those he set up is of interest
for what at the moment interests natural scientists
and philosophers. He points out that when time
is measured by the motions of the heavenly bodies, and when these
% --- p. 55 ---
\setcounter{footnote}{0}%
\opage{55}appear different, both as regards speed and direction,
from each individual globe, then each star gets its
own time, different from that of the other stars.

After having roamed for several years from one German
university to another, he settled for a time in Frankfurt,
where he composed a whole series of writings in Latin,
in which he gave a more systematic presentation of his
doctrine than in the half-poetic Italian writings. In the midst of
this activity he suddenly breaks off and goes to
Italy (1591). People in Frankfurt seem to have been uneasy
at having this heretic in their midst, and at the same time
a young Venetian, Mocenigo, who was interested in his
writings, had invited him to live in his house. It was
longing for the fatherland, weariness of the barbarism and
the unreceptiveness of the northern countries to the new thoughts,
that led him back. And he had not given up the hope
of reconciliation with the old Church, without needing
to give up his conviction. It was to prove a
great disappointment. The reaction in Italy was in full swing, and the Inquisition
was vigilant.

Mocenigo was disappointed in his teacher and suspected him
of keeping his most important thoughts hidden from him. When he
realized that his guest was an arch-heretic, he reported
him, on the advice of his confessor, to the Inquisition, and
Bruno was imprisoned by it (1592). For a year he sat imprisoned
in Venice. In the interrogations he at first appeared with great
boldness, told the story of his life, and developed his philosophical
ideas. He held that he had attacked the doctrine of the Church
only indirectly, although he admitted that he had not
observed the customs, and he explained that in the doctrine of the Church
he saw a symbolic expression of what he himself by another way
had arrived at. But people were not satisfied with this, and
his hope of an accommodation with the Church burst. Then he collapsed
and declared himself willing to make a complete recantation.
What the outcome of the trial in Venice would have been
cannot be known. But now the Roman Inquisition stepped
in, and in order to please the Pope the government of Venice
% --- p. 56 ---
\setcounter{footnote}{0}%
\opage{56}handed him over to Rome (1593). An entirely new trial began.
His writings were studied through. It seems that it
was hoped to strike a great blow by getting this great heretic
to deny expressly his scientific doctrines,
not only his theological heresies. Among his
judges were the zealous persecutor of heretics Cardinal Sanseverino,
who had celebrated St. Bartholomew's Night as a day of rejoicing,
and Bellarmin, who later also helped to condemn Galilei,
and who perhaps precisely at Bruno's trial had his eyes opened to
the dangerous consequences of Copernicus' hitherto unnoticed
doctrine. What happened in the seven years in which Bruno sat
imprisoned in Rome, one does not know; the records are said to have disappeared
from the Vatican archives. Only the very last proceedings
are known, since a German scholar copied them out while
Pius IX in 1849 was for a time driven out of Rome. When we again
hear of Bruno, he stands firm, refuses to retract anything.
On the 9th of February 1600 he was then degraded,
excommunicated, and handed over to the secular authority
(with the customary hypocritical request to avoid
bloodshed). On this occasion Bruno said: ``You
who condemn me are more afraid than I''. On the 19th of
February he was burned, and his ashes thrown into the Tiber.

A German scholar who had gone over to Catholicism was
present at the execution and has described it in a letter.
He ends his description with the words: ``Now he can go
to one of the other worlds he believed in, and tell
how people are treated in Rome who are like him!'' ---
It seems to have been the possibility of other inhabited worlds
than our earth that especially attracted notice, because God
could become man only once. In the Middle Ages
for the same reason belief in antipodes was a heresy.

After having followed Bruno on his course of life and learned
to know some of his most important thoughts, we shall now try
a characterization of his personality.

According to the interrogations he was a small slender man with a
% --- p. 57 ---
\setcounter{footnote}{0}%
\opage{57}frail body, pale face, a pensive and lively, but
yet melancholy expression. --- From his writings, with their direct
and indirect self-characterization, we get the impression of
a character full of contrasts. Strong sensual passions
dwelt alongside high intellectual striving; pronounced
self-esteem, which could often find bombastic expression,
did not exclude self-forgetting enthusiasm and absorption
in the great ideas, the infinite horizons that had opened
before him. This last trait he illustrates in several places
in pictorial fashion. As Actaeon boldly went hunting
with his hounds, but before evening himself became the hounds'
prey, so it may go with the thinker when he delights
in his thoughts; they can at last gain such power over
him that all regard for himself vanishes. A boy in
playful high spirits looses a boat on the shore, but the current can
carry him out onto the great sea, and he must learn from experience that there
are greater powers than those he can handle in play. Through
other images the conflicting feelings are made vivid
in him who strives for knowledge: at one moment he feels
himself small and wretched before the high object of his knowledge,
feeling the weakness of his thought; at the next he lifts
his head again in heroic striving, which carries him out
over all hindering barriers. The heroic soul would rather
sink in the struggle, or at least feel the pain of its own imperfection,
than lower its sights from the ideal goal. Bruno has
on the whole a sense for states of the soul that are determined by
strong contrasts. Here he stands --- as in his own way
Shakespeare does --- as a modern psychologist who does not, like the
ancient thinkers, dwell on single clear traits of character or
fear that harmony of the soul could not be preserved
despite strong, conflicting contrasts. It was not
only the outer world that became larger in the Renaissance; the
inner world too became richer and deeper.

When Bruno only in broad outline gave a sketch of the picture of the world
that in the following centuries was carried through in
% --- p. 58 ---
\setcounter{footnote}{0}%
\opage{58}detail, this is connected with the fact that his imagination is
developed quite as strongly as his thinking. And yet
he is clear that experience and experiment must step in, in order to
confirm his bold thoughts. He wishes precisely to open eyes
to the necessity of new investigations. He is also
clear that even the boldest and most earnest research
will not be able to exhaust existence.

\piecerule

% --- p. 59 ---
\setcounter{footnote}{0}%

\opage{59}\essayhead{SPINOZA}

\essaysource{24 NOVEMBER 1632 --- 21 FEBRUARY 1677}

\romanhead{I}

In the lectures Henrik Steffens gave at the University
of Copenhagen in 1802, and which became so important for
Danish spiritual life, he gave an enthusiastic description of Spinoza's
life and character and then declared: ``Never can I imagine
this glorious man in his calm solitude, with a
light in his heart that no accident could drive away, with
an inward peace that no persecution could disturb,
without a ray of his light awakening the fairest hope
in me.'' --- What Steffens has here declared, all will agree on,
however they may stand toward Spinoza's doctrine. The
inward absorption in the work of thought, --- the fidelity
with which the truth won is presented without regard
to the opinions of the multitude and of the mighty, --- the purity of
mind that placed the spiritual good, which free knowledge is, above
the outer goods that so easily draw many to themselves, --- all
this will awaken admiration everywhere where preconceived opinions
do not rule in the judging of the inner
life of human beings. But it is not enough to explain the enthusiasm
with which so many from all sides in the world of the spirit
celebrate every day of commemoration that leads thought back to him.
On the 21st of February 1677 he died in his humble dwelling in The Hague,
and when the memory of him is on that occasion brought forth around
the whole world, it is not merely on account of his pure,
ideal, and disinterested disposition. When the international society
Societas Spinozana, founded some years ago,
buys the house in which he lived at his death, in order to make it
a memorial to him beside the statue that
% --- p. 60 ---
\setcounter{footnote}{0}%
\opage{60}was erected opposite it fifty years ago, it is because
there has come forth from his solitude one of the foremost works
of human thought, his ``Ethica'', which in brief and forceful
form sets forth ideas in which the basic thoughts of the new philosophy and
of the new natural science were interwoven.

Historically Spinoza stands as the completer of the thought
of the Renaissance. He unites Giordano Bruno's enthusiasm for nature
and great horizon with the demand for strict grounding
of each particular, such as Galilei, Kepler, and Descartes
had asserted. The unity in nature, which for the Renaissance
became one with the Deity, had according to Spinoza its ground
in the inner connection that science
seeks to demonstrate among all the individual things it investigates.
It is no explanation of a phenomenon or an event
--- Spinoza maintains --- to find their cause in God,
for God is assumed to be the cause of all things, and no explanation
of the individual thing can therefore be immediately derived
from him. The individual thing or event must on the contrary
be shown as a member in a series, within which every
member presupposes the preceding members. The explanation of a
particular then comes to lie in its relation to other
particulars. Where such a series can be formed, a law
makes itself known. And through the interconnection of the laws
one arrives at the idea of a Deity that works in the interior
of things through the lawfulness that conditions their
essence and their existence. What with the earlier thinkers of
the Renaissance was an enthusiastic vision, became for Spinoza
through the study of the new natural science, which was
the seventeenth century's great gift to mankind,
a thought that could be laid as a foundation for scientific
work. Spinoza states outright: ``The more we understand
the individual things, the more we understand God'', --- a proposition
that was Goethe's, that great Spinoza-worshipper's, favorite proposition.
What Spinoza states as a member of his whole system
recalls, in part in the wording, propositions in Kepler and
Galilei. --- Kepler had (in his famous work on the planet
Mars) demanded that the causes appealed to in % --- p. 61 ---
\setcounter{footnote}{0}%
\opage{61}science should themselves be demonstrated in experience, not merely be
objects of faith or speculation. He thereby corrected
his own earlier mystical explanations. The expression
Kepler uses of a cause demonstrable in experience (vera
causa) also occurs in Spinoza. It was at that
time a technical expression, and is still used by Newton. It
states the principle of natural causes. --- Galilei maintained
(in his dialogues on astronomical systems of the world) against
scholastic opponents that it could not avail to appeal
to God as the cause of a single particular phenomenon,
whose interpretation the astronomers disagreed about, for God
in his omnipotence could just as easily let the sun move around
the earth as the reverse. And in analogy with this Spinoza teaches
that in science one cannot appeal to a being
that is assumed to lie at the ground of All, as the cause of a
single particular thing. Each single thing is to be explained by
its connection with another thing, this again by
its connection with a third thing, and so on to
infinity. (Ethica I, 28).

As regards the causal relation itself, Spinoza maintained, in
agreement with Descartes, that there cannot be more in the effect
than there was in the cause. If there were more, there would
be something for which the cause had not been found!
For Spinoza the causal relation stands as a relation of identity or
equivalence between different phenomena. This
conception was on the whole characteristic of the thinkers of the
seventeenth century. It is developed more closely by Huyghens
and Leibniz, and in the nineteenth century it became of importance
for the discovery of the law of the conservation of energy.
Robert Mayer, the first in the series of men who each
by his own way discovered this law, is said constantly to have had the proposition
causa æqvat effectum on his lips; it stood in his
eyes as self-evident, and it guided him in his decisive
observations and experiments. The other natural scientists too,
who independently discovered the conservation of energy, started from
the thought of a relation of equivalence between cause and effect
as self-evident. But the self-evidence dates
% --- p. 62 ---
\setcounter{footnote}{0}%
\opage{62}only from the seventeenth century, and Spinoza is the one who has
stated it most clearly.

The fusion of the mysticism of the Renaissance and the new
natural science became for Spinoza the foundation of a doctrine of
an eternal subsistence, an absolute substance, which he called
God or Nature, which lies at the ground of All and reveals
itself both in the inner force with which each individual
being seeks to assert its subsistence, and in the bond that
ties all beings in the world together into one great whole.
Our knowledge begins, to be sure, with the collecting of
single observations, but through comparison among
these and through determination of their mutual connection
the inner ground of All can be found, and the culminating point
of all knowledge was for Spinoza a vision of the particular
in its rational connection with the whole, --- a concluding
intuition. The particular in its isolation, without connection
with the whole of nature, was for him only an abstraction.
But the concluding intuition is reached only through strict
work with the particulars, and Spinoza himself declares that
it is only very little he has managed to understand by this
way\footnote[1]{In \textit{Den nyere Filosofis Historie} I have given a fuller characterization of Spinoza's philosophy, and in \textit{Spinoza's Ethica. Analyse og Karakteristik} (Videnskabernes Selskabs Skrifter. 1918) I have gone through his principal work.}.

\romanhead{II}

But how did Spinoza arrive at this philosophy? --- This
question leads to bringing forward his biography and his
personal course of development.

Spinoza was born in Amsterdam on the 24th of November 1632.
His parents were Portuguese Jews, belonging to the so-called
Marannos, who under Spanish rule had
been forced to give up their Jewish culture and their
Hebrew language. His parents, like so many other
Marannos, had fled to Amsterdam, where their practice of religion was at first overlooked
and later tolerated. Spinoza's first name was
% --- p. 63 ---
\setcounter{footnote}{0}%
\opage{63}Bento; with this Portuguese name he signed
himself on legal documents. As a student at the Talmud school
in Amsterdam he was called Baruch, and in letters and in one
of his writings he called himself Benedictus (like Bento
corresponding to the Hebrew Baruch). He studied with zeal
the Hebrew religious books and the Talmud, and people expected
an outstanding Talmud teacher of him. But gradually
he aroused a painful attention among his co-religionists
by his historical and philosophical criticism, and at last he was,
at the age of 24, solemnly expelled from the synagogue. Only now
did he truly make acquaintance with the new science
and philosophy of the West, which, as we have seen, have left deep traces in his
own philosophy. When the Protestant clergy too were scandalized
by his criticism of the Old Testament, the
authorities were induced to expel him from Amsterdam. Afterward he lived
in various places in the country and supported himself by
grinding optical glasses, in which he attained great skill. At last
he took up residence in The Hague, and here he died of a chest disease
on the 21st of February 1677. His pure and harmonious
conduct of life under humble circumstances, his disinterestedness, and
not least his calm death were for the theology of the time a great
problem. Could a man with his views live a decent
life and die in peace? There were many rumors in circulation
about his dissolute life and about his death in despair.
But a worthy pastor who came to live in Spinoza's lodgings
and with the same landlord and landlady, set himself in a biography of the philosopher
to refute these rumors. It is a quite small Dutch
genre picture one gets from his description, which was based
on accounts from Spinoza's landlord and landlady,
who were very fond of their lodger. Spinoza often came down
to them and conversed with them over a pipe of tobacco. Even
on the morning of the day that was to be his last,
he was down with them and conversed with them, and they were
greatly astonished when they came home from church in the afternoon
and found him dead. He had died in peace, in the presence
of one of his friends, who was a physician. The good pastor's testimony
is all the more important in that he nourished the
% --- p. 64 ---
\setcounter{footnote}{0}%
\opage{64}greatest abhorrence of Spinoza's doctrine and was much scandalized
that bills came to the estate that spoke of him as
``the late Mr. Spinoza''.

Let us return to the time right after the expulsion from
the synagogue. We know very little about the years that passed from
this point until his correspondence with various
friends began (1656---1660). It was, however, partly the time
when his philosophical ideas matured, partly a period of ferment,
as regards his personal way of taking life. In
both ways was prepared the wonderful harmony between
doctrine and life that places him in the first rank of those who have
labored for free research and ideal conduct of life.

In this period Spinoza made the transition from
theology to philosophy. As a rabbinic theologian he had found in
the concept of God, as the Hebrew religion conceived
it, a principle of understanding. God contained
for him the explanation of all things. The change that now
took place in him can be expressed by saying that whereas God had formerly
been for him the principle of understanding, it was now the principle of
understanding that became God for him. The transition
came about through his asking what it really is to understand
something. The answer was found in the new science and philosophy
with which he now made acquaintance. The point was to find such
causes of things as themselves made themselves known in experience.
It was the principles of Kepler, Galilei, and Descartes
that he took up into his philosophy. In the presuppositions of such
knowledge of causes the principle of all understanding is given,
and this principle essentially determined his concept of God. He
certainly felt the oppositions between the two concepts of God,
between which the transition took place, very sharply. But
although in a special work (Tractatus theologico-politicus)
he set himself to impress the difference between theology
and philosophy, he nevertheless thought himself entitled to retain the
same word, since he now uses the word God synonymously with
the words substance and nature. He thought he had a concept
of God, even if he could not have any corresponding
image. His concept of God was both the principle of
% --- p. 65 ---
\setcounter{footnote}{0}%
\opage{65}the subsistence of existence and of its intelligibility. It
filled his mind with the highest joy that existence was
to so high a degree accessible to scientific understanding.
The summit of the spiritual life was for him what he called the
intellectual love of God (amor intellectualis dei),
which was called forth by the thought that true knowledge was
possible in this world, which has its ground in God as the one who
is one with nature. To philosophize was for him a religious
act.

But the same years of transition, in which these ideas matured
in him, were at the same time occupied by an inner struggle of
another kind than the intellectual struggles, which yet also
led him to see in the work of thought the highest spiritual
good. By his own confession there was a time
when partly the craving for sensual pleasure, partly the craving for
honor and riches stirred strongly in him. He felt
himself drawn in different directions. But he experienced during
these struggles that peace came to the mind when he only
could collect himself in earnest reflection (modo possem penitus
deliberare!). In his own personal life he has thus
experienced the significance of understanding, and the joy of knowledge he
has come to know both through the conflict of thoughts and through that
of the drives.

It has been held that his upbringing in the religion of his people
was decisive for him in his whole life, and that it was
the echo of his religious childhood that later
worked to give his conduct of life a stamp of purity and
loftiness. But there can be no doubt that the life
in the world of thought, to which he rose through inner struggles,
has had a decisive influence on his life and
shown its sustaining power in it. When a seafarer comes from
the southern hemisphere up across the equator, it is new stars
that guide him, even if he (perhaps with piety) remembers
that other stars formerly determined his course. --- I
cannot here help thinking of my teacher Professor
Hans Brøchner (who incidentally wrote a thorough monograph
on Spinoza in 1857). He too stands as a fine
% --- p. 66 ---
\setcounter{footnote}{0}%
\opage{66}testimony that a noble human life can be lived in
the light of the new stars that the work of thought has led to discover.

\romanhead{III}

Put to a test by the epistemology of our time,
Spinoza's philosophy stands nearest as a presentation, given in dogmatic form,
of the ideal of scientific knowledge, not as
a conclusive solution of problems. But whoever wishes to immerse
himself in his work will notice, beneath the firm, dogmatic formulation,
a working thought. Therefore Lessing,
one of the first who understood Spinoza's significance, could declare
that if he were to call himself after anyone, he would
call himself after Spinoza, although he himself had uttered the
famous proposition that not the possession of truth, but
the constant striving after truth was the highest we human beings
can attain. Spinoza himself has said that many riddles
remain, even if one enters into his philosophy. He
thought, after all, that he had understood only very few subjects in
a perfect way. And he could not imagine that the fullness
of existence could be exhausted by the only two forms
of existence that our experience shows us, the mental and the
material, whose mutual relation presents so many
difficulties. There must, he held, be countless other
forms of existence, even if the knowledge we are able
to win of the forms of existence accessible to us is
enough for us to come into inward connection with the
innermost ground of existence.

It was the life of thought that for Spinoza was the highest, that
which gave value to all other sides of life. From his youth
on he had, after all, experienced what significance the life of thought can
acquire when it goes into depth. And he stands as a great model
through his inward dedication to intellectual work.
There are other sides of life that in our eyes do not
come into their own in him. Thus art, poetry, and
the free life of feeling, which in its deepest and most inward forms
may be called religion.

% --- p. 67 ---
\setcounter{footnote}{0}%
\opage{67}Nevertheless he had the capacity and the desire to observe human life,
as it showed itself in him and around him. Like
a whole series of the thinkers of the Renaissance he found in
the urge of self-preservation, in the urge to hold fast and deepen
his being, the foundation of all human mental life.
In all beings there works, according to Spinoza's faith, the infinite being
subsisting in itself and determined by itself, which he
calls God or Nature, and it is a ray from the fullness of this
divine being that stirs in the urge of finite
beings and leads to self-assertion. Pleasure and displeasure,
joy and sorrow, hope and fear, love and hate are different
forms of expression of the urge of self-preservation that
stirs in us all. And Spinoza gives a description, masterly not only for his
time, of the development of these feelings. Man
can be tossed about by the passions like a
ship on a stormy sea, if he does not seek a clear understanding
of them and their causes. Through such self-understanding
there can develop a strength of spirit, a courage
(fortitudo) toward life, which appears partly as
courage to live (animositas), partly as that magnanimity
(generositas) that takes up the burdens of others besides its own. In these different
forms true freedom expresses itself, which is not a
mystical faculty but a trait of character that presupposes
that man has inner conditions for subsisting and unfolding
himself. When self-understanding becomes scientific knowledge,
the foundation of true freedom is laid so much
the more firmly.

Spinoza had an understanding, rare for his time, of
religious standpoints other than his own. Through his
historical study of the religion of his people he became a precursor
of modern research in the history of religion. He saw the
significance of the positive religions in this, that in their conception of God
they held up an ideal of life (exemplar vitæ), in obedience
to which a good and just life could be led.
What was decisive with regard to religious conceptions for Spinoza,
who in this stands as a precursor of men like Carlyle
and Kierkegaard, was that the single individual with full and honest
% --- p. 68 ---
\setcounter{footnote}{0}%
\opage{68}adherence (sine ulla hæsitatione, sed integro animi consensu)
follows the ideal of life he acknowledges. It was, according to
his biographer's statement, in this spirit that Spinoza conversed
with his landlord and landlady, when he discussed with them what
their pastor had said in his sermons. --- In a special
work (the Tractatus Theologico-politicus mentioned above)
he seeks to show that when religion is conceived in this
way, the state will be able without danger to grant full freedom of religion and
of thought. Here as on other points he has uttered
ideas that point far beyond his time.

\romanhead{IV}

In spite of all, Spinoza is a child of his time. His thinking
comprehends in concentrated form problems that now must be treated
one by one, so that afterward it can be tested
whether and how the solutions of the individual problems
can be brought into agreement with one another. It
is the division of labor in the intellectual domain, and in
the domain of spiritual life generally, that makes it necessary for us now to philosophize
differently from Spinoza. But it will always lead to
enrichment and consolidation to go back to the quiet thinker of
the seventeenth century, the time when modern science
was founded. And it strengthens faith in the sustaining power
and significance of the work of thought when one dwells on the picture
of him in his simple dwelling, working on his monumental
work amid misunderstanding, indifference, and
fanatical persecution. The sacred fire burned in his interior
and made him forget the cold and the darkness outside. The spirit
that animated him must not vanish from the world of the spirit,
however much both the tasks and the means of their
solution change.

What persists in the world of thought as in the world of nature,
is not the forms that at a given time were at disposal,
but the forces that lay at the ground of the origin of these forms,
and which for a time, but also only for a time,
% --- p. 69 ---
\setcounter{footnote}{0}%
\opage{69}laid themselves to rest in them. To them we seek back in order to
strengthen our hope for the future. Even if we nourish a firm
faith that spiritual values will ever anew arise,
we cannot decide in advance of what kind these values
will be. But we can get an inkling of it by making
ourselves familiar with the spirit in which the great ones of the past worked.

\piecerule

% --- p. 70 ---
\setcounter{footnote}{0}%

\opage{70}\essayhead[Pascal and Kierkegaard]{PASCAL AND KIERKEGAARD\footnote[1]{On the occasion of the three-hundredth anniversary of Pascal's birth (19 June 1923) the Revue de Métaphysique et de Morale published a Pascal number, for which, at the request of the editors, I wrote this comparison between the two religious thinkers, which appeared at the same time in Danish (in ``Tilskueren'').}}

Two centuries apart, two thinkers met, both of whom clung to
Christianity with the whole inwardness of their souls, in the manner in which they
reckoned with their religion, as it had developed through the ages.
Both maintained that there was an irreconcilable
opposition between the spirit in which Christianity
entered the world and the Christianity of more recent times.
The opposition had arisen from the attempt
to enter into an intimate relation with ``this world'', that is
to say, with the culture developed on a purely human foundation.
The reckoning was undertaken by both of them with such passion,
and carried through with such consistency, that there
would have been material enough for a great revolution in the world
of the spirit. But a tragedy attaches to them both:
they broke down when they had spoken their last word. The
spiritual struggle — within them and with their surroundings
--- had consumed them. They died young, Pascal at the age of 39
(1662), Kierkegaard at the age of 42 (1855). They
left behind a great riddle, more sharply formulated than ever
before, but unsolved as before. Yet a heroic deed in the world
of the spirit is never performed in vain. They have, with might
and main, put the question whether the religious
traditions of the Christian churches can still be made the basis of thought and
life by sincere and honest seekers after truth.

On the part of the churches no answer has been given to this
% --- p. 71 ---
\setcounter{footnote}{0}%
\opage{71}question. They think they can continue their practice without
taking any notice of it. They rely on what is indeed also
a fact, that no other
organization has so far arisen that is able to awaken larger circles
to a serious conduct of life and at the same time to bring consolation and relief
under the hard conditions of life. But the question is not
thereby removed from the world. It is told of Scipio Africanus
that when he was accused in the popular assembly of
disorder in his accounts, he said: ``On this day
I defeated Hannibal at Zama; --- let us go to the Capitol
and thank the gods!'' The people followed him --- and forgot the rendering of accounts.
The doubt still stands --- and since the time of the two thinkers
has been thoroughly substantiated by
critical investigations --- whether the churches have a right to consider
themselves in continuity with the deep movement of the spirit that
1900 years ago broke forth in a corner of the great world empire,
and whose guiding thought was that life must now be led
in tense expectation of an imminent supernatural
irruption of a new kingdom, an expectation in relation to
which life in family and state, in art and science became
of vanishing value, in so far as it was not a positive
hindrance.

It is not merely for the sake of this reckoning, significant
as it was, that we remember the two men. On their
way to the reckoning they developed thoughts of great importance,
and --- as so often in the realm of thinking --- the way
is no less significant than the goal. On the way they
have cast light on the nature and conditions of the human life of the spirit.
They have made a multitude of significant psychological
observations. They have contributed to what we would now
call epistemology, in that they have tested the power of thought
at the very limit of thought. They have shown that even where feeling
and passion have the floor, it is nevertheless thoughts that, like
guiding stars, determine the direction of the movement of the soul.
They are knights of thought, and their struggle is of interest, whether
one believes that it has led to victory or to defeat.
--- Their paths did not coincide. They are two different
% --- p. 72 ---
\setcounter{footnote}{0}%
\opage{72}personalities that we have before us here, each with his peculiar
inner and outer presuppositions, and the historical conditions
under which their thoughts were brought to a point were highly different,
though analogies are not lacking.

Already in their style they differ. It must here be
remembered that Pascal's ``Pensées'', his most important work, is not
a finished work, but a series of drafts, which he wrote
down in the intervals between his attacks of illness,
often interrupted by them. They are terse, compressed,
vibrating with an emotion that now lured the thoughts
forth, now was called forth by them. There is a unique
force and clarity in these thoughts, and these qualities
would hardly have been lost had the book been
worked out. Kierkegaard's style does not have this striking
force, this epigrammatic character. Authorship was
for him, in his first period at any rate, a cure: as long as
he wrote, melancholy had no power over him. He
therefore took his time. And besides, he loved language
for its own sake; he played with it and did not always think
that the reader perhaps did not have as much time as
he, but longed to be able to survey the course of thought and see
where it leads. Episodes of a lyrical and nature-describing
kind are interwoven, as well as examples from daily life.
The passion that also lies behind his production
is felt more as a constant trembling than as momentary
emotion. Only in the very last period of his authorship
does this trembling become an earthquake. ---

If in this attempt at a comparison I should
come to dwell comparatively much on the Danish thinker,
it is because I am writing for French readers. To wish to
contribute to the understanding of Pascal in a French journal
would be to bring owls to Athens, so rich and thorough is
French Pascal scholarship. I wish here in particular to express
my admiration for Ste Beuve's ``Port Royal'', with which I have now,
after an interval of almost 60 years, renewed my acquaintance,
and which, in spite of the work that has been done in the meantime,
appears with the same freshness and thoroughness as
% --- p. 73 ---
\setcounter{footnote}{0}%
\opage{73}before. I must also mention the magnificent edition of
Pascal's collected works that Léon Brunschwicg has prepared
in collaboration with Pierre Boutroux and Félix
Gazier, and which in its introductions and notes provides many
indispensable pieces of information. In particular Brunschwicg, in the introduction
to the ``Pensées'', gives an extraordinarily interesting account
of the line of thought that Pascal may be assumed to have
intended to carry through in his Apology. In what follows I cite
the ``Pensées'' according to the arrangement of the fragments
proposed by Brunschwicg. --- Kierkegaard's works are at present
appearing in a second collected Danish edition. They are also available
in German (from Diederichs, Jena). The works comprise 14 volumes.
But besides this a complete edition has been begun of
Kierkegaard's journal entries and papers (by P. A.
Heiberg and V. Kuhr). So far 11 volumes have appeared, and there is
material for 7 volumes still. It is thus a very voluminous
production that is in question here. In French there is
a fine essay by Delacroix in the Revue de
Métaphysique et de Morale (1900), and a speech by me
at the University of Copenhagen on the hundredth anniversary of Kierkegaard's birth
(5 May 1913) is printed in the same journal
(1913). I also take the liberty of referring to my monograph
on Kierkegaard as a Philosopher (1892), which has been translated into
German.

Of any influence of Pascal on Kierkegaard there can
hardly be any question. In the catalogue of Kierkegaard's
library Pascal's works are found (in German translation). But
it is only at a late stage of Kierkegaard's development
that Pascal is mentioned in the journals (according to what Mr. Archivist
Heiberg has kindly communicated to me). It was in the autumn
and winter of 1850, and by then he had already published
the writings (``Sygdommen til Døden'' 1849 and ``Indøvelse
i Kristendom'' 1850) whose content could most remind one of
Pascal. He also read various works
on Pascal, in particular Reuchlin's work on Port-Royal, and became
aware of the difference between older editions
of the ``Pensées'' and the edition prepared by Faugère, in which
% --- p. 74 ---
\setcounter{footnote}{0}%
\opage{74}the most revolutionary utterances were first available. (The
German translations that were found in Kierkegaard's library
were made from the older editions.)

--- As I now attempt a comparison between the two
thinkers, I dwell successively on temperament and character,
on intellectual presuppositions, on the relation to
Christianity, and on the question of which paths
stood open when the two men had spoken their last word.

\sectionhead{1. \textit{Temperament and Character.}}

Pascal's character is described by his two sisters (Jacqueline
in her letters, Mme Périer in her biography of her brother)
in agreement. He had a fermenting or
boiling temper (humeur bouillant), and his lively spirit
brought with it such impatience that it was at times
difficult to please him; his eldest sister marveled
that in spite of his fiery temper he had been preserved from the excesses of youth.
It is certainly characteristic of him
when in his Pensées (fr. 139) he thinks he has discovered that
all the misfortune of men comes from one single thing, that
they cannot keep quiet. In the remarkable ``Discours
sur les passions de l'amour'' (of which he is probably
the author) the beauty of the passion of love is praised, precisely
when it expresses itself in an overwhelming manner, and passion
in general is described especially by the fact that all thoughts rush
off in one single direction. There is here a presage of the
passion of faith that seized him after his definitive conversion.
The basic stamp of his character was preserved, however different
the stages of development he passed through.

Although a double conversion was needed before he came
to his last, decisive stage, he was nevertheless never irreligious.
Through his upbringing a serious religious
keynote prevailed in his soul; but alongside it
there stirred an early awakened intellectual passion, which led
him to epoch-making mathematical and physical investigations,
and for a time he also took part in worldly life,
% --- p. 75 ---
\setcounter{footnote}{0}%
\opage{75}a witty and literarily interested social life with gifted
men of the world. In this worldly period he had the
experience that seems to underlie the treatise
on the passion of love. According to Strowski (Pascal et son
temps. II p. 278), the object of his passion seems to
have been a beautiful and distinguished lady, who loved adventurous
and restless moods, but who gave him no occasion
for anything he would have to repent. Yet it was later to
stand before him as a great sin that his mind had been
occupied at all with anything other than God. For Pascal there was
here an opposition as between life and death. He belonged to
those characters who could reach their final stage only through a decisive break; this followed from his ``humeur
bouillant''. The catastrophe came on the night (23 Nov. 1654) when
it dawned on him, in an ecstatic state, that the living
God demands everything, so that an unconditional renunciation had to come
into force. His previous life now stood before him as a life
in emptiness, but he rejoiced that he had now found
the source of life.

Against the background of what he had thus experienced,
the attempts of Jesuit confessors to moderate the demands of the
great ideals according to the needs of worldly people had to offend
him in the highest degree. It was impossible to join God
and the world (joindre dieu au monde) (fr. 935). The world
is quite willing to have religion, but it must be a gentle and mild religion
(fr. 956). In his ``Lettres écrites à un provincial'' and
in the last ``Thoughts'' that he wrote down, he speaks in ever
stronger terms of the pernicious path on which the Church has set
out. The Church is no longer Church if it does not have
God with it when it looses or binds. (Si elle absout ou
si elle lie sans dieu, ce n'est plus l'Église. fr. 870.)

During the struggle into which Pascal thus entered, and while
he was in full swing with the publication of the ``Provincial Letters'',
a young relative of his was cured of an eye disease
that the physicians had declared incurable, by the
diseased eye being touched by a thorn from Jesus' crown of thorns, which belonged
to the relics in the church of Port Royal. This miracle confirmed
% --- p. 76 ---
\setcounter{footnote}{0}%
\opage{76}him in his conviction of being on the right
side. The age of miracles is by no means over, as his opponents
thought. And when the ecclesiastical authorities deny
the truth, God must intervene directly (fr. 832; 839).
Confidently, therefore, Pascal appeals from Rome, which condemns
his writings, to the tribunal of Jesus Christ (fr. 920).

In his own life Pascal carried through the absolute renunciation.
True, he did not enter a monastery, but he lived his last
years in the strictest asceticism, occupied solely with works of charity,
but at the same time anxiously guarding against anyone
of those closest to him attaching their love to his person
and thereby robbing God of the full devotion that was due
to Him. His long illness he felt as a relief,
although it prevented him from completing the apology for
Christianity of which the ``Pensées'' contain drafts. Spiritual
sickness is indeed the natural condition of Christians, and it
is easier to remain in it when one is also physically
ill!

It is a somewhat different portrait of character that presents itself
to us in Kierkegaard. There is greater continuity in his
development, though it exhibits decisive transitions. He
cannot be counted among ``the twice born''. Neither ecstasy
nor miracle plays a role in his life.

Like Pascal he had been brought up in a serious Christianity.
But in his youthful years he was absorbed in poetry
and philosophy. It was the age of Romanticism, and the idea of a harmony
between art, science and religion dominated
minds. Romantic thought had recently been systematized by Hegel, and Kierkegaard was for a time absorbed
in his ideas. If it was the case that all oppositions
finally meet in harmony, one could for the time being
let moods and fantasies freely unfold; they too
belong to the rich, manifold world that at last
could find its explanation in the basic Christian ideas.
Kierkegaard noticed early, however, that a duality was present here.
In his journals from this time he used
Goethe's lines about Gretchen as a kind of motto: ``Halb
% --- p. 77 ---
\setcounter{footnote}{0}%
\opage{77}Kinderspiel, halb Gott im Herzen!'' And there was even a
time when he was drawn into a wild life, which left behind
remorseful memories, memories that had to feed the melancholy
that lay deep in his nature, and which could sometimes overwhelm
him so that he felt himself on the verge of insanity.

He now saw it as his task to form a harmonious
whole out of the oppositions that stirred in him ---
play and seriousness, frivolity and melancholy.
He discovered here by way of his own personal experience
that in practice it was not so easy to reach a ``higher unity''
as it was in Hegel's system. This inner work made
him lonely. He felt isolated. His melancholy made
it (he thought) impossible for him to enter into relations of life
that rested on open and glad devotion. He broke off his
engagement to a young, joyous woman, and he felt himself
unable to take on an official position. There was
for him a glaring opposition between other people's easy
way of taking life and his own struggle with dark
powers. Had he been able to live in the Middle Ages, he would,
he thought, have gone into a monastery. The inner work
of working out his personal life he has beautifully
described in the following manner: ``I sit and listen to the tones
in my inner being, the glad beckoning of music and the deep
seriousness of the organ. To work them together, that is a task, not for
an organist, but for a human being who, in the absence of greater
demands on life, restricts himself to the simple one, of wanting to
understand himself.'' --- Now ``Gott'' and ``Kinderspiel''
no longer dwell peacefully side by side in the heart. The opposition
between them sets a task that can lay claim to
all the energy of life. What appears here is repeated
throughout Kierkegaard's whole life. He did not move
from one extreme to another. The oppositions were always
there, but they sharpen as his development
proceeds, and the work of overcoming them
therefore becomes greater and greater. From his own experience
he could take the standard for the height of the personal life.

When it came to checking melancholy and relieving the
% --- p. 78 ---
\setcounter{footnote}{0}%
\opage{78}inner tension, there was a means at his service.
He needed only to follow his urge to literary production.
It was a liberation to him to devote himself to poetic
and philosophical production. Like the princess in ``Tusind og én
Nat'', I saved, he says, my life by telling stories, that is
to say, by writing. In the first period of his work as an author
(1843---1846) some of his most widely read
writings came into being, for example ``Enten --- Eller'' and ``Stadier paa Livets
Vej''. They concerned what the struggle in his inner being was about,
--- to become a personality, to win a firm center around
which all the elements of life could be ordered. And thereby
these writings could at the same time serve to influence the slack age,
which was the epigonic age of Romanticism and of speculative philosophy and was inclined to volatilize the oppositions of life
and thereby lower the level of the spiritual life. Kierkegaard
did not dare to say that this influence had been
the original intention of his work as an author.
He was honest enough to admit to himself that it had been
a tension in his inner being that needed a diversion. Whereas
Pascal in his ``Lettres à un provincial'' set himself with full awareness
the task of fighting against a lowering
of the ethical level, Kierkegaard came more involuntarily
upon work in this direction. But neither was he
yet, like Pascal, at the height of his religious development.

A new period in Kierkegaard's work as an author (1847
—1855) was called forth precisely by his inner life taking on a
more and more pronounced religious character. ``I feel now,''
he writes in his journal (1847), ``a need to come to myself in a deeper
sense by coming nearer to
God\dots{} Something stirs in me that points to a metamorphosis\dots{}
I shall therefore keep still.'' One sees
here the opposition to Pascal's ecstasy: a half-unconscious
unfolding of a new direction of life, toward which he stands
waiting. The year after he writes: ``Now I am at
faith in the deepest sense''. Christianity now acquired a significance of reality
for him that it had not hitherto
% --- p. 79 ---
\setcounter{footnote}{0}%
\opage{79}had. After concluding the first period of his work as an author
he believed for a while that his literary
activity was now concluded, and he thought of withdrawing
to a solitary region in the country. But the more strongly
awakened religious consciousness then awakened a sharper criticism of
existing Christendom. In his first period as an author
he had aimed at the slackness in the individual
way of taking life; now he found in the Church a compromise,
a softening of Christianity according to worldly
inclinations, while the original ideals were at most
given a place in the background. Christianity was in
his eyes in fact abolished by being made
the religion of mildness and consolation. ``And now there lives in the
existing Christendom a spoiled, proud and yet cowardly, defiant
and yet soft race, which on occasion hears these
mild grounds of consolation set forth, but hardly knows whether
it will make use of them, even when life smiles most beautifully,
and which in the hour of need, when it is seen that they are after all
not really so mild, takes offense.'' This utterance is found in ``Indøvelse
i Kristendom'', which together with ``Sygdommen til Døden'' is
characteristic of the second period of Kierkegaard's authorship.

On the part of the official Church no notice was taken of
the assertion of the sharp opposition between the strict demand of original
Christianity to be wholly absorbed in the
one thing needful and the idyll of modern Christianity built
under the shelter of the dogma of reconciliation. Then Kierkegaard began his
passionate struggle against the existing Church, the most violent
spiritual struggle that has been waged within Danish intellectual life.
His last word was: ``The Christianity of the New Testament
does not exist.'' In the midst of the strife he fell ill and died.

\sectionhead{2. \textit{Intellectual Presuppositions.}}

Pascal cultivated mathematics and physics and has a fine
name within these sciences. From these
studies he drew his conception of what scientific cognition is.
% --- p. 80 ---
\setcounter{footnote}{0}%
\opage{80}A systematics like Descartes', with its attempt to trace knowledge of nature
back to certain simple propositions and thus
to reach a rational foundation for the work of human reason,
he considered impossible. He has aptly been
called an experimental positivist. Already in his ``worldly''
period he realized that strictly mathematical thinking
(esprit de géométrie) cannot, through its deductions,
lead to a full understanding of what goes on in the human
mental life. The mental phenomena are too
manifold, too composite and too changing for that. In the
domain of mental life the point must be to trace the manifold
phenomena back to certain types, to enter into each
of these by itself, to put oneself into their motives and interests.
Mathematical thinking is not suited to this. A fine
sense of types and nuances (esprit de finesse) alone makes
it possible to find the essential elements and to follow them
in their connections. Mme Périer precisely praises her brother for
his ability to put himself into the presuppositions of those
he addressed. He knew, she thinks, ``all the springs
of the human heart''. This sense had been
developed in him in his intercourse with men of the world,
and it served him later, when the point was to win
souls for the true religion. Speculative proofs do not lead
to the goal, he realizes; whether they are tenable or
not, their effect is limited to those moments
of our life when we are purely intellectually interested; afterward
they are easily forgotten. What matters is to take men
where they stand, and then lead them to get further by
their own ability; one is, after all, best convinced by
reasons one finds oneself. One must lead them to see what
strange mixture of greatness and wretchedness their nature
contains. Thought can lead them up into higher regions,
but drives and passions draw them down again into the mire.
Man must be led to see what an incomprehensible
monster (monstre incompréhensible) he is. In the true religion
man will find an explanation of the riddles of his
nature and at the same time means and ways to remove the discord in his
% --- p. 81 ---
\setcounter{footnote}{0}%
\opage{81}inner being. A consideration on which Pascal lays special weight is
that man must feel lonely in the infinite world,
from whose depths no voice is heard that could give understanding
and awaken hope. Pascal himself, through unrest, wonder
and anguish before the great riddles, was driven to seize as the only
way out the one that Christianity opens. Christianity
admittedly contains mysteries alongside the revelations
it brings. But a complete unveiling of the secrets
of existence would exclude the spiritual work,
the struggle and the courage to venture, to stake everything on
one card, which is after all the condition for being worthy
of the highest.

According to Pascal there are two main types, as regards the natural
human beings' relation to the great riddles. The one
type holds to thought and duty, to the great and firm
in man, and regards everything else as indifferent. The
classical representative of this type is the Stoic Epictetus.
The other type takes the swell of life as it comes,
enjoys the single moments and otherwise adheres to custom and
usage, having once and for all become clear about the transience of all things.
It finds a secure repose in the insight that one
basically knows nothing. Montaigne is the classical representative
of this type.

The study of these two types was for Pascal a preliminary exercise for
the vigorous grip with which in his ``Pensées'' he seizes
and shakes the unbeliever in order to bring him to the decisive
point where only one single way out of despair shows itself,
--- the one he himself had personally found. It was for him
a disappointment that there was so little interest in the study of
human nature, its types and their practical
significance. ``When I began the study of man,''
he says, ``I realized that the abstract sciences were not suited
to man, and that I removed myself more from my
true position in the world by immersing myself in them than other
men by being ignorant of them.'' He had then
expected to find many companions in this new study, but discovered
that in this domain there was even greater
% --- p. 82 ---
\setcounter{footnote}{0}%
\opage{82}indifference among people than toward geometry (fr. 144). He himself
now threw himself with passion into this study, while
his earlier scientific interests more and more
vanished. In a letter to the famous mathematician Fermat
he writes, toward the end of his life (10 August
1660), that he still regards geometry as the highest domain
in which spiritual powers can be exercised, but that he
now considers it quite useless, especially as he has entered
upon quite different studies, which are so different from his
earlier studies that he has hardly any memory left
of them.

Yet he constantly maintains that the greatness and nobility of man
rest on thought, and that it is this that we must hold to. We
do not fill space like the great masses around us, and
our force is vanishing in comparison with theirs. We are like
reeds, but we are thinking reeds, and though the external masses may
crush us, they do not reach the dignity that thought
gives. (fr. 346---348). But the thoughts to which Pascal attributes such
great significance were more and more only those thoughts that
serve religious self-understanding. Just as in his ``worldly''
period he had described the passion of love as a
concentration, a rushing (précipitation) of all thoughts
in one single direction, so religion had now become
his passion, and only those thoughts that led in this direction
would he now acknowledge.

Kierkegaard comes with entirely different intellectual presuppositions
than Pascal. Brought up among the epigones of Romanticism,
he regarded for a time Hegel's philosophy as the highest
form of thinking, whereas for Pascal geometry was
the ideal. In a system of thought within which each member
of necessity led over to another, every subject was to
find its place. Poetry, religion and science stood as different
forms or stages of one and the same divine
truth. That was his youthful dream. In association with his
teacher, Professor Poul Møller (the most Danish of Danish
authors), he worked his way out of this dream. His
% --- p. 83 ---
\setcounter{footnote}{0}%
\opage{83}significance as a philosopher rests on his having been one of the first
to criticize speculative philosophy from an epistemological
standpoint. (Uvidenskabelig Efterskrift. 1846).
He shows that even if an ``absolute'' knowledge could be given,
the single individual would still always have to begin the line of thought that
was to lead him up to this high stage from a definite
point within experience, and why more from one point
than from another? And can this factual starting point
itself be derived from the ``absolute'' knowledge that is to be
the result? Will there not moreover always be certain presuppositions
on which even the ``absolute'' knowledge rests, and will one
ever be able to be sure that one has got back to
the ultimate presuppositions? --- Kierkegaard came to see that
the problem of knowledge is the most central of the philosophical
problems. He lays special weight on our living in time.
We live forward, but we understand backward. A conclusion
can therefore not be reached, and Kierkegaard declares himself in agreement with
Lessing in the proposition that beyond the constant striving for
truth human beings do not get.

Like Pascal, Kierkegaard maintains that personal
life needs thoughts, but that these thoughts concern what
is our abode in the course of life. The thought of God is an expedient in need,
not a subject for purely intellectual interest. He
distinguishes sharply between ``essential'' and ``inessential'' knowledge,
and to the inessential he reckons natural science,
philology, history and pure philosophy, --- all thoughts that are
without significance for how we take our stand toward the great
questions of life and death. There is no purely intellectual
way of taking life, unless one wishes to make oneself
into an abstraction. As regards the different ways of taking life
he employs a procedure similar to
Pascal's, in that he tries to characterize certain types.
We can here keep to those he has presented in his most
popular writings (``Enten—Eller'' and ``Stadier paa Livets
Vej''). The aesthetic type makes the enjoyment of the moment the
highest. The art here is to preserve a certain elasticity, so
that one can seize each new moment. The point is to break off
% --- p. 84 ---
\setcounter{footnote}{0}%
\opage{84}in good time so as not to be held fast in a lasting relation of life.
Therefore a certain rotation of crops is to be recommended. The aesthete is like
a tangent to the circle of life (or rather: circles). Only at
a single point does a momentary contact take place; afterward
another point of contact is sought. Hence: fleeting love, but
not marriage; comradeship, but not friendship; fleeting
occupation, but not citizenship; intellectual enjoyment,
but not science! Whereas the aesthete is a tangent,
the ethicist describes a circle about a center and is determined
by the relation to this. It is the type of fidelity.
Ethically, life is lived in fixed relations of living together and in constant work.
Above both types rises the religious life, whose striving
and whole direction is determined by a goal that lies beyond
all moments and beyond all time, incommensurable with
all other human ends, indeed even at its pinnacle,
in Christianity, in the sharpest opposition to
whatever can be conceived in purely human terms as an end or
task. The more pronounced a character the religious life
takes on, the less does man live in his natural element,
but becomes like a fish on dry land. And finally
eternity does not stand (as in Plato and Spinoza) as the great,
ever-present background for time in life, but breaks
in in the midst of time, as the Godhead itself steps forth as
a human being at a definite point of time and in a definite place
in space. Then this revelation, and the relation to
it, becomes decisive for all evaluation of life and its relations.

It is a comparative philosophy of life that Kierkegaard here
has given. It may recall Pascal's characterization of Montaigne
and Epictetus and of the Christian type of life in contrast
to them both. And as Pascal declared sickness
to be the natural condition of Christians, so Kierkegaard comes
to the result that the highest type of life, the Christian,
entails tension and suffering, because the oppositions
between which life is to be led stand in the most cutting
contradiction to one another. And this type of life is precisely
for that reason the highest, because it makes the greatest demands
on the synthetic energy that all personal life requires.

% --- p. 85 ---
\setcounter{footnote}{0}%
\opage{85}According to Kierkegaard it is by a leap that one passes
from one type of life to the other, and the greatest leap is
that which leads a soul placed in despair over to the Christian
type, But not even at the most decisive point does
the leap in Kierkegaard take on the character of ecstasy. And a certain
continuity he does not exclude, since the standard (the energy of
the personal synthesis) is to be common to the whole series
of types. Basically Pascal has indeed the same standard,
that which is taken from the relation between the fullness of life that
is at one's disposal and the force and concentration with which
the full content is ordered. It is now also the only
standard that can be applied to the personal and historical
differences in men's ways of taking life,
and it must be applied also by those who find
the high points of life elsewhere than the two thinkers we
stand before here. --- An objection we must already make here
to Pascal's and Kierkegaard's application of the standard.
They lay exclusive weight on the inner, subjective
resistance that is to be overcome, thus on the inner work
by which unity and concentration are reached in the single individual's
mind. But one must also ask about the significance of this inner work
for larger circles of human life, for the
society in which they live, perhaps for all humanity.
This neither Pascal nor Kierkegaard would enter
into. On the contrary they maintain that there is a yawning abyss between
the type of life that in their eyes is highest and human
cultural work.

\sectionhead{3. \textit{The Problem of Christianity.}}

Pascal appears within Catholicism, Kierkegaard within
Protestantism. This must of course in many
ways influence their course of development as well
as their definitive standpoint. But it is of great
interest for the history of religion to see that at the end of
their striving to penetrate to the core of their religion
they come to stand before one and the same problem,
% --- p. 86 ---
\setcounter{footnote}{0}%
\opage{86}a problem that can be formulated thus: Can the conception of life and the conduct of life of
primitive Christianity be maintained
together with the appropriation of, and work for, the spiritual
and material culture that Christianity has not produced,
but which --- before and after the rise of Christianity
--- has developed according to its own laws? It is, after all, also
a problem that is common to Catholicism and Protestantism.
Before we go further, let us see briefly how it
arose.

The first Christian generations lived in expectation of
the near return of Christ, which was to bring the conclusion of the whole
course of the world through a judgment of the world. In this expectation
consisted the original Christian faith, and both conduct of life
and conception of life were determined through and through by it. In view of
the prospect of the near catastrophe, which meant the breakthrough
of a new heaven and a new earth, all purely human
and social tasks had to lose their significance. Therefore
the New Testament has no talk of art, science
and political life, and even family life is acknowledged only with a certain
reservation. The ethics of the New Testament had to become what
Albert Schweitzer has aptly called an ``interim ethics''.
The point was to gather all attention toward one single
point in the near future. The coming of the Kingdom of God, for which prayer
is made in the ``Our Father'', meant for the first generations something quite
definite, not merely a vague prayer for spiritual goods.
And among the believers the expectation was joined with an
enthusiasm that often rose to ecstasy, and for which every
sacrifice became easy. Everything in ordinary human life became small
in face of the great thing that was awaited.

In contrast to a standpoint such as the primitive Christian one, both
Catholicism and Protestantism appear as organizations
for the long run, built on the prospect of a long
course of development for the religious life as for human life
in general. One became less and less prepared for a swift
departure. Even if one sought to hold fast to the original
expectation, its fulfillment was nevertheless relegated to an ever
more distant future, and the expectation gradually lost
% --- p. 87 ---
\setcounter{footnote}{0}%
\opage{87}its practical, all-determining power. And both the culture handed down from
classical antiquity and the new culture
that developed independently of Christianity, one sought
(after due resistance, often of a bitter kind) gradually
to appropriate. But in spite of everything both churches maintain that they are
in continuity with primitive Christianity and seek, each in
its own way, to justify the claim to hold fast to the
old ideals in the new times.

The most grandiose solution is that of Catholicism. It rests
on a distinction between degrees of perfection.
The monk and the nun give up participation in the life of the world in order
wholly and undividedly to live for the Beyond and in the one thing needful.
Through them the connection with the primeval age of Christianity
is maintained. The laity then carry human
life further under the guidance of the confessors, through
whom the Church exercises a controlling influence on the
single individuals. In the prelates of the Church Augustine (De
civitate dei XX, 9) already saw a partial fulfillment of the Apocalypse's
prophetic description of those who sit on the seat of judgment
in the millennial kingdom.

Pascal was clear that at all times in the Church there had to
prevail a practical acknowledgment of the primitive Christian remoteness
from all cultural life. He lived in his last period as a
monk, even though he did not enter a monastery. The time until Jesus'
return had to be a time of tension and suffering.
The scene in Gethsemane was in his eyes the expression of the constant
situation of true Christians. One must live in contemporaneity with
the suffering Savior. The only time Jesus in his anguish
sought help and consolation from human beings was in Gethsemane. Then
the apostles failed him. But Christians still have the
duty that was then neglected. Jesus' mortal agony lasts to the end of the
world, and the faithful must keep awake until then.
(Jésus sera en agonie jusqu'à la fin du monde; il ne faut
pas dormir pendant ce temps-là. fr. 583). It is thus not,
as most often in the Church, the victorious, but the suffering Jesus
to whom the gaze is to be directed. The Garden of Gethsemane is the
% --- p. 88 ---
\setcounter{footnote}{0}%
\opage{88}counterpart of the Garden of Eden. Fallen man has to dwell
in Gethsemane and can only thereby have prospect
that Eden will be opened to him.

In the practice of Jesuit confessors Pascal now saw a criminal
adaptation of Christian life to the life of the world,
and he directed a sharp protest against it in his ``Lettres à
un provincial''. Despite errors in detail he is
right from a strictly Christian standpoint. It is a different
ethics that is proclaimed in the New Testament and in Jesuit
casuistry, and even from a purely human standpoint
there could be misgivings about the results of this casuistry.
But besides the struggle that Pascal in the Provincial Letters
waged against the mild religion that had gradually replaced
the Gethsemane mood, he has in a special treatise
made a comparison between the Christianity of the first times
and that of our time. When the distance here is so great, he finds
the cause of it in this, that now baptism precedes instruction
in what Christianity is, whereas originally no one
went over to Christianity without being thoroughly informed
about its doctrine and without having renounced
the world. At that time one could tell the difference between the world
and the Church, but now the Church, through infant baptism, has made it
possible for worldly life to spread within it. It
was with a good intention, Pascal adds, that the Church introduced
infant baptism, namely out of fear of excluding the children
who died before baptism from the Savior. But what has here been
done out of care for the children has become corruption for the
grown-ups! ---

It was, as we have already seen above, not only on
this point that Pascal's break with the Church took place. He was led
to reject the authority of the Church, both as regards doctrines
and as regards the establishment of facts. Together
with his Jansenist friends he maintained that the propositions
that the Pope had found in Jansenius' work ``Augustinus''
and condemned as heretical were not to be found at all in the said
book, and he denied the Pope the right to establish something as
a fact that everyone could and must convince himself of.
% --- p. 89 ---
\setcounter{footnote}{0}%
\opage{89}Already earlier he had (in a treatise on the
vacuum) decisively rejected the intrusion of faith in authority
into the problems of empirical science. And in the 18th Provincial Letter
he states that if Galileo's doctrine of the motion of the earth
is false, it is not the Pope's decree that establishes
this, and that if proof can really be given that the earth
moves, then it goes round both with him who
denies the motion and with him who assumes it. Pascal
was not a Copernican, but he touched here on a point very dangerous for a
Catholic at that time, the very one that had
brought Bruno to the stake and forced Galileo to perjury. A spirit
rebellious against ecclesiastical authority forms a subterranean
connection between Pascal as natural scientist and as
religious thinker. He even states that the Pope hates and fears
the scholars who have not sworn an oath to him (as monks
or priests) (fr. 873). This utterance comes about on the occasion
of the Pope's condemning the doctrine of the omnipotence of divine
grace and of the complete dependence of the human will
on God's will, this Augustinian
doctrine that Jansenius had renewed in sharp opposition to
scholastic and Jesuit theology. So the Pope himself is thus
not orthodox, and he can no more in matters of faith than concerning
facts be an infallible authority. But the Pope
has so much else to attend to that he has
no time to acquaint himself with matters of faith (fr. 882)! And when he
condemns the truth, one must cry out all the louder,
and appeal must be made from Rome to Heaven! We
must obey God more than men (fr. 920).

This last utterance Pascal's Jansenist friends found
it best to suppress when they published his ``Pensées''.
It would have been a deep sorrow to Pascal to notice that they
here did not give honor to the truth. Only death had prevented
him from publishing his last thoughts himself. --- On his deathbed
he received the sacrament from the hand of his parish priest.
His opponents saw in this a recantation; but it
meant only that he did not consider his relation to the true
Church broken. ---

% --- p. 90 ---
\setcounter{footnote}{0}%
\opage{90}Kierkegaard had developed within the Protestant
Church, from the first with great reverence for its leaders,
the living as the earlier ones. But in the years 1847—
50 his eye sharpened for the opposition between original
and modern Christianity, and in two of his most profound
writings (``Sygdommen til Døden'' 1849 and ``Indøvelse i Kristendom''
1850) a great reckoning with the Church is prepared. The main idea
in the said writings was that, especially in Protestantism, one
had forgotten the pattern of the Reconciler. One
took the reconciliation as a given fact and settled
cozily in the world as it is. Kierkegaard now appeals
to the representatives of the Church to make the concession that
there was a great distance between modern Christianity and
the Christianity of the New Testament, and that there had occurred a falling away
in relation to the heroic age of Christianity. But when
on the part of the official Church one was silent about this and even on given
occasion expressly proclaimed the claim to
be in continuity with ``the days of the Apostles'', he issued a
violent protest against this, especially in his agitational pamphlet ``Øjeblikket''
(1855). In the midst of the strife he collapsed and
died. He died without the sacrament of the Lord's Supper, as he did not
wish to take it from a pastor's hand. He is here a step
further out than Pascal; but he has not rejected the sacrament
altogether, and here he still, for his part,
hangs together with the old Church.

``Sygdommen til Døden'' consists in the despair that
inevitably is called forth by the high task being set
for man in Christianity, to live a life of the spirit in the
sharpest opposition to the life that natural feelings and drives
would lead to. Where this task has not come to
consciousness, perhaps a life of childlike joy in
life is led. But anxiety will again and again announce itself and put an
end to security. Man can then, in weakness, in reserve
or in defiance, try to subdue the anxiety and
withdraw from listening to the great demand; but despair
will always be present, hidden. Only by despair
breaking out openly and unfolding is there
% --- p. 91 ---
\setcounter{footnote}{0}%
\opage{91}possibility of getting beyond it. Only then is the condition
present for grace to come into force. But within Christendom
one has, now in thoughtlessness, now in romantic
enthusiasm, now in hypocrisy, let up on the great opposition
between the spiritual and the animal, and sought to level
it instead of working to overcome it.

``Indøvelse i Kristendom'' impresses that as long as Jesus
has not come in his glory, the Christian has to take
the humiliated Jesus as his pattern. The point is to look
away from what one believes oneself to have learned through dogmatics about
Jesus' being God. One must make oneself contemporary
with Jesus of Nazareth, the man who said he
was the Son of God and therefore was mocked and executed.
Instead, like schoolboys, one has sneaked to a result
and cheated oneself out of doing the sum on one's own. ---
The parallel with the Gethsemane mood demanded by Pascal
lies close at hand.

It became more and more Kierkegaard's conviction that
men had committed a great betrayal of Christianity.
One took the mildness and grace and forgot the
necessary presuppositions. And he maintains that, especially in Protestantism,
this betrayal lies plainly in the open. Catholicism
would never be able to sink so deep into spiritlessness
as Protestantism. And yet Catholicism has already
compromised. The meaning of Christianity collapsed when it was regarded
as something extraordinary to follow Jesus in
earnest by withdrawing from the life of the world. It is a senselessness
that there should be two kinds of Christians. But what
had thus begun was continued when Luther broke out of the monastery
and got busy making a new Church, which in a comfortable
way leveled the opposition between Christianity
and the world. ``Luther, Luther, Luther, you have a great responsibility!''
And yet Luther had gone through the school of despair.
But later generations then trust in the achievements of their forefathers
and have no need to perform achievements themselves. Protestantism
is a plebeian tendency. It levels the difference
between the high and the low in the world of the spirit. And
% --- p. 92 ---
\setcounter{footnote}{0}%
\opage{92}family life is largely to blame for this. Through
infant baptism all in the so-called Christian countries were made
Christians. ``And that'' --- exclaims Kierkegaard (Øjeblikket
VII) --- ``one dares to offer to God under the name of Christian
baptism! Baptism --- it was by this holy act that
the Savior of the world was consecrated to the work of his life, and after him
the disciples, men who had long since come to years of discretion,
and who now, dead to this life, promised to live
as sacrificed ones in this world of lies and wickedness.''

Kierkegaard's polemic culminated in the setting up of one
single thesis (in contrast to Luther's 95): ``the Christianity of the New Testament
does not exist.'' And to this was added a warning
against ``those who walk in long robes'' and against participation
in the public worship of God, in which God was made
a fool of!

The appearance of Pascal and Kierkegaard has great interest for the
history of religion. From a pronounced sociological standpoint
it has been maintained that religion is essentially a social phenomenon.
And in fact the history of religion is interested chiefly
in such religious phenomena as rest on
traditional cult. How, then, can individual personalities
whose striving and thinking lead them to a break with
their religious community (without their entering another
community) acquire significance for the history of religion? One could answer the
question by assigning each of them to his ``sect'',
Pascal to Jansenism, Kierkegaard
perhaps (as has been done in
an excellent French sociological work) to Pietism. This
would, however, founder on the fact that Pascal precisely in his last, decisive
turn separated himself with great pain from his earlier
fellow believers, and that Kierkegaard has only satire and scorn
for the Pietists and expressly separates his cause from theirs.
Rather one could say that they break with the ecclesiastical tradition
prevailing in their age only in order to maintain an
older tradition, in their eyes alone legitimate. But with
their appeal to what in their eyes was original they stood alone in
% --- p. 93 ---
\setcounter{footnote}{0}%
\opage{93}their time, and it was in the name of individual conscience
that they protested. There was no social authority to which
appeal could be made. Pascal found confirmation of his doctrine
in a miracle, but this miracle was not ecclesiastically sanctioned.
And Kierkegaard expressly declared that the religious
type that was for him the only valid one no longer
existed. Only in their relation to the sacrament on the deathbed
is there a bond left, which nevertheless did not prove strong enough
to bring about the connection with the Church
again. --- But it is an untenable theory that lays so great
weight on tradition and cult that one forgets what
goes on in the single souls who pay homage to these social
powers. It must become an essential task for the study of religion
to trace the development and attitude of religious personalities.
And then not least such individual
processes of development as lead out to the boundary of the religion in question
precisely through the attempt to live through
this religion in spirit and in truth. And such attempts
are precisely present in the religious development of Pascal and of Kierkegaard.

\sectionhead{4. \textit{Possible Ways.}}

What the two men would have done had they happened
to continue their lives, of that nothing can be known. But it
may be of interest to examine what possibilities
can present themselves to such human beings as, with
full conviction and assent and supported by their
own experiences of life, have followed the two spiritual combatants
right up to the point where they had thought their last thought.

There are three possibilities that present themselves here.

One could try to lead one's life by resuming the presuppositions of
primitive Christianity, in the great expectation
of the near breakthrough and with the giving up of
all participation in the work of culture and in the interests of ordinary
human life. There are perhaps --- who knows? --- those
who go this way. It was the way that in the Middle Ages
Francis of Assisi in his own fashion followed.

% --- p. 94 ---
\setcounter{footnote}{0}%
\opage{94}Or one could make the concession demanded by Kierkegaard
about the distance between original and modern
Christianity. But one would then have to go a step further and
acknowledge that modern Christianity is really another
religion than primitive Christianity. One then avoids
the great contradiction of thinking oneself to live by ideals
that one in fact denies. One must then conceive one's religion
so that, however essential a factor it is for life,
the conduct of life and the conception of life are partly guided by ideals
that have an entirely different origin and lead in an entirely different
direction from the religion one pays homage to. It may be difficult
to make these different guiding stars belong to
the same world-system. But the honesty that the two great
religious thinkers demanded is not violated in advance.

If one cannot go either of these ways, there remains the further
possibility of trying to resume the Greek
striving to ground an ethical conduct of life on human
nature and human experience of life, by investigating
what ideals and what ends must
consistently present themselves on this foundation. This is in fact the way
the endeavors for a philosophical ethics in the last centuries
have gone. I once (in my Etik\textsuperscript{5} X, 3) expressed this
thus: ``The Platonic concept of justice as personal
harmony we can adopt to express the inner
connection, the higher unity of self-assertion and devotion.
The ancient concept of harmony must be widened and deepened
by means of the ethical experiences that lie expressed in
the evaluations of character of Stoicism, Christianity, the Renaissance and modern
humanity. But such a change in the content of the concept
is quite possible without bursting its frame. Our ethics is Greek ethics, as
the later ethical development through its manifold
oscillations has been able to alter and correct it, but which will never
be able to be displaced as long as there is to exist anything that
can rightly be called human ethics.'' --- On this
foundation one cannot work without a belief developing at the same time
that such work is always possible, ---
that existence thus has room for it, indeed perhaps first
% --- p. 95 ---
\setcounter{footnote}{0}%
\opage{95}comes to full development through it. The closer shaping
of such a belief can come about only through poetic
symbolism. A great and deep poetry therefore itself belongs to the
values whose persistence and unfolding must be hoped for.

In relation to the two other indicated standpoints the position is
that one --- instead of renouncing all natural development
of human life, all independent culture, all life in
family and society, in art and in science, and instead
of undertaking an external mixture of Eastern religion and
Western culture --- takes one's stand resolutely on the foundation of human nature
and of Western culture and confidently
works on, as far as one can reach. There may be problems
enough here, and there may come great disappointments, so that
many lose heart. There may be strife about the ways in matters of
detail, and about the grounding of the thoughts that guide us;
but the main lines will constantly be discernible anew. ---
And nourishment may very well be drawn and instruction gathered
from other directions of life. The pathos with which the believer
awaits the ever-near Kingdom of God can find a more or less
pronounced counterpart in every striving after high ends.
Without tension no great work is carried through. If we compare
the significance of Christianity for European human life
with the significance of Buddhism for Asian human life,
the proposition can be defended that I have put forward in
my Religionsfilosofi, that Buddha has made Asia mild, but
Jesus has taught Europe a great Excelsior. But above all
the spirit of love of humanity, which Christianity
has proclaimed and (not always with full consistency) placed at
the head of the virtues, will also find its place in a purely human
ethics. Aristotle already stated that the highest form of justice
has the stamp of love, and conversely
it holds that love is a real virtue only when it
unites with wisdom and thereby becomes justice.
On the humane standpoint one stands in general in a
similar relation to Christianity as to other spiritual tendencies;
one appropriates from them what accords with
constantly continued striving. The great world religions have, like culture,
their origin in human
% --- p. 96 ---
\setcounter{footnote}{0}%
\opage{96}nature and human circumstances of life, and whatever of the noble,
great and beautiful has developed within the religions belongs
therefore in the end to the human race as a whole, not to a
single sect, even if this has ever so wide a diffusion.
---

Whichever of the three ways one then goes, there is one thing
one can learn from the two mighty figures we have here dwelt
on. They have both shifted the center of gravity in the life of the spirit
from the external authorities and from the content of doctrine to personality,
its experiences and its demands. Pascal motivates the claim
that Christianity is the true religion by saying
that only in it is the deepest need of human nature
satisfied, in that it at once explains and reconciles the greatest
oppositions of life. And for Kierkegaard the thought of God was
an anchor of distress under the storms of life, the necessity of which he does not understand
who always sails in smooth waters. It
therefore depends, according to the thought of both, on the right, deep-going
experiences being made in the course of life. But
first and foremost the point is not to deceive oneself
and others by thinking oneself to pay homage to a faith that has not come out
of one's own experience of life and one's own need. ``Pour les religions''
--- says Pascal (fr. 590) --- ``il faut être sincère; vrais
paiens, vrais juifs, vrais chrétiens.'' And in his last struggle
against the existing Church Kierkegaard says: ``I dare not
call myself a Christian; but honesty I will, and for it
I will venture.''

But there is a point at which those who go the third way
must be decidedly at odds with Pascal and Kierkegaard. They are
both clear that life, also the religious life, cannot
do without thoughts that hold fast to and guide the direction of life.
But they distinguish so absolutely between thoughts of life and scientific
thoughts, or, as Kierkegaard puts it,
between ``essential'' and ``inessential'' knowledge.
To the ``inessential knowledge'' all science is reckoned. And yet there
are scientific thoughts that can become of the greatest importance
for personal faith. The boundary cannot be drawn once for
all. Quietly and imperceptibly for most, the magnificent
% --- p. 97 ---
\setcounter{footnote}{0}%
\opage{97}development of science, both natural science and the sciences of spirit, in the last centuries
thus exerts its influence on religious conceptions. It is in general
a false conception that personality and science
should stand in opposition to one another as resting on
entirely different interests. Science is due to personal work
and springs from personal need, and the joy of knowledge
belongs to the noblest feelings of man. Even the forms
in which scientific thought moves, without always
being conscious of it, lead back on closer investigation
to personality, in that they are conditioned by the constitution of the
human spirit. On the other hand personality is developed
through scientific work, gains
in clarity about life and its conditions and learns to set itself tasks
that lead far beyond egoistic interests. And
as it goes with science, so it goes with all culture.
Man is both greater and smaller than the culture given at any
time: greater, because all culture must in the end justify
its value by contributing to the development of personal life,
--- smaller, because culture sets tasks that require the subordination
of individual interests to that which serves
the lasting progress of humanity as a whole.

They are two solitary figures we have dwelt on. Successors
they have not had, and a school they have not formed. But
the analogies between the thoughts that guide each of them
and between their appearance in face of the existing order that each
for himself stood before contains instruction for all times.
From their solitary stages they caught sight of great, decisive
problems in the domain of the history and the psychology of religion, and they have staked their lives on
fixing these problems so firmly that they cannot be forgotten.
And oblivion will so much the less become the lot of the two men
as they have set forth their problems not merely with
passion, but with a wealth of thought and a fullness of mood
that in and of itself lets their works belong to the
most precious treasures of humanity.

% --- p. 98 ---
\setcounter{footnote}{0}%

\opage{98}\essayhead{SIBBERN AND KIERKEGAARD}

Both of these thinkers have had a great influence on my
development, and even after I believed that I had come to
years of discretion in philosophical matters, I have often returned
to them. And when this past winter [1925] some young
students invited me to give a few lectures
on Danish philosophers, it was above all these two men on
whom I dwelt. It is the accounts I gave on that occasion
that I am now working up into a short essay.

Kierkegaard was born in the same year in which Sibbern became Professor
of Philosophy at the University of Copenhagen (1813). The difference in age
between them was thus great, and the difference in character
no less. And yet they stood in a personal relation
to one another. Kierkegaard for a time attached himself to Poul Møller,
whom he loved and admired, but after his early death
he stood for a time quite close to Sibbern. This is known not merely
from a letter he wrote from Berlin to Sibbern (1841). But one
of Sibbern's daughters has told me in a letter that it belongs
to her childhood memories that she walked with her hand in
her father's, while Kierkegaard walked on the other side, eagerly
talking and now and then stopping in order to get said what
lay on his heart. She also remembers that Kierkegaard
often visited her father and liked to sit in the sitting room
in the twilight, when the fire shone from the stove. Of the content of
their conversations nothing is reported, but the difference between the two
men's characters and points of view must again and again have
made itself felt. Sibbern knew from his own experience times
of inner struggle and ferment; the description of such times
% --- p. 99 ---
\setcounter{footnote}{0}%
\opage{99}is among the best things in his Psykologi. But in spite of everything he kept
an open, naive view of life. As a psychologist
he had understanding and sympathy for all mental phenomena,
and it was a favorite saying of his that just as
for the botanist there are no weeds, so there was
nothing mental that the psychologist could not have an interest
in studying. Not only by reason of this intellectual interest,
but also by his firm faith in life, he constantly
caught sight of a harmony behind all oppositions. Kierkegaard's
melancholy and never-resting reflection made him walk
along abysses, from which only a convulsive self-assertion
saved him. His sharp eye for the oppositions of life
led him to the thought of an absolute paradox, and the capacity
to suffer became for him the highest form of the will. But
despite the great opposition between Sibbern's open eye
for everything human and Kierkegaard's passionate confinement
to the narrow way, which for him led to the one thing
needful, they nevertheless met as proclaimers of the principle of personality,
of the thought that the truth that has
value must spring from man's own experience and reflection.
This principle was, remarkably enough, pronounced in one and the same
year (1846) by the older and the younger thinker. They
hardly knew that the principle of personality had been proclaimed with great clarity
some years earlier (1842) by the Swedish thinker
Erik Gustaf Geijer. All three Nordic thinkers broke by
this principle with Romanticism, with its life in the world
of fantasy as well as with its dependence on what is historically
handed down. It is of great importance to note that it
was not merely strict scientific and historical research
that led away from Romanticism and the conservatism connected with it, but also the individual personality's
claim to free, independent development.

Sibbern (1785---1872) plunged in his youth into the
romantic current, which was felt as a refreshment after
the dryness of rationalism. But he emerged again, and his
whole life was a striving after free research on life with
% --- p. 100 ---
\setcounter{footnote}{0}%
\opage{100}preservation of the sense for the greatness of the past that had distinguished
Romanticism.

Sibbern was first and foremost a descriptive psychologist. I
shall here dwell only on the topics that he describes best. To these
belong, as already mentioned, periods of ferment. He was
clear that the deeper life opens in a mind, the
sooner it may come to pass through a chaotic mixing and stirring
of the elements in its interior. And such ferment
easily leads to a morbid preoccupation with oneself, to brooding
and anxiety. He held that such irritability and nervousness
were especially characteristic of ``our time, in which partly the
greater struggle for physical and civil subsistence has worked
toward self-reflection, partly individuals no longer, as
formerly, feel themselves, either in a religious or in a civil respect,
borne by a firm order of things in life''. This
statement is from 1828; there are surely many who think that it
fits the present day. In contrast to ferment and self-reflection
Sibbern then sets sympathy and admiration, the joy
in life as it unfolds around us, and absorption
in the great and the exalted, which leads the individual
out beyond himself.

Psychology was for Sibbern a part of a general
theory of life, to which physiology too makes its contributions.
Psychology and physiology belonged for him closely together. One and the
same fundamental activity appears to us in two forms.
The word soul does not mean a being subsisting by itself, but
a living and working that always has a corresponding counterpart within
the organic processes. Sibbern was neither a spiritualist
nor a materialist. The conception of unity maintained by Spinoza and Leibniz,
to which he adheres, had met
him in his predecessor Treschow.

For Sibbern, according to this conception, mental life was subject to
the same great laws of nature as everything else, in particular
the law of development. Early it became for him a main point
that all development proceeds sporadically, from different starting points,
and that the whole in all domains comes
to be through the fusion of processes that for a time might
% --- p. 101 ---
\setcounter{footnote}{0}%
\opage{101}seem to have nothing to do with one another. He illustrates
it by examples from both the inorganic and the
organic world, and in the domain of mental life he finds
here an explanation of why ferment, doubt, and disharmony
appear so easily; they are due to the way in which
need and capacity so often first appear as different,
perhaps opposed, so that it becomes the task to harmonize
them. In the social domain the sporadic appears
in men's different needs and conceptions and in oppositions
between peoples.

Sibbern's world-view is conditioned by this conception.
A great development goes on from chaos to harmony.
The world, and especially human beings, has perhaps not come very
far along this way. As for human beings, they are
perhaps still only in the embryonic state, or at most in the age of rowdiness.
And Sibbern then looks hopefully toward the continued
course of development. He was, as his close friend A. S. Ørsted
said of him, an incorrigible optimist. ---

In his conception of the position and significance of philosophy
Sibbern stood in opposition to the Hegelian philosophy, which in
Denmark was essentially embraced by aestheticians such as Heiberg
and theologians such as Martensen. A system completed once and for all
he considered impossible. Even if it should succeed
in gathering the knowledge of the age into a whole, then, while
this systematic work was in progress, experience would have
advanced, and there would in fact have come a new age
with new content. Moreover Sibbern finds here too an application
for his doctrine of sporadic development. Every cognition
unfolds from a single point in existence. And
thus every philosophy too stands at a definite stage
and can only present how existence looks
from this single point. We then understand why there are disputes
between philosophers. It is a sign that our cognition
is in its becoming.

Also as regards philosophy's relation to other spiritual
domains Sibbern has a conception other than Hegel's.
Art and religion are not imperfect forms of philosophy,
% --- p. 102 ---
\setcounter{footnote}{0}%
\opage{102}but independent forms of spiritual life, whose significance
it is the philosopher's task to clarify.

In his youth Sibbern was strongly captivated by positive Christianity.
He praised it for the fullness of spirit to which it opened
access, and for the rest and refreshment it could
bring. Yet he maintained the independence of philosophy over against
Christianity as over against all other forms of
life of the spirit. Gradually he was led to distinguish sharply between
the developed dogmas and the experiences of life on which
they were originally built. And his idea of the sporadic
in all spiritual development had to lead him also to the conviction
that one and the same faith does not correspond to all stages
or to all individual experiences and forms of life.
His own experience of life led him away from Christianity.
Everything that has spiritual value begins from within, as
subjective. That is the way of the spirit. Later formulations
and organizations may take place and an apparent
objectivity be reached. But when one dogmatizes, one
easily removes truth from its living root and makes a
mummy of it. And the circumstance that great communities
unite around certain dogmas does not make these any
the more objective, if by ``objective'' one understands the
same as by ``valid''.

These thoughts Sibbern set forth in a university program
in 1846, one of the most important documents for the history of the religious problem
in Denmark. In his Utopia (written in the years
1858---1872), to which I shall presently return, he depicts
the human beings of the future, as he hopes they will
be minded three hundred years hence. They will no longer call
themselves Christians, because the history of the church in their eyes is full of
presumption and domineering. Yet they find much that is great and
beautiful in Christianity and therefore distinguish between Christliness
and Christendom. But they read in the great book of
nature and of life, not merely in the Bible. They live out of a power that
stirs in their interior, and they have the faith that their own
striving forward and the divine working go together as one.

The social question, which Sibbern in the said Utopia
% --- p. 103 ---
\setcounter{footnote}{0}%
\opage{103}(whose title is ``Meddelelser af et Skrift fra Aar 2135'') treats,
had already before the Revolution of 1848 aroused his
attention. The first preface to the Utopia is dated
1846. And in a university program of 1849 he urged
the statesmen to familiarize themselves in good time with the
idea of communism, which is the expression of a higher, distributive justice, in
contrast to production and competition, which contain
no guarantee that the proceeds really benefit
large circles. He extols the social question
in contrast to the constitutional (parliamentary) one,
which mostly occupied minds. In the future society he depicts
in his Utopia there is no question at all of state and church affairs,
any more than of legal, commercial, and property affairs.
Work is organized in such a way that all take part both in
spiritual and in material activity. The proceeds are common. From public
storehouses everyone takes what he needs.
The conditions are depicted in detail with many naive traits.
The naivete shows itself on the large scale in this, that no one asks
what security there is that the ideal organization
will be maintained, any more than how it has come
to be. The question, difficult for all ideal depictions of the future,
of how the transition from the present
economic and spiritual culture to the ideal state can
take place, the naive utopian avoids by letting human beings
(with the exception of a small chosen band) fall into a long
sleep, during which the old institutions collapse.
On waking one is then not minded to set
them up again. Nor is it asked how the drive
to work and the initiative toward new directions can be preserved
under the new conditions.

Sibbern's Utopia has its greatest interest in this, that whereas
Romanticism had been turned toward the past, in which it
found immediate revelations of the eternal, the mind
of the old Romantic now shows itself, with sharp criticism of the remnants
of the past, turned forward toward new possibilities. And
like other utopias (Fourier's, St. Simon's, Owen's) from
the same time Sibbern's bears witness to a warm feeling for those
% --- p. 104 ---
\setcounter{footnote}{0}%
\opage{104}who suffer under the imperfections of present social life.
---

As I have once said before, in no
Danish author have psychological sense and a feeling for humanity
been joined with such freedom of spirit as in
Frederik Christian Sibbern.

Through Poul Møller Kierkegaard stands in connection
with Sibbern. The latter brought him a greeting from Poul
Møller's deathbed, and Kierkegaard attached himself, as we have seen,
for a time to him. This personal relation has interest in
bearing witness to the continuity in Danish philosophy.

In his Papirer (ed. by Heiberg and Kuhr. Fifth Volume
p. 20) Kierkegaard says: ``Danish philosophy, if there
is ever to be talk of such a thing, will differ in this
from the German, that it does not begin at all with
Nothing, or without any presupposition, or explain everything by
mediating, since on the contrary it begins with the proposition that
there are many things between heaven and earth which no philosophy
has explained''. This proposition, he thinks, by
being taken up into philosophy, will supply the proper corrective
and at the same time cast a humorous-edifying emotion over
the whole. --- What Kierkegaard here expects of a Danish philosophy
that might come fits, in my view, well
the Danish philosophy as it has developed from Holberg's
time on. In the whole series of Danish philosophers before Kierkegaard
--- Holberg, Sneedorff, Steffens, Treschow, A. S. Ørsted,
Sibbern --- we find\footnote[1]{In my book \textit{Danske Filosofer} (1909) I have tried to show this more fully, and in \textit{Udvalgte Stykker af dansk filosofisk Literatur} (1910) one will find testimonies of it. } as the chief peculiarity
a sense for the individual and in close connection with it
an emphasis on experience. For one cannot understand individuality
except with the help of experience. There may have been
disagreement about the extent to which the problematic appears
in the world. Thus Holberg thought that with the widening
of experience we should face still more problems
% --- p. 105 ---
\setcounter{footnote}{0}%
\opage{105}than we have hitherto to work with, whereas the more optimistic
Sneedorff thought that many of our present
problems will fall away with widened experience. But on the whole
the characterization given fits. And Kirkegaard
himself marks a paradoxical high point of the prevailing
tendency, as regards the emphasis on the individual;
but this advanced position, as regards the one question,
is purchased in his case by a narrowing of the broad
basis of experience on which Danish thinking earlier built.

I shall here speak only of Kierkegaard as philosopher, and
the chief writings here are ``Filosofiske Smuler'' and ``Uvidenskabelig
Efterskrift''. This latter work, which appeared in 1846,
concluded the rich production in aesthetic, religious, and philosophical
direction which Kierkegaard carried out in the first half
of the forties, and which he later, as can be seen from his
papers of 1849, made an attempt to derive from a definite
preconceived plan, an attempt which his great love of truth
nevertheless led him to give up, as he became clear
that it was the involuntary urge and capacity to produce
that was the driving force. The satisfaction of the strong
urge to produce was for him at the same time a means of mastering
the deep melancholy that ruled in his innermost being.
``Uvidenskabelig Efterskrift'' occupies a central position in
Kierkegaard's production, inasmuch as it gives a retrospect
on his preceding works and supplies information
for a right understanding of them. In philosophical respect
it has the great significance of clarifying his position as
epistemologist and as philosopher of life, the two points of view
from which we shall here consider him. For a while he had
thought that this work should form the conclusion of
his production as a whole. But a deepening of his religious
life led in the following years to a new and mighty unfolding
of religious authorship, which again led to his
coming to feel himself in stronger and stronger opposition to
the usual conception of Christianity and to the established
church.

% --- p. 106 ---
\setcounter{footnote}{0}%
\opage{106}--- For Kierkegaard thinking was characterized by
constantly seeking and finding difficulties where plain
understanding, and often the speculative systematization too,
found the way smooth and passable. It is difficulties, riddles,
which in their highest degree become paradoxes, that keep
thought alive. And since life constantly presents new riddles
and tasks, whose solution demands spiritual work,
no system can be completed. Life is not finished, but
unfolds ever further, and thought, which explains life,
always comes afterward. We live forwards, but we understand
backwards.

Sibbern had already taken up a similar position in
opposition to speculative philosophy. But in Kierkegaard
the line of thought is developed with greater emphasis and consistency.
With enthusiastic assent he seized Lessing's words
about the eternal striving after truth as the highest
that human beings can attain, and he found in this saying an expression
of the self-knowledge every honest human being must come
to, and therefore an expression of an ethical view of life. But
also in purely theoretical respect Kierkegaard takes up the question.
He demands a revision of the problem of knowledge in
Kant's spirit and thereby stands as a forerunner of the great
work that was done toward the end of the century
within epistemology.

He asks whether the fundamental ideas of science, the so-called
categories, which constantly come to light in
special scientific work, can be found through abstract
speculation or whether they are due to experience [Erfaring]. It was a
question he wished to treat more closely. This wish
he did not see fulfilled, since his philosophical production
was interrupted by the turn in his spiritual life mentioned above.
But basically he has given the answer to his
question by maintaining that understanding always comes
after experience.

I may perhaps venture to insert here that I have tried to discuss
this question in my book ``Den menneskelige Tanke'',
in which I show that the categories are found by analysis of
% --- p. 107 ---
\setcounter{footnote}{0}%
\opage{107}scientific work, as it develops through
the ages. I lay special weight on the fact that what is
a matter of course to the inquirers of one age has not always stood so
for the inquirers of earlier ages. A presentation of scientific
fundamental ideas, a theory of categories, can therefore
mean only a stage in the development of human thinking.
Categories can arise, and they can die. --- My
Religionsfilosofi too builds on this conception of the relation
between experience and understanding.\footnote[1]{When Kierkegaard thought that for the philosopher ``a humorous-edifying emotion'' will make itself felt in the mind when he becomes conscious of the limitation of knowledge, then the experience I have had in my work of thought has, for my part, confirmed this, and I have sought, both in \textit{Den menneskelige Tanke} and later in \textit{Den store Humor} (see that work p. 117---123), to give expression to a mood of the said kind.}

Since the presuppositions on which all justification rests
cannot themselves be justified, one comes, according to Kierkegaard,
only by a leap into scientific thinking.
Here there makes itself felt what Kierkegaard calls a dialectic
of beginning. I must begin my reflection at
a single point, with a single assumption, which only later
on, through continued experience and reflection, can come
to stand as justified, and thus for the time being stands as arbitrary.

On this point Kristian Kroman has developed Kierkegaard's
thought further, in that in his book ``Vor Naturerkendelse''
he maintained with great sharpness that the first
principles of science are fixed by an arbitrary act in the
forms of the postulate. About my own position on this question I shall
later have occasion to speak. ---

When Kierkegaard is characterized not merely as epistemologist
but as philosopher of life, it means that he has
brought forward points of view that are of importance for the
way human beings take life, and that these points of view have
sprung from deep psychological insight and from his own dearly bought
experience. In our time there rages much, especially in Germany,
% --- p. 108 ---
\setcounter{footnote}{0}%
\opage{108}that calls itself philosophy of life, but which is only an attempt to make
speculative and dogmatic assumptions
valid in the twinkling of an eye, notwithstanding that they have lost their theoretical
foundation. They are attempts to solve the problems without
posing them in earnest. But precisely the earnestness and energy with
which Kierkegaard has depicted life's relation to thought,
and especially to those thoughts that life itself cannot do without,
has now also abroad more and more drawn attention
to him and to his points of view, which I shall now
try to characterize.

As stated above, Kierkegaard forms the conclusion of
the characteristic tendency in older Danish philosophy to
maintain the individual and lay weight on experiences, but
a conclusion that is characterized by the fact that individuality
(or, as Kierkegaard expresses it, subjectivity)
is stressed to the utmost, while the element of experience at the same time
is limited in its significance.

As, according to Kierkegaard, it was by a leap that we came
into the world of thought, so it is also by a leap that we come
into the view of life in which we land. It is personal
need, the feeling of spiritual distress, finally despair, that
leads to seizing that sheet anchor on life's stormy sea
which every view of life is. Here the Lessing
proposition about the eternal striving is sharpened to the proposition that personality
(subjectivity) \textit{is} the truth. That is to say,
that the truth that a human being can find on life's way has
no other ground than precisely personality's own need.
No science can help us here. Every view of life
needs thoughts, and so far knowledge. But Kierkegaard
here makes a sharp distinction between essential and
inessential knowledge, a distinction determined by
the opposition between the thoughts that are necessary to us
in order to keep ourselves up in the struggle of life, and such
thoughts as only intellectual interest leads one to
immerse oneself in. To inessential knowledge he counts all objective
science, theoretical philosophy, history and philology,
as well as natural science. He sensed that
% --- p. 109 ---
\setcounter{footnote}{0}%
\opage{109}natural science, by its closely connected points of view and methods,
could become a similar danger to his outlook
as the speculative philosophy and theology he had long
combated. From the view that life is something individual,
not accessible to abstract theories, he detached
it from any connection with science.

The question then presents itself whether every view of life
is as good as every other, or whether there is
any standard by which an evaluation can here be made.
What is high and what is low, --- what is great and what is
small in this domain?

Already earlier the question of such a standard had come up
in Danish spiritual life. It was in the dispute between
Baggesen and Grundtvig, in the last analysis a dispute between
philosophy and theology. Grundtvig agreed with his opponent
in fighting the battle on purely human ground. ``Here we
do not ask much about confessions!'' he exclaims. But
when the dispute between the two spiritually moved men petered
out, it was because they could find no common human
standard.\footnote[1]{On this dispute see \textit{Danske Filosofer,} p. 66---72.} Kierkegaard has tried to maintain
such a standard, and his significance as philosopher of life
rests not least upon this.

Kierkegaard's standard is taken from the very concept of personality.
In determining this concept one will always
come back to the fact that it presupposes energy to comprehend
a greater or smaller manifold of impulses
and experiences, so that a unified bearing toward
the surrounding world can become possible. The fuller and more
comprehensive the content that is to be comprehended, the
greater the energy presupposed, if the inward concentration
is not to fall away. It is basically the concept of synthesis,
deriving from Leibniz and Kant, that lies here at the
foundation, and philosophy of life will here always prove to stand in debt
to theoretical philosophy.

By this standard Kierkegaard then divides the different
views of life or, as he with an unfortunate word
% --- p. 110 ---
\setcounter{footnote}{0}%
\opage{110}calls them, the different ``stages''. --- Concentration
may be limited to the single moment and what can fill
it (the ``aesthetic'' stage). --- One may be
apparently absorbed in the single situation and yet
inwardly stand free toward it. It is the jest hidden in earnest
that characterizes irony. --- One may in inward
earnestness immerse oneself in one's duty (the ethical stage). --- One
may regard life's content and tasks as unable
to fulfill the mind's whole striving (humor). --- One may
finally see life's real task as lying beyond all experience,
time as vanishing in relation to eternity. That
is the religious stage, which in Christianity appears
in paradoxical form, because there the eternal itself appears in
time and in a human shape one may meet on
one's way.

As, according to Kierkegaard, one comes by a leap over
from life in the world of thought to the faith or view of life
by which one leads one's life, so it is also by
a leap that one comes over from one view of life to
the other. It is continuity under all its forms that Kierkegaard
is out to destroy --- remarkably enough precisely on the basis of the
inner, closed coherence that for him rightly characterizes
personality!

Both as epistemologist and as philosopher of life
Kierkegaard's great significance is conditioned by the paradoxical
sharpness in the posing of problems. It made him a strict
teacher for the generation that developed under his
influence. And traces of this will be noticeable even where
development has led far away from his standpoint.

From my present standpoint I will remark the following.

Every leap must stand for thinking as a task,
not as a last word. Why does one leap precisely
at this particular place and in this particular direction? And
does not the leap itself signify a continuity, even if it signifies
a fall! --- Philosophy, which has the task of investigating
the relation between the different sciences, must
also discuss the leaps, the abrupt transitions between them.
% --- p. 111 ---
\setcounter{footnote}{0}%
\opage{111}It finds these transitions conditioned by the adoption of new
presuppositions, but then the necessity of these
new presuppositions becomes something to investigate. And the same holds
of the transition from thought to life, a transition which moreover
always presupposes a preceding transition from
life to thought. Here, then, there is a constant reciprocal action.

In the transition from thought to life one cannot, as
Kierkegaard wished, distinguish once and for all between essential
and inessential knowledge. Kierkegaard's standpoint
here signifies a narrowing in relation to the broad
foundation on which earlier Danish thinkers, especially men
such as Holberg and Sibbern, built their views of life.
The widening of the world-picture was for Holberg, the great development
in nature from chaos to cosmos was for Sibbern, very
essential elements both in their faith and in their knowledge.
What significance has science, then, if
it does not, perhaps by many detours, in the end gain influence
on human views of life and of the world? And
such influence is indeed constantly exerted, and in all quiet,
even if only an eye trained in the school of the history of philosophy
discovers it. ---

I have now for the last time spoken of the two Danish thinkers
who have meant most to me. And it is a joy to me that
the invitation to occupy myself with them once more came to
me from a circle of young students.

% --- p. 112 ---
\setcounter{footnote}{0}%
\opage{112}\essayhead{PHILOSOPHY \\ AND ECCLESIASTICAL RATIONALISM}

Dr. Michael Neiiendam's interesting study of
Christian Bartholin has clearly illuminated an ecclesiastical tendency
which was for the most part derided by posterity or at any rate little
understood, but which is now subjected to a purely historical
illumination. The main figure's character and thought, which are
so characteristic of his time, the latter half of the eighteenth
century, emerge clearly and distinctly.

For philosophy ecclesiastical rationalism has in itself
very little interest. It rests on a popular philosophy
without clear principles, and it likes to build on crumbs that
had fallen from the table of the properly philosophical inquirers. On the other hand
it is of interest to see, as the study also
shows, that stricter philosophical thinking contributed to the dissolution
of the untenable attempt to solve the religious
problem that underlies ecclesiastical rationalism.
But in the gallery of personalities which the study
brings forward, and which, as far as I can see, makes good contributions
to the understanding of the spiritual conditions of the eighteenth century
in Denmark, there appear two personalities whom I do not find
thoroughly and justly treated. From a philosophical
point of view there are essential objections to be made to the author's
presentation.

1. I do not see that it would have been necessary for the author
to mention \textit{Lessing.} There has, as he himself says,
at that time been no influence of Lessing on
Danish theology, apart from the fact that his polemical writings
are occasionally used. Such an influence was exerted, on the contrary,
later through Kierkegaard's understanding and
% --- p. 113 ---
\setcounter{footnote}{0}%
\opage{113}enthusiastic mention of Lessing, and if the author had paid
attention to it, he would have expressed himself differently about
Lessing than he does. He lets him cooperate with the Enlightenment
and strengthen its critical sense. This is in
any case not the whole truth. Lessing belongs, as from
other sides do Rousseau, Hume, and Kant, to the personalities
who --- even while they are in many respects
permeated by the spirit of their century --- point beyond
their contemporaries to a richer and deeper conception of the problems.
What is most characteristic of Lessing
is his sense for genuine and original content of spirit,
wherever it is found. This sense led him in his dramatic
criticism (``Hamburgische Dramaturgie'') to bring forward
Shakespeare's mighty poetry in contrast to French-classical
tragedy; --- it led him (in ``Laokoon'') to draw
a boundary between the visual arts and poetry; --- and it led
him (against rationalism as against orthodoxy) to
point to the inward and immediate movements of religious life
in contrast to the reflection and the learning
which religious documents can give occasion for. But
Lessing also insisted that if one wished at all to
use thoughts in one's view of life, one ought to think
them to the end. There were two different currents that
stirred in him from his youth on, the one nourished
by his historical sense, the other by his philosophical
thinking. He did not succeed in bringing these two currents
to unite. But who, basically, has succeeded?
--- By making them both count against the religious
problem he has maintained the right of personality and of
independent thinking against the unclear attempts at mediation
of which not only ecclesiastical rationalism but
orthodoxy itself constantly makes itself guilty, and he
thereby acquires a significance that reaches far beyond his time.

The attempts at mediation that ecclesiastical rationalism undertook
were more hateful to Lessing than the old orthodoxy, because
they (as he wrote to his brother) ``under the pretext of
wishing to make us rational Christians make us irrational
% --- p. 114 ---
\setcounter{footnote}{0}%
\opage{114}philosophers''. But the (then) modern representatives of orthodoxy too
(such as Schumann and Goetze) had to count for him
among the unclear mediators. The free religiosity he hoped
for would --- he thought --- rather unfold when
one let the old faith speak openly and honestly. He
did not want the original religious content bound up
with speculative ideas, as had happened early
in the church. Positive religion is something concrete and historical,
and on that no eternal truths can be built. Positive
religion is built throughout on what is assumed to
be historical truth, and can neither lead to nor be derived
from general ideas. The proper proof of religion is
inward appropriation of a historical content. --- If Lessing
had been more of a psychologist than he was, he could have
spoken more clearly. Yet he expresses himself clearly enough against
intellectualism in the religious domain.

It was indeed Lessing's faith that religion would
have to pass through a development by which it increasingly
came out beyond the purely historical and therefore more
or less arbitrary foundation. But this development
had to take place slowly. He advised against extinguishing one's lights
before the sun had risen. And this caution had precisely
its root in his sense for personal appropriation, for what one
has later called personal truth. His historical sense
moreover taught him that neither the so-called truths of revelation
nor the so-called truths of reason were the
same at all times.

Against the orthodox Hauptpastor Götze Lessing maintained
that Christianity does not derive from written tradition,
but from the events these accounts are
about, and it was also in dealing with him that Lessing had occasion
to utter the famous proposition that the constant striving
after truth is to be preferred for a human being to the
constant possession of truth. Later, over Jacobi's
sentimental semi-theology, he had occasion to defend the
great and quiet movement of thought that lies behind Spinoza's
crystallized systematics. In both places his
% --- p. 115 ---
\setcounter{footnote}{0}%
\opage{115}genius in conjunction with his congenial insight has brought forward significant
points of view. It is characteristic of Lessing
that it is two clashes that have called forth his most significant
utterances. He was the great seeker, belonged
to no confession or to any system. And it is
in any case not right, as has been done in Dr. Neiiendam's
otherwise so interesting book, to present him merely
as a representative of the Enlightenment.\footnote[1]{See further on Lessing \textit{Den nyere Filosofis Historie}\textsuperscript{3}, III p. 16---25.}

2. The other figure, who is distorted in the book just
mentioned, is the Danish author \textit{Johannes Boye} (1756
--- 1830), whose writing (``Statens Ven'') is there called ``the most important
source for the understanding of the domestic Enlightenment''. But
Boye is not given his due by being
characterized as a man of the Enlightenment. He is a Renaissance nature, maintains antique
clarity and strength and shows how ethical ideas develop
under the struggle for existence as the only
possibility of putting an end to the sole rule of physical power.
Although he combats Kant's grounding of morality
in pure reason, his manly conception of life, according to
which spiritual development goes through strife and labor,
is quite in Kant's spirit. ``Nature,'' he says, ``brings up its
children strictly, but it does so in order to produce heroes.'' --- Boye
lays great weight on the historical-psychological presuppositions
for the coming-to-be of a purely human morality. But
Kant himself has done that at a significant stage of
the development his ethical ideas underwent. Man's
disposition as a rational being can according to Kant come
to development only in the course of the generations, which through
the strife of the passions leads to the idea of a society
within which the freedom of the individual through being bound by law
becomes compatible with the equal freedom of others. Without
strife the human faculties would not be developed. The trees
in the forest, says Kant, seek to rob one another of light and
air, but precisely thereby are compelled to seek light and air in
higher regions. (Cf. further on this point Den
% --- p. 116 ---
\setcounter{footnote}{0}%
\opage{116}\textit{nyere Filosofis Historie}\textsuperscript{3}. III, p. 74---78). Neither Kant himself
nor his outstanding Danish interpreter Anders Sandøe
Ørsted was clear about the relation between psychological-historical
investigation and ethical theory of principles --- and this relation
is a problem to this day. Ørsted proceeded from pure
rational morality as fully finished, Boye from the constant struggle
for existence. And it was by no means a defeat that Boye
suffered in bringing forward this side of the problem.\footnote[1]{See further on the remarkable thinker Boye was, \textit{Danske Filosofer,} p. 33---35. --- At about the same time as my attempt at a rehabilitation of Boye (in the book mentioned) \textit{Vilhelm Andersen} gave a very interesting characterization of him in ``Typer og Tider''. }

% --- p. 117 ---
\setcounter{footnote}{0}%
\opage{117}\essayhead{A MODERN BUDDHIST}

\essaysource{Arthur Pfungst: Gesammelte Werke. I. 1926.}

Through my late colleague, the famous Pali scholar
Viggo Fausbøll, I came to know the remarkable man
whose collected works, which are to comprise three volumes, are now beginning
to appear. Pfungst's acquaintance with Fausbøll
sprang from his great interest in Buddhism, which
for him not only had historical significance, but in essential
respects determined his view of life and gave him
expression for what he had come to by his own way.
Fausbøll, as is well known, edited the Buddhist writing Dhammapada
(The Words of the True Doctrine) from the manuscripts
that Rask had in his time brought here from Ceylon, and later the Jataka Book
from the same sources. There was a great contrast
between the two men, whose friendship was built on the common
interest in the ancient Indian books of wisdom. Fausbøll
was a quiet scholar, absorbed in his scientific work,
which gave him a fine name in the circle of Indologists and even
made him known and esteemed among the Indian Buddhists,
although he did not subscribe to their doctrine. He admired Buddhism,
not least because under its
rule there had never been committed such abominations and cruelties
as under Christianity. But he was not a Buddhist,
since he could not give up faith in the immortality of the soul
and in a personal God. This view of his he expressed
openly in a short autobiography he wrote in the year 1900
in the Buddhist paper ``Journal of Maha-Bodhi Society''.
Pfungst admired the candor of this statement,
since the paper mentioned was a propaganda paper for Buddhism,
and since the doctrines Fausbøll held to were
% --- p. 118 ---
\setcounter{footnote}{0}%
\opage{118}heresies from a Buddhist standpoint. That the paper without
more ado reprinted Fausbøll's statements was indeed precisely a sign
of the great tolerance that Fausbøll praised in Buddhism.
Arthur Pfungst was a man of the world, rich, heir to
a large factory in Frankfurt am Main, strongly absorbed in all
movements within the different domains of culture, politically
and religiously radical, poet and businessman at once.
His business trips took him around the world, and he
often came here to the North, where he never forgot to visit
Fausbøll. He was lively and cheerful, and only from statements
that came out occasionally did one learn that behind his
cheerfulness and his sparkling wit lay the conviction that
life was not worth living. Buddha's teaching was Pfungst's
own religion, and he worked for it not only as
poet, but also in a series of essays which he
collected in 1904 in the book ``Aus der indischen Kulturwelt''. One
might ask how such a duality, such an
opposition between the inner and the outer, was possible. But
from what I have known of him, I believe that one of his
closest friends, Dr. Angermann, who has written the
``biographical introduction'' to the first volume, is right when
he states:

``For the pessimist the unpleasant side of things
stands fast as the most self-evident thing in the world, so that from that
side no surprises come for him, whereas from
the other side good surprises may come and the eye
may be sharpened for what is joyful and playful that life may
bring. The rejection of this world stands fast for the pessimist,
but the more he knows that within its bounds
he must not make life more difficult for himself and others
than it already is. In other words, the pessimist's humor
is the deepest and finest; it is in its innermost essence
the sage's pure and good disposition toward others and himself.
Therefore in Pfungst all that was sharp and wounding was also lacking,
and this was just what All who knew him praised as the
most beautiful quality of his humor.'' ---

It is surely not necessary to remark that the word humor
% --- p. 119 ---
\setcounter{footnote}{0}%
\opage{119}in this statement is used in another sense than the popular one,
that of the capacity to make people laugh. Nor will it
be necessary to recall that there are several different
forms or kinds of humor besides the peculiar
kind that Pfungst's friends found in him. And when we shortly
shall hear more closely about Pfungst's peculiar relation
to Buddhism, it will appear that his humor in reality
had other elements at its foundation than pure
pessimism, elements that bring it into close kinship
with the humor described by Thackeray and George Eliot,
which is due to a union of comic sense and sympathy
(wit and love).

It was hardly the study of Buddhism in Pfungst's case
that made him a pessimist. Nor has the reading of
Schopenhauer been decisive in this respect. Nor
has it been his own disappointments and sorrows that led
him to his total feeling toward life. He has, as far as is
known, not been able to say with Edgar in King Lear that he
has come under the blows of fate and through his own sorrows has
learned to have fellow-feeling with others. And there has hardly stirred in
him, as in Schopenhauer, early an unruly sensual
craving and an equally unruly hot temper, which could let
his imagination apprehend existence as a tormenting fire. If
one should conclude from his great epic poem ``Lakaris'', which I shall later
discuss, that a strong lust for life in its author
was subdued through disappointments and sorrow, then
those nearest him have in any case known nothing of it.
His biographer thinks that his view of life was
founded in the so often fateful years of transition, in
which there is often great receptivity to the tragic
sides of life. The world-sorrow that easily arises thereby, and which
for most is only a passing phase, acquired lasting
significance in Pfungst, and it found early expression in his poems.
He had grown up in a rich and happy home and had had
a bright childhood. His mother's tenderness and zest for life could
make up for the energetic father's pessimistic tendencies.
It must have gone with Arthur Pfungst in a similar way
% --- p. 120 ---
\setcounter{footnote}{0}%
\opage{120}as, according to the legend, it went with Buddha himself. Buddha lived indeed
in splendor and glory, until one fine day he caught sight of
the misery that may be in the world, of sickness, old age,
and poverty. That glance he never forgot. Pfungst too
was convinced that he had cast a glance into
life and its conditions, a glance that had to be determinative
of the light in which all that life brought had to stand.
But this glance engendered neither bitterness nor passivity.
He is said himself to have expressed his standpoint thus, that he
was a pessimist in his thought, an optimist in his work. He did
not withdraw from life, but waged a struggle to alleviate
human suffering and in particular to counteract oppression
and burden from the side of the state or the church toward
free striving in the different domains. As a native of Frankfurt
he felt himself in many respects in opposition to the
Prussian spirit, and he spoke with great joy of how
he had succeeded in warding off an attempt by the
Prussian government to gain authority over a bequest
that had been established for the benefit of a free-religious community. Through
him I once got an article accepted in the ``Frankfurter Zeitung''
on the conditions in North Schleswig, and he himself reviewed in
that paper my lecture on ``Fichte's Reden an die deutsche
Nation og de danske Nordslesvigere''. In his periodical ``Das
freie Wort'' he fought against faith in authority and clericalism
as well as against political compulsion. He felt himself indeed
a stranger toward human beings' natural faith that everything
will turn out well and that it exists for their sake. One of his
poems has for its title: ``Ich bin ein Fremdling auf der
Welt gewesen''. But when it was a matter of fighting suffering and
injustice, he did not feel himself a stranger. Even though he did not
have the naive optimist's faith that victory could be won
once and for all, his lust for work
and for battle was not weakened thereby. The struggle was to him more beautiful than
victory. One of his poems has precisely the title: ``Der
Kampf ist schöner als der Sieg''.

As so many oriental religions, Buddhism has
a passive, self-renouncing character. It preaches
% --- p. 121 ---
\setcounter{footnote}{0}%
\opage{121}compassion, but it warns also and above all against attaching one's
heart to what is finite and therefore in the end without reality.
It says in the Dhammapada (Tuxen's translation)
(Ch. 16): ``Love begets care, love begets
fear, he who has freed himself from love knows neither
care nor fear.'' And further (Ch. 26): ``He
who gives up all connection with human beings and all connection
with gods, --- he who is without all connection, ---
him I call a Brahman.'' But just as modern Christianity
became possible only because the expectation of Christ's
imminent return gradually gave some room to other
expectations and to striving for the development of human
culture, so a modern, European Buddhist too must
take up into his practical conduct of life elements that were not
found in original Buddhism. Intellectual, artistic,
social, and political tasks acquire positive value, a
value that the yearning toward Nirvana or longing for a heaven
beyond did not know. The oriental world religions
arose under cultural conditions that were very different
from the cultural conditions under which their modern adherents,
especially the western European ones, live. When Tagore was here,
he found that we here in the West have too many problems, tasks,
and difficulties. Come over to us, he said; there
everything becomes so simple. This oriental way of thinking was
far from Arthur Pfungst's. He took up the struggle where
he stood. But when he then praises the glory of the struggle, he
is more of an optimist than he as a Buddhist is allowed
to be. For even if he does not believe in victory for his humane
endeavors, he must yet believe that there are forces to
keep the struggle going, and already such a faith or
such a hope brings more light into the world than there
according to Buddhism can be. For Buddha nothing in the world of experience
had such a reality that one could
build on it hope of attaining what has true value.
Hermann Oldenberg compares Buddhism's denial
of every positive value in the world with the Brahma speculation's
``no, no'' toward every attempt to say what Brahma is.

% --- p. 122 ---
\setcounter{footnote}{0}%
\opage{122}The modern Buddhist stands toward Nirvana as the
modern Christian stands toward the expectation of Christ's
imminent return, which was the foundation of the ethics of Christianity.
A fulfillment of this expectation in the midst of the
striving work of culture would halt it, show it in
its meaninglessness. In a correspondence between two Danish
pastors, which I reviewed some years ago in ``Tilskueren''
(later in ``Oplevelse og Tydning''), the difficulty thereby
grounded for modern Christianity appears
in a naive manner. One dares not give up the primitive Christian expectation,
but how would it go, if this expectation
were suddenly fulfilled, with the progressive unfolding
of national life that the two friends both hope for?

The two great world religions in more recent times subsist only through
a constant taking up of new elements, and the great problem
then becomes how these elements can be united with the
elements that were the decisive ones in the great times of beginning.
In the very organization of worship, dogmatics, and
priesthood, without which these religions could not become
``positive'' religions, there lies testimony that there has been
taken up much and many things that were not found in the original,
and that would be disturbed and pushed aside if
this original were suddenly put into force again. From this
side Dostoevsky has in a story, whose content is set
in the medieval church, depicted the consternation
that Christ's sudden return would arouse in the existing
church, which has arranged itself for the long run and is not at all
arranged for Christ himself to take the helm.\footnote[1]{In ``Oplevelse og Tydning'' (p. 205), where I have already mentioned this story, I ascribed it to Thor Lange as author. But it stems from Dostoevsky's wonderful work ``Brødrene Karamazov''. (First Part p. 227 in the French translation). } A corresponding
problem could be posed for ecclesiastical Buddhism, which
now subsists in the East as one of the most venerable religious institutions
in the world. As long as religion builds on authority
and tradition, a problem of this kind will again and
again come up. Only a view of life that builds on
% --- p. 123 ---
\setcounter{footnote}{0}%
\opage{123}one's own experience of life will be able to draw nourishment from the past's
great thoughts and great figures without being entangled in
any contradiction. ---

But I return to Arthur Pfungst and to the present
part of his works. This first part contains
his poems. I consider them here from the side
of their content of thought, especially with regard to whether motives
for the peculiar way of taking life, which
the author came to early, are indicated in them. The poems suggest that even
though the poet's biography cannot report definite experiences
that could have conditioned his faith, he has nevertheless
gone through a strict school and reached his standpoint through
self-mastery and resignation. He states in an epigram
that there is a dominion in which the desire to rule
can find satisfaction, even if one has no subjects,
namely dominion over oneself. And in a poem whose
title is ``Zurück!'', he confesses that he once wanted
to be like other human beings, absorbed in life and its
motley manifoldness, to rejoice in sunshine and in love,
---

\begin{quote}
jedoch in meinem Herzen rief es „Nein!“\\
mein Genius stand vor mir und sprach „zurück!“
\end{quote}

In the great poem ``Laskaris'', which takes up the greatest part
of the present first part of his works, a course of
life through hope, guilt, and sorrow is depicted, and as a kind of chorus
it is proclaimed again and again, as an introduction to the
individual cantos, that ruin and dissolution threaten all that
one hopes, all that one rejoices in. Laskaris is a
young Greek from Cyprus. He is initiated by the old alchemist
Philaletes into the art of making gold, but to the initiation
the old man attaches the word that the treasure his art can
produce, and which all human beings strive after, has by no means
made him himself happier. Therefore he urges
Laskaris to use it only as a means of remedying
human want and sorrow. But the young man feels
% --- p. 124 ---
\setcounter{footnote}{0}%
\opage{124}a craving for power and honor and pays no heed to Philaletes' wise
speech. Through changing fortunes --- at Augustus the Second's
splendid court in Dresden and in Charles the Twelfth's camp ---
he shines by his art and procures honor and wealth.
But he betrays his self-sacrificing wife for the dazzling
Aurora von Königsmark, and later through his thoughtlessness he causes
the death of wife and child. At last he wants to return
to his native place, but perishes when the ship is wrecked
on the coast of Cyprus.

It seems to have been the author's meaning to impress
that only he can say No to life who has really
tried it. Philaletes was from the beginning clear that
his wise speech would not keep his young friend from
plunging into life. Only he, he says, who \textit{has} lived,
can disheartened say that life's goods are not worth the trouble! It
is a warning to young pessimists, but as it is
shaped through the poem, it applies also to young optimists,
who blindly trust that everything must turn out well. Beyond
the opposition between the two moods, which will probably always
fight against each other in the minds of human beings, points the thought
that it is above all a matter of knowing life through one's own
experience and one's own work. But from this follows naturally the
thought, which Arthur Pfungst would not enter upon, that experience of life
and work of life can never be concluded in such a way that an
absolute evaluation of life in its totality should become possible.
In an unfinished dramatic poem on the same subject as
``Laskaris'' the young Niceketas receives from the gods the right to the fulfillment
of three wishes, and his third wish then becomes: ``Let
the world perish!'' But then the gods answer: ``You do not know
life in its fullness, not the joy, not the love.
Wait then at least a year before you demand your wish
fulfilled!'' Niceketas agrees to this, but declares that his
wish after the course of a year will be the same as now.
The drama was not completed, so it is not known what
fortunes the young man was to go through, and what
verdict he would afterward have pronounced. But in the Song of the Spirits, which
closes the Prelude, it says:
% --- p. 125 ---
\setcounter{footnote}{0}%
\begin{quote}
\opage{125}Nicht im Grimme darfst du scheiden,\\
Weh', dein Herz ist liebeleer!\\
Wer da zürnend flucht das Leben,\\
Bringt ja Rettung nimmermehr!\\
--- --- --- --- ---\\
Lern' wie man das Leiden wehrt\\
\textit{Lern' das Leben wieder segnen!}
\end{quote}

When in the summer of 1909 I visited Pfungst in Frankfurt and
spent an evening with him, I had the impression that there
had come over him something more subdued and reserved.
He was not quite the same as when on
his visits up here he sparkled with wit and zest for action, with eagerness
for intellectual and practical endeavors, and when his
humor could express itself in so delightful a manner. Perhaps
he was already then suffering from the heart disease of which
he died three years later.

Through Dr. Angermann's beautiful biography and through
renewed reading of his poetry those who knew him only
through a few conversations can gain a total picture they will not
forget. His memory will live through his works and through
a foundation he established for the promotion of intellectual
and philanthropic purposes. It will live especially among those
who met him on his way, and who by immersing themselves
in the present work have won an understanding of him.

\piecerule

% --- p. 126 ---
\setcounter{footnote}{0}%

\opage{126}\essayhead{``BAKKEHUS OG SOLBJERG''}

\essaysource{Lecture to Nordic teachers at Hindsgavl \\ 30 July 1923.}

\romanhead{I}

Although I am not a historian of literature, I have nonetheless had to
promise to speak here about my late friend Troels Lund's
``Bakkehus og Solbjerg,'' and I was told that this was precisely
what was wanted, because the author was my friend. I do not
believe that the book needs to resort to a friend's contribution in order to be seen
in its significance. And besides, one may very well be good
friends and yet see many things very differently. Nor shall I
conceal what objections I have
to make, both as regards its presentation and as regards
the points of view that underlie his
judgments.

``Bakkehus og Solbjerg'' deals with the history of Danish poetry
in the first half of the last century. It
is not a learned work of literary history. Sharp criticism of it will surely be made
by professional literary critics, and
the work will pass through a purgatory similar to that through which his great
work of cultural history, ``Dagligt Liv i Norden,'' had in its day
to pass at the hands of critical historians. But his
work of cultural history did win recognition in the end.
Our foremost critical historical scholar said of
him that ``his lively power of imagination, his psychological
acuteness has served him as a researcher,'' and that ``he has found
much in the sources that was there, but that others only
saw when he brought it forth,'' and further that ``the work bears witness to
his deep love for people and fatherland, marked by
\opage{127}having been written by a man who was a fine expression of something
of the best in the Danish national spirit.''

A similar characterization can, I believe, be given of
his work of literary history, on which he labored right
up to his death and which he had to leave unfinished. It is the attempt of a man
deeply gripped by the history and spiritual life of his
fatherland to express what he owes to the rich and great poetry
that unfolded here in Denmark in the first half of the
last century. And he does not conceal that
there was an inner need in him which sought and found satisfaction precisely in this poetry.
He was a harmonious nature.
What life brought him of resistance or sorrow he did not shut
out, but let it fade away through continued hope
and striving, and he always saw light ahead. He declares
himself, too, in the present work to be a ``seer of light.''
But he wishes only to give a contribution to the great debate on the value
of life from his point of view. ``I really know only
how what was seen has been mirrored in my bright mind. But, good Lord,
let others then teach darkly. I can only teach
from myself what life has taught me.''

The period of Denmark's history with whose daily life his
work of cultural history dealt (the sixteenth century)
was a period of ascent, and his bright outlook found confirmation in many respects
by occupying himself with it. But he saw it too
brightly. It must nonetheless have borne within it the seeds of the
unhappy time that in the seventeenth century broke in over the
country and brought it close to ruin. Something similar
holds, in my view, of his work of literary history as well:
He has caught sight of a stream of light and follows
it, but does not ask whether the stream's sources are sufficient
to feed its future course. In one of his most remarkable
writings, which he called ``Livsbelysning,'' he has
given an account of his general conception of how
men's view of life has developed. This is conditioned by
the great oppositions that experience presents between light
and darkness, pleasure and pain, good and evil, life and death. In
the expressions light and darkness he sums up all these
\opage{128}oppositions, and for him the greatest advance
in the great development of mankind is therefore that the bold
conclusion was drawn: not darkness, but light is the victor!
It is in the religions of Western Asia and Greece that this great
step is taken. One begins by taking the luminous heavenly bodies
for gods. But from there thought is raised to the view that they
are only the dwellings of the gods of light, and at last the faith in
a single invisible God of light emerges. And for a long time in vain
the belief in the Devil, a remnant of the faith in dark powers
resting on fear, forms a check upon the faith in light. This
check too, which was so strong in the Christian Middle Ages,
is overcome. And as regards Danish spiritual life, this
breakthrough takes place, in the author's opinion, precisely in the
great poetry that unfolded in the first half of the nineteenth
century. In ``Bakkehus og Solbjerg'' he has wished to demonstrate
how the great world-historical movement from
darkness to light passed through one of its phases within the narrow frame of Danish
spiritual life.

The genesis of the book is characteristic of its author.
When he had completed his great work of cultural history,
he felt a need to live in the world of his own memories
and thought of writing his autobiography. This plan occupied
him for several years, and he often spoke of it. But
the more he immersed himself in his memories of childhood and youth,
the more he came to dwell on the summers
that he spent with his family out at ``Bakkehuset'' on
Valby Hill, and on the memories that were thereby awakened
of a period in Danish poetry in which this region played
no small part. A series of poetic figures
thronged forth before his imagination, who had either lived
at Bakkehuset or had had their spiritual haunt there. He
grew tired of occupying himself with himself and chose to
immerse himself in these figures, whose works had meant so
much for his spiritual development. He does himself an injustice,
however, when he says that in the face of the rich pictures
that his memory and imagination called up in him he felt
\opage{129}like the wooden stick on which the cake-women string the
cherries they sell to children. From his youth there was
an inner life of thought in him that sought expression. I recall
with pleasure the time in my young days when we came together
in my garret on Gammeltorv (diagonally across from his
father's house) and studied Plato together. Troels Lund
was much interested in this study, which in his case
led to a spirited doctoral dissertation on Socrates.
But he was more taken with the artist in Plato
than with the thinker Plato. It was above all Plato's capacity to
let a thought take shape in a picture and be animated by mood
that appealed to him. Thinking and research filled his
life from his youth onward, but received their value for
him above all as a foundation for the unfolding of mood and imagination.
And on this rested the value that cultural-historical and literary-historical
subjects acquired for him throughout his life. Often,
it cannot be denied, and this holds also of his work of literary
history, the poet in him got the upper hand over the researcher.
But in this work, which he prepared in his old age,
his mind found rest in immersion in what stood
before him as the most important thing in the spiritual life into which he had been
initiated from his childhood on, --- the great, by his people
sprung world of poetry. It stood before him as a
world of light, which had not only been able to guide and nourish
his own personal life, but also, while the World War
and its consequences cowed the countries, could portend that the powers of
light could not be overcome in the inner life of souls.

Before I go into the special content of the book, I must
explain its name, which is due to a symbolism characteristic of Troels Lund.
Denmark has no mountains, but we are all the
gladder for the hills, which make possible wide views over
our flat country, where light and shadow, because of the great deal of
wind, often change very quickly. Just outside Copenhagen
we have Valby Hill. Valby has long since been a suburb of Copenhagen,
and villas, tall houses and factories cover the
\opage{130}formerly so idyllic place, to which so many memories
from the literature of the last century attach. Oehlenschläger
lived as a young man at Frederiksberg Castle, where his
father was castle steward. In the parish clerk's house by the church
he wrote ``Helge,'' one of his very finest works, and in the nearby
Bakkehus, at the Rahbeks', he had made the acquaintance of
his bride, Kamma's sister. In his old age he had a summer residence
at Fasangaarden in Frederiksberg Gardens. In Bakkehuset
on the eastern slope of the hill lived the Rahbeks, whose
home for a long time was a center of Danish spiritual life.
Almost all renowned poets frequented it, and after
the couple's death both Heiberg and Grundtvig had
summer residences there. That Troels Lund's parents lived there for a
number of years in summer brought the place and its history
so vividly before his imagination, when he immersed himself in
the memories of his life, that instead of an autobiography he came
to give an atmospheric survey of the history of Danish poetry
in the nineteenth century.

Solbjerg was the name of a village that in olden days
lay on the hill, but had to yield its place when Frederiksberg
Castle was built. For Troels Lund Solbjerg means a
mountain or a hill where the sun was worshiped. And both Bakkehus
and Solbjerg became symbolic expressions for him. Bakkehus
means a home on the way up toward the worship of light
that extends through the history of mankind, and that,
according to him, worked its way forward in a peculiar manner in
Danish poetry in the nineteenth century, a small example of the
striving toward the light that is the innermost essence of the human spirit.
Troels Lund has kept to this example because it
was the one he knew best. Other peoples he knew, he says,
too little to be able to write about their homes
of light-worship. Not only here in the North are there many Solbjergs
and many Bakkehuses, but every place receptive to the light
is a little Solbjerg.

% --- p. 131 ---
\setcounter{footnote}{0}%
\opage{131}\romanhead{II}

According to Troels Lund's conception, it lay in the nature of the circumstances
here in Denmark at the beginning of the nineteenth century
that poetry had to become ``the boldest and most knowing
guide'' in the spiritual domain. ``It went ahead and
often uttered the captivating and liberating word, while science
still hesitated and the church anathematized.'' ``Through
poetry the spiritual development in Denmark has proceeded most easily,
while science, theology and philosophy followed
behind.''

In these words the basic idea of the work is given. And it lies in
the quoted words that the darkness which the light is to defeat is not only
that which is due to misfortune, pain and sorrow, but also that
which is caused by the power of dogmatic traditions to bar
access to the light that the human spirit has in itself the strength to
reach. The work then gives a gallery of the country's poetic figures,
in part with brilliant characterization and with a multitude of
individual traits. Troels Lund had, both as a historian of literature
and as a historian of culture, a great gift for entering
into the figures of the past and letting imagination furnish
them with a wealth of traits. As regards the details
one asks --- in this one must grant his often so severe
critics that they are right --- often where he actually has them from, ---
and often asks in vain. But what is the main thing for
us here is the series of seers of light whom he wished to present to
us. It is impossible to reproduce the presentation in detail;
we must content ourselves with a few examples.

First the Danish poet king, Adam Oehlenschläger,
is to be mentioned. In him a new view of life already announces itself in the youthful poem
``Jesu Liv i den aarlige Natur.'' Whereas in the church
one has often illustrated its doctrine by pictures from external
nature, the poet here uses, conversely, the church doctrine's
ideas as symbolic events in nature. He
appeals to the readers' religious childhood learning in order to render
the life of nature intuitive and vivid to them. This
marks a turning point. Nature begins to come into
\opage{132}its own. One follows its direction, and the ecclesiastical ideas
are used only as supports for the imagination.

Special attention is given to Kamma Rahbek, who was called by Ingemann
the mother of the Danish poets. Not only did she accompany
young poets on their way with heartfelt understanding, but
she carried poetry through in her life, was an evangelist of nature, in
constant development and striving. In her loving mind and in the
hope that builds on love she found the solution
of the riddles of life. She was a religious nature, but did not go along
with the sharpening of orthodoxy that came in the wake
of Romanticism.

This very orthodoxy, with its strong emphasis on the doctrine of the Devil
and of hell, was overcome through the great seers of light.
Just as Denmark's greatest jurist, A. S. Ørsted, in keeping
with his general view of how the legal provisions
of the past should be applied to the conditions of new times,
maintained that confessional writings of the church, which
themselves owed their origin to a particular age, must be interpreted
according to the need and experience of the new times, so
even Grundtvig, the new orthodoxy's crassest prophet, gradually became,
above all under the influence of his heartfelt participation
in the awakening national life here at home, a great seer of light,
for whom the doctrine of the Devil did stand as a historical presupposition
for the teaching of the New Testament, but was not a
tenet valid for all times. Thereby he became a
great spiritual liberator and leader for so many among the people, a
mighty champion of true enlightenment.

But the high point of free view of life is reached, according to Troels
Lund, in Ingemann, Hauch and Paludan-Müller.

As regards Ingemann, ``Holgers Sang til Livet'' is above all
to be mentioned. It is ``a hymn to life, as full and inward as
had not yet been heard in Denmark.'' In it Troels
Lund found the best expression of his own bright view of life, culminating
in the last verse:

\begin{quote}
But when under the sun the soul does not get what it wants,\\
then there are other suns and other stars.\\
And if all suns and stars are someday extinguished,\\
yet the spring of life flows, where eternally it flowed.
\end{quote}

% --- p. 133 ---
\setcounter{footnote}{0}%
\opage{133}Troels Lund rightly emphasizes the great infinite
horizon in Ingemann's poems of life. While the Middle Ages'
limited picture of the world still governed Grundtvig's
ideas, there can be traced in Ingemann, as later in
Hauch and Paludan-Müller, the great expansion of the picture of the world
that had become the consequence of Copernicanism. Holberg
had already --- with much caution --- expressed himself about
the influence that Copernicanism must have on
the whole conception of the world. But I believe that Troels Lund is
right that only in Ingemann, Hauch and Paludan-Müller is
this influence traced in Danish poetry and in the religiosity
that found its expression in Danish poetry.

In Hauch similar thoughts appear in a more inward
manner and in harmony with his lively sense of nature. He
is the poet of fidelity and of strict honor. Paludan-Müller's
faith is a religion of hope. For him hope is not merely
the eye of the soul, but the kernel of life in us. The poem ``Nattevagt'' (1837)
contains his great faith in light and life, despite all
illusions and limitations that meet us on the way. Later
he described in ``Adam Homo'' man's weakness,
self-deception and wretchedness, but also the power of love
to raise even the most wretched soul upward toward the light: Alma, who
herself has access to the kingdom of blessedness, nonetheless follows Adam
faithfully through purgatory.

Troels Lund's beautiful work remained a torso. It was broken off
in the middle of the third part by his death, the death he went to meet with
cheerful courage, as befits a seer of light. This third part
was to deal with Søren Kierkegaard, but got only as far as
relating his youth. The continuation would have set
Troels Lund a great problem, for Kierkegaard declared
not only the poetry of Romanticism to be a mere dream, but maintained
also that the idyllic conception of Christianity,
on which the Romantics built, was a deception. I do not know
how Troels Lund in the face of this mighty figure
and this powerful protest would have been able to maintain
and carry through his doctrine of the victory of the bright view of life
in Danish literature. The sharpening of the problems of life that
% --- p. 134 ---
\setcounter{footnote}{0}%
\opage{134}Kierkegaard brought about demanded in any case a serious
trial. Perhaps one can from a hint at the end
of the second volume suppose that Troels Lund essentially saw
Kierkegaard's significance in his overthrow of the visible
church. Then there could be more room for poetry's
free flight; but it would indeed then also have to take place on
other terms than before.

\romanhead{III}

Only a few words must I add, from my own standpoint,
to the characterization I have tried to give of my
late friend's work. It is a great task he had set
himself, to show how a lofty poetry of life developed among
the poets of the Danish Golden Age. And with great gratitude
I too recognize how much we owe in this respect
to these poets. But it was the age of Romanticism.
One took the realities, in both nature and history,
lightly, and poetry hovered high above the waters of actuality.
It cannot be said that Romantic poetry was the precursor
of science and church doctrine. In many respects it gave a reinterpretation of the ecclesiastical ideas, and
its boldest flight was due to the Renaissance's expanded picture of the world,
which only now was properly brought forth outside the
domain of the special sciences. And poetry cannot at all
be a precursor of science in the sense
that it should be able to explain to us something that science
cannot yet explain. It cannot fill in
the gaps of science when it is a matter of understanding or explanation
of reality, just as little as science can ever
reach up to the height of poetry when it is
a matter of expressing the mood, the emotion that
the great oppositions and collisions of life can arouse.

The poetry of life that we need, and that every age needs
in its own way, must not, like Romanticism, hover over
reality, but have its root in free and deep personal
experiences, spring from minds that have been tried in the struggle of life and
% --- p. 135 ---
\setcounter{footnote}{0}%
\opage{135}of thought. It must not suppose that it can solve problems
at which thought has come to a halt, nor must it
be merely a conversion into poetic form of received ideas,
but in a bold upsurge it is to give great symbols
of that struggle of life, that clash of light and
darkness, which will endure as long as life lasts. Light can
be victorious, but the victory must always be won anew. That was probably
really also Troels Lund's opinion, though he often seems
to revel in the pure light, where all oppositions have faded away.
He has said that it is precisely in the dark times
that faith in the light is to show itself.

For my part I agree with him that there are in
any case two wonderful examples of poetry of life
that we can take with us from the Danish Golden Age, --- Ingemann's
``Sang til Livet'' and Paludan-Müller's ``Nattevagt.''
The manliness that the first poem expresses, and the deep feeling of unity
for which the second finds words, are important elements
in every strong and inward feeling for life. If we have found in life,
as we know it, and as we have pondered it with the
thoughts at our disposal, something that was
worth working and fighting for, and have put our best
powers into this, then we need no external supports for
our hope. We can then say with a poet from a somewhat later
time than that from which Troels Lund took his examples:

\begin{quote}
Fight for all that you hold dear,\\
die, if it comes to that!\\
Then life is not so hard,\\
nor is death either!
\end{quote}
\piecerule

% --- p. 136 ---
\setcounter{footnote}{0}%

\opage{136}\essayhead{THE BOOK OF THE PHILOSOPHERS' STONE}

This is a book about growing up. It therefore begins
naturally with childhood. And the way in which the children's
life in the little village is depicted will perhaps for many
readers be the best in the book, that to which they most like to turn
back. Not because the children's life is depicted as sheer
joy and harmony. There too we find both coarseness and
gentleness, and even the gentlest mind can be seized by the savagery of a beast of prey
when coarseness has outraged it. But in childhood
the difficulties are resolved relatively easily. The horizon is
limited, the tasks fit the abilities fairly well, and
the multiplicity is not greater than that a receptive mind
can think it sees it as a unified whole. Those who are not
receptive in this way do not for that reason lose the coherence
of their minds. They meet every wave that life sends
toward them, and swim on upon its back. They do not ask
``Why?'' Their motto is ``Why not?'' ---
if they have any. They follow the tradition into which they have been led,
and their trust in it is not shaken, precisely because they
take things as they come, and have no need
to cast a glance out over life as a whole. Thus the book depicts
a young couple, Ellen Nielsen and Hans Olsen.
One day in a writing lesson at school they happened to look up at
each other and thereby smiled, and since that day the one
of them has only laughed or wept when the other did.
One fine day, after they had grown up, they happened to
talk together about this and agreed that it would always
go so with them. Then they wandered hand in hand
through life, in common happiness and in common work. They believe
what they have learned; why should they not? And
% --- p. 137 ---
\setcounter{footnote}{0}%
\opage{137}they are even sure that, even if they had not learned it,
they would nonetheless have believed it: ``it comes over them of itself.''
For others it does not go so easily. There is the mild and
helpful Holger Enke, who, in exasperation that the
young woman whose purity and innocence he admires is seduced
by a libertine who takes flight, himself rapes
and murders her, and then after a failed suicide willingly
goes to the penitentiary. When he is pardoned on account of good conduct,
he returns to the village, takes people's fear
and contempt as a part of his punishment, finally wins new esteem,
and in the close of the book proclaims the wisdom of life he
has reached. To him we shall return later.

It is the clerk's son and the pastor's son on whom we must for the time being
dwell. The clerk's son, Jens Dahl, through his personality
and his thoughts poses the book's real problem. Already
as a child he learned from his own experience the great difference
between what he called ``the Open'' and ``the Closed.''
For him who is ``open,'' differences of time and place
do not hold. Everything is near. His little brother, who is dead, is constantly
present. It is not possible to express ``openness'' in
words; the one would after all have to follow upon the other, which is an
imperfection just as great as that one must take step
after step in order to get from one place to another.
In ``the language of heaven'' we say everything, without a single word. Here
all riddles are solved --- because they have not arisen at all. But
when one is ``closed,'' one must think things over and go
step by step. And then all the difficulties come
and the laborious work. The seventeen-year-old boy felt
himself, by virtue of his ``openness,'' to be a genius, sure
that everything that was ``closed'' could be seen through
and overcome from out of the ``Open.'' And it did not warn him
that he fell for the first temptation he met. But
when he later came back to the village, he had had
experiences that let him look at his old comrades
partly with melancholy, partly with envy, on account of the
naive, involuntary way of taking life that they had retained.
They went from step to step on the ladder of life, and the one
% --- p. 138 ---
\setcounter{footnote}{0}%
\opage{138}step followed of itself upon the other. In their decency
and unconsciousness they were not far from the world of the ``language of heaven.''
--- Before we follow him further, we must tell
of his friend, the pastor's son.

Christian Barnes had early lost ``openness.'' The
insight that he had gained as a schoolboy in the market town into
what can hide behind the ``closedness'' of the grown-ups had
made a strong impression on him. And an even stronger effect was had by
a secret vice into which he had fallen. His need to
analyze himself and others led him away from the ``language of heaven.''
But although he had left his childhood faith, he
nonetheless had the conviction that there was a natural religious
feeling that does not stand and fall with any particular form of religion.
And he gives a long lecture on the philosophy of religion about this
to his friend Jens Dahl, when this one too has
come to Copenhagen as a student. (There are surely too many
lectures in the book. But I leave it
to the literary critics to complain of that, if it is to be complained of.)
What Christian Barnes seeks (all the while he is preparing
for the master's examination in English) is to become healthy, to
get the sound view of life that leads to free religious feeling.
He also succeeds in this under the influence of an energetic
English clairvoyant, who sends him to America in order that
as a cowboy he may find out what physical fatigue
is, and thereafter acquire a surplus of vital energy that
can be used for sound thinking. At the same time he learns from her
to collect his mind in stillness and devotion. On the other hand he does not find
in America the philosophers' stone, indeed he even thinks that
it is not worth having, after it had been provided with high wisdom
that it was well hidden. When he has come home
from America, he meets again the woman who as a little
girl talked with him over the garden fence, when on Thursdays
he had an hour between classes. She had not forgotten him.

Before he went to America, he had given his friend
Jens Dahl the fateful advice to study mysticism and
theosophy. He said to him: ``You who want to enter the Open
and live in the world of the language of heaven, which is hidden from
% --- p. 139 ---
\setcounter{footnote}{0}%
\opage{139}ordinary mortals, are you not, quite like these [mystics
and theosophists], one who hopes to find the philosophers' stone?'' ---
Jens Dahl was prepared for this study through his
childlike attempts to understand ``the Closed'' --- i.e.
the manifold, the surprising experiences, the hard
and difficult in life --- by ``the Open'' --- i.e. the
involuntarily emerging picture of the whole that can form
in a childlike mind. Now he gave up all other studies in order to
immerse himself in the old mystics. And he soon felt
at home in their world. It was as if he here met
his own again, --- as if he were called back to the days
of childhood.

It is perhaps not unnecessary to explain what the
philosophers' stone historically means. The expression stems from the alchemists,
who sought a substance that could be transformed into
all other substances, above all into the most precious, silver and
gold. And this wonderful substance, it was then believed, could
also cure disease and prolong life. All material
goods could be obtained by it. The alchemists were also
often mystics, and conversely. But the mystics who are
in question in the present book sought not a substance, but a
thought, the thought of what lies at the ground of everything and
holds the world together, --- of something that the popular religions
have sought to express in their figurative language, but that
one grasps properly only through energetic spiritual concentration,
especially when one can become one with it and see everything in the light
of it.

Jens Dahl thought, as said, that he here found again what he had
experienced in the ``open'' hours of his childhood. Yet the
experiences the mystics spoke of were of a higher and more difficult
kind, and only when he had received instructions for
daily spiritual exercises in the direction of concentration and self-absorption,
did ecstasy come over him one day. After a severe
% --- p. 140 ---
\setcounter{footnote}{0}%
\opage{140}self-examination that let him see through his whole wretchedness,
he finally saw his personality, as he had hitherto known
it, sink away into nothing, but felt himself in return as borne
by a loving power, the same that bore the whole world,
and not only borne by it, but one with it, without
any thought, any word being able to express the light that
streamed through his whole being. He now thought he had
learned by experience what his friend had so often talked about, --- the
religious feeling that lies deep in the nature of man and
is in itself independent of all dogmas and all external forms.
Through his strong concentration he also gained the capacity
to influence others through hypnotizing and suggestion.

Dahl was convinced that what he now experienced was
something purely natural, conditioned by definite causes. It
must therefore, he thought, be possible to bring about the
ecstatic state again. But only imperfect and fleeting
glimpses emerged in him. The constant, convulsive exercises
made him nervous; the forced exertion caused
only uncanny hallucinations, and at the same time sensual
drives were aroused, just as had already been the case in
his years of transition, precisely when he had been most
``open.'' In this state he was impressed by the calm and
assurance with which his theosophical friends carried through
their methods. They went through a kind of preliminary school and could
thereafter receive direct instruction from masters from the
world of spirits itself, emissaries from the great Indian leaders.
Dahl then also went through a new training in ecstatic
dreams, in which he moved on the ``astral plane'' among
the departed, whose time in purgatory he witnessed. Soon it
came to the point that he did not know the difference between dream and waking.
He had hallucinations in the street and had to be conveyed
home by the police. In a moment of reflection he saw
that he had misused the ``open'' states that
had been granted to him, in that he had sought only enjoyment in
them --- just as he had treated the women he had
met on his way as herbs from which he had
extracted a precious elixir, while they had individually been nothing
% --- p. 141 ---
\setcounter{footnote}{0}%
\opage{141}to him. In vain had the English clairvoyant
admonished him not to strive so much to see the angels of Heaven, but rather to see to it that the angels he met here
on earth got their full and proper significance in his life.
He had, as he himself finally saw, ``made
God's garden into a pleasure garden, and it was therefore denied him
to work in it.'' --- One day, as he sat on the shore of a
lake, he saw on the opposite shore a female figure
beckon. He wanted to go straight across to her, but the water
closed over him.

If one asks what human interest the depiction
of Jens Dahl's development and downfall has, it must doubtless
be sought in this, that in a single example there are presented the dangers
that threaten in the transition from the dreams of childhood and early
youth to the reality of life and the demands
of life. There are, as we have heard, those who in a way for them
selves imperceptible manner glide over into the reality of life.
Dream and goal work in them as mildly animating forces,
and it does not come to a crisis, to say nothing of downfall. But
there are natures for whom the ``openness'' of childhood, its
naive faith in the ease of solving the riddles of life, leads to
hard and fateful collisions. They do not see that
it is only through serious work, a work that must be
chosen in accordance with personal need and ability,
that the conversion of the dreamed ideals into reality can
take place. And it is then not merely the given reality
that is to bend before the dreamed ideal; it is quite as often
this ideal that must assume new forms in accordance with the new experiences
that force themselves upon one. Here an inner struggle is waged that
is not always noticed by others, but that can be of great significance
and mark the innermost life. In all domains
there can be occasion for such a struggle, --- in science,
art, religion, in the relation to individual people and to
society, and not least in the domain that at all
% --- p. 142 ---
\setcounter{footnote}{0}%
\opage{142}times has been a favorite theme of poets, the erotic
relation between man and woman. The book about the philosophers' stone
also indicates a parallel between Jens Dahl's relation
to the reality of life and his particular relation to the women
whom life brings in his way. It could have been examined here
whether there was not a direct connection between
mysticism and eroticism. In Jens Dahl eroticism appears rather
as an opposite to mysticism, a kind of contrast effect.

But the reader cannot avoid asking himself why
he is to be led into all this mysticism and theosophy, which
makes it necessary for the author to let one of his characters
give a lecture, so that the reader can be à jour with
what is to be presented. Does life not offer relations
enough that can be used to illuminate the significant
transition that is the book's real theme? It is possible
(of this I cannot judge) that mysticism, yogism and theosophy
play so great a role in the human life of the present that it
might lie near for an author to choose his subject here. And
it cannot be denied that the author gives characteristic
pictures in the mystics and theosophists with whom he brings the young
man into contact. They are living figures each
in his own right, and each takes his religion in an individual
way. It is moreover one of the theosophists who utters the
most liberating word about the philosophers' stone. He is a poor
craftsman who has abandoned his wife and his children
and has combined his theosophy with striving to
produce a perpetuum mobile. In the end he comes
to see the impossibility of both, and he returns
to his craft, where work can and must go piece
by piece, step by step. And when he is asked whether he
does not believe in the philosophers' stone, he answers: it is a milestone
on the way, --- and when asked, which milestone, the answer is:
the next one on the way and always the next one! It is Voltaire's
old watchword about cultivating one's garden, only that the
good theosophist has seen that the garden does not always lie at the
same milestone.

% --- p. 143 ---
\setcounter{footnote}{0}%
\opage{143}The significance of the book would in my eyes have been greater
if it had moved more in the daylight of life, where
in the end all real tasks lie, even if we may need
the brooding of the night and the stillness of the night in order properly to solve them. And
even if one precisely wishes to depict how one can dream
oneself away from real life, one need not fetch
one's examples from theosophy. There are unfortunately examples
enough besides. ---

As regards the depiction of the mystical phenomena, there
are a couple of remarks I cannot refrain from making.
It ought to have played a greater role in the presentation that
ecstasy, even if preparatory exercises have been made, nonetheless
always comes involuntarily, as a surprise, and that it
is hopeless and self-contradictory to build one's life on the attainability
of that kind of states. A closer study of
biographies of mystics leads moreover to the conjecture that
ecstatic states are tied to certain periods of life
and fall away after their close. In the famous ecclesiastical
mystics this does not come out so strongly, because they
with great ease glide over from experience to interpretation. And
since they take their interpretation from the ecclesiastical ideas
as a matter of course, they do not always lay strong weight on the
difference between the actual ecstasies and the vivid and
captivating moods into which they can come when the
mighty images of ecclesiastical tradition strongly seize them.
Differently will \textit{he} be placed who, without such support, wishes to
live himself into the world of mysticism. For him the opposition
between what one might call the great
and the small moments will be unendurable, and it will be able to be overcome only
when work is undertaken to carry through in the valley of experience
what has been beheld up on the heights. But
it is precisely such work that modern mystics
and theosophists shun. Mysticism has brought great and good things
with it, above all by leading away from blind absorption in
dogmas. It has often led precisely from dogma to life. But
it can itself become a danger, and it is perhaps that for
many in our days. However much we must hope that
% --- p. 144 ---
\setcounter{footnote}{0}%
\opage{144}concentration, enthusiasm and self-absorption will never cease
to exert their fertilizing influence on human
spiritual life, it is nevertheless not always fortunate that they assume the form of ecstasy,
and they will presumably also more and more seldom
do so.

There is one person in the book whose childhood history is not
told. It is ``the Candidate.'' Occasionally we learn
that he has made journeys around the world. His
adventures on this journey are said not always to have been innocent;
but he can nonetheless say that he has nowhere done
anything that was contrary to the laws of the country concerned. Now
he has settled down in the village from which he comes, and
has bought his late father's house. He finds the region to be
the Garden of Paradise, if one only looks at it with the right
eyes. Each one must procure these eyes by his own
way. With this outlook on life the Candidate then wanders
about the region and helps great and small along, when
things get tight for them. A harsh missionary he makes ridiculous
in the eyes of the population. A married couple he helps in
a cunning way out of the ``incompatibility of temperaments''
that threatens their life together. A young woman who, overwhelmed
by the wickedness of life, wants to take her own life, he saves
at the last moment. He strengthens people in their religious conception,
when he notices that it really accords with
their interest and need. The pastor of the parish, who for the rest is very
well disposed toward him, says of him that he always has the
religion that the person he has just been talking with has need
of having. It is the Socratic art he practices, and despite
the sympathy that stirs in him, it is a pleasure to him
to practice this art, and he can ``chuckle'' when it
has succeeded well. This Socrates of the village is a
highly enlivening, and often also touching, element in the narrative.
One is always glad when he takes matters in hand. If I
remember rightly, he already appears in one of the author's earlier
narratives.

% --- p. 145 ---
\setcounter{footnote}{0}%
\opage{145}If the Candidate did not always take his pastoral work
solemnly, there was nonetheless one figure in the presence of
whom, despite all his art and all his knowledge of human nature, he
came to feel small. This was Holger Enke, the rapist
and murderer, he who after having come
out of the penitentiary went his pilgrimage in his home district,
where his deed had been done, and slowly worked his way forward
despite suspicion and contempt. What he had to suffer was to him
a relief, and in ecstatic rapture he even declared
that he willingly went under eternal damnation, if it
should be so. Forgiveness he did not seek. The Candidate
drew from this the conclusion that Holger had overcome hell,
had come out beyond the opposition between sin and
grace and rested in the assurance that everything he had
been and done was an object of the eternal understanding of the Deity.
Psalm 139 gave him his philosophy of life: From
the hour of his coming into being everything was written in God's book. The days
were determined before they came, also everything that happened in the days.
This eternal thought, the source of the great laws of the course
of things in the world, he bowed before, and what it taught
him he applied not only to himself, but also to
others. When he then met the man who in his time had
seduced the young girl he admired and loved, and who
thereby had become the cause of his crime and misfortune,
he let him calmly go his way with the words: ``I did not
create you.''

In the last chapter of the book Holger develops his whole pilgrim philosophy
for the Candidate: Even if one came up to Heaven,
one would nowhere find God, precisely because he is everywhere
and permeates everything. When man stops on his way
and thinks himself to be something, he thereby makes God into nothing,
and every thought we make of God removes us from him.

Doubtless we hear here rather the author than Holger
Enke. It is the author who in Spinoza's spirit converts
Psalm 139 into a speculative system. The real
Holger would have stopped beforehand and held to the
great images that his religion presented to him. Perhaps
% --- p. 146 ---
\setcounter{footnote}{0}%
\opage{146}he would, without being aware of it, have interpreted his
images in his own way, on the basis of what he
had experienced, forfeited, suffered and understood. But abstract
concepts like being and non-being he had hardly wrestled
with. And the Socrates of the village would have been gripped by
the depth of his mind and the power of his spirit without it.

I have set myself the task in this review of
examining whether the author has succeeded in converting the
understanding of human spiritual life, particularly of its religious
forms, which he doubtless possesses, into real poetry.
I do not believe that he has everywhere succeeded. My task
has made me keep to the principal characters and
their thoughts. But around them stands a series of figures,
each with his life course and his way of taking life.
And here I have found genuine poetry. When I again take up
the book, as I doubtless shall, I shall first and foremost
seek out these minor characters, who appear partly in the village,
partly in the market town, partly in Copenhagen. In the depiction of
them a knowledge of human nature and a love of mankind
are displayed, which for me determine the mood
the book leaves behind. It lay only outside my task,
unfortunately also outside my ability, to characterize these
figures.

\piecerule

% --- p. 147 ---
\setcounter{footnote}{0}%

\opage{147}\essayhead{ANKER LARSEN: ``FOR AABEN DØR''}

Attentive readers of Anker Larsen's earlier books
will find this book, for all that it brings of the new and fresh,
no surprise. The subject which in his earlier
books was treated, in the eyes of many, with the greatest love
and skill, the life of happy children in the country, plays
a decisive role in the new book too, though it appears
and works more directly. The child's full surrender
to its states and moods, its inward absorption
in what single moments can give, its untouchedness
by criticism and reflection make it for the author an ideal
for adults, whose minds are often weighed down or narrowed by
difficulties and problems that do not always have their
good ground. The road to Hell is in our day often paved
with groundless cares.

Already in the author's perhaps most widely read book, ``De Vises
Sten'', and no less in ``Martha og Maria'', the child's mind
stood not merely as a poetic subject but was held up
as the expression of a state in which we could open ourselves
completely and feel wholly free. There he already takes up the old saying
that we must become children again. The fundamental thought of ``De vises
Sten'' is given in the opposition between the open and the
closed. For him who is ``open'', differences of time and place
do not count. He lives wholly in a valuable Now, and everything gathers
for him into a whole. In the ``closed'' states, by contrast, one is
captivated by distances, by obstacles, by a multiplicity
that cannot be gathered into harmony.

Already by its title the new book recalls what was
a main subject in the preceding ones. The difference is that this
subject is now taken up more directly, partly in autobiographical
% --- p. 148 ---
\setcounter{footnote}{0}%
\opage{148}form. One follows the author's presentation with great interest,
and it is of no small psychological significance. Anker
Larsen relates how, at a certain point in his
life, after he had tried both theology and theosophy, and
after he had become entangled in what he calls ``daily-newspaper culture'',
a need arose in him to relive
what had made his childhood so bright, so peaceful
and so whole. More and more the longing for the
vanished time filled his mind. ``The first step on the road home''
he came to take one day during a walk in a forest.
He felt at once full of health and strength and intoxicated
by the mood of the forest. In this paradisal mood the
image of a lovely forest path arose in his soul. It appears
from his account that this image arose only
after the mood had unfolded. The joy in the image
then increased the very mood whose symbol it became. He
had had his fill of looking at what surrounded him and had
withdrawn his eyes, and only when he looked up again did he see
before him ``a little forest path so fresh, pure and fairy-tale-like that it
must have been one of the paths in the Garden of Paradise''. Only later
did he become aware that the path belonged to the region of his childhood.
There was really no particular likeness between the
path he was walking on and the one that now stood before him. --- I believe
that his description must be exact. It was the mood
that overwhelmingly became the cause of the image's arising,
and the likeness between the two paths was in any case
of subordinate significance. On the other hand there
was a likeness between his mood now and his moods from
his childhood, yet in such a way that the former did not reach the latter
in full life until the sight that had been tied to a
childhood experience arose involuntarily. The need that
in every strong mood stirs for expression, a kind of
mental mimicry, led to a sort of vision, in that the sight of what was
present at the moment did not sufficiently satisfy
it.

The older mystics distinguish sharply between the experiences
themselves and the images in which these experiences find expression.
% --- p. 149 ---
\setcounter{footnote}{0}%
\opage{149}One can infer from a likeness between the images that arise
in different mystics a likeness in their
experiences. In my book ``Oplevelse og Tydning'' I have particularly
brought this out in the comparison I make
between two women mystics, the Spanish nun Teresa
in the sixteenth century and the remarkable Protestant mystic
in the twentieth century, whom her psychological
confessor calls Cécile V. There is an analogy between Teresa's
image of the iron that glows in the fire but otherwise lies as
a black mass, and Cécile's image of the plank in the sea water,
which now shines lit by the evening sun, now is seen merely as
a black board. Both mystics have felt the oppositions of life and of the mind
so strongly that these oppositions involuntarily
found symbolic expression in their visions. On this
point there is a characteristic difference between Anker Larsen's visionary
experiences and those of those other mystics, an opposition
to which I shall return later. Of course in
his case too it is a sharply felt opposition that gives rise to the need
and the mood. But, by his account, it has
precisely been the striving to get beyond division, toward
unity and toward whole and free life, that has led him into
the states which his visions symbolize, and he has not
felt a need to let the very oppositions he strives
to get beyond find expression in the visionary symbols. He longs
for the naive, undivided life of childhood, and he finds satisfaction
in dwelling on an image that makes him
forget everything else, the strife and the divisions too. ``The deep
happiness of reality'', he says, ``fills me; my state
spread out and became One with all similar states''.
--- And so he walks, free and glad, along his reborn
forest path.

The author lays strong emphasis on the fact that what he saw before
him were not memories that could be called up at will.
The glimpses came of themselves, and the state itself had
the character of a presence, not of fantasy. The state took on in
him its character of having been called forth by longing for the
holy simplicity of childhood, and the glimpses acted as
% --- p. 150 ---
\setcounter{footnote}{0}%
\opage{150}exhortations to begin life afresh. They had, as he expresses
it, a will-determining power, set him a goal
that cannot be changed. This goal is to take up life as
it now happens to lie before one, --- to will one's fate.
If one meets resistance, it is regarded as a means, an
instrument for overcoming oneself. The freedom that broke through
in the single glimpse becomes real only by being held fast
in the course of time.

Gradually the glimpses took on a more lasting character and led
to a clear vision that lifted above time and space and opened
the way to eternity. ``Eternal life'', says Anker Larsen,
``is not a life on the other side of death. Eternal life
is a state in the soul of man; it is life's deepest
reality, which unfolds in his whole being and before
his wondering eyes. Alas, all words are useless for describing it.'' It seemed as if the things of the world too
now opened themselves to him and showed him their whole true reality.
The old ``closed shape'' was only a
disguise, not their real reality. To the
``closed'' shape belongs especially the opposition between life
and death. It falls away for the open gaze. The author
recalls that in one of his books he lets a boy sit
under an elder tree and speak ``the language of heaven'' with his departed
brother. That boy has been in ``an inner state
in which all his experiences were present'': ``There is
no distance; everything is \textit{here.} There is no past and future,
everything is \textit{now}.'' ``Death does not exist at all.'' ``He who has
experienced the Now sees no bottomless abyss between `Life' and
`Death'.'' ``I am not interested in duration, but in
eternity.''

The thought Anker Larsen strikes here recalls Shelley's
``Sensitive Plant'', where the poet maintains that beauty
and love are raised above the oppositions between life and
death. The profound English thinker Herbert Bradley, who died a few years ago,
expressed his assent to Shelley's
beautiful poem, and he adds, what will appeal to our author,
that the paradise in which those things grow is not
% --- p. 151 ---
\setcounter{footnote}{0}%
\opage{151}any region or place in time and space, but that it is \textit{here,} ---
that it is everywhere where a finite being is raised up to the
higher life that is now waking reality. ---

The openness that was won, as lasting illumination
replaced the single glimpses, gave rise not only to visions but
also to a need for exclamation, for words, --- clear as the author
otherwise was about the inadequacy of words in face of the best
and the highest. He found an outlet in the exclamation ``God!'' This
word contained for him ``jubilation, being, presence, personal
relation''. What was to be expressed was that ``something mighty and good
came to him''. It was like a deep soulful
breath, --- as if warmth, strength and security
streamed into the soul. But not even the language of the gods would
be able to express such experience, and the author does not wish to
discuss the various theological and speculative concepts of God.
``Here,'' he says, ``it is a matter of a practical and actual
life, and not of views of life.'' What has been experienced gives rise,
as we have heard, to a need and courage to let the Eternal
(the Now that \textit{is}) permeate life in time (the Now that
\textit{happens),} to clarify and purify the limited ends, so that work
stands as a good, and is not merely practiced for the sake of utility.
The life-feeling thus called forth becomes devotion,
without special hours of devotion being needed. ---

Anker Larsen feels akin to mysticism, as
it appears within the various religions. He
often refers to the writings through which we come to know it
as his ``family documents''. Yet his mysticism
is of another kind than that which appears in those he
names as his spiritual kinsmen. It is a kind of open-air religion,
and ``the open door'' may often recall ``the
green field'', in which one of our finest and boldest lyric poets
found expression for his thoughts about life. And yet
it seems as though he wishes to withdraw from life and its struggle;
he is too quickly done with life in the world of change
by speaking of ``daily-newspaper culture'' and ``interests of the day''.
And it is no accident that the images in which his moods
find an outlet are images of rest. The ``forest path'' lies still,
% --- p. 152 ---
\setcounter{footnote}{0}%
\opage{152}and there is no work to be done there. This is connected, as already
remarked, with the longing for the idyll of childhood
that was a main motive of his mysticism.

And yet it turned out that what he experiences in the highest
moments is not enough for him. The eternal Now is for him
not fully real; it needs to be ``realized'' in
time. Thus it is (what most mystics, by the way, forget)
not only the temporal that needs the Eternal; the
converse holds too. And it is a main difficulty
for all mysticism (and for all speculation) that even if a way has been found
from time to eternity, none has been found from eternity
back to time. In the great, gathering moments it is
as if life and all things open themselves. But why are
they otherwise closed, --- and why is our own mind so
often closed? How can the closed state be brought
to an end? ---

Anker Larsen wants to give us experiences, not views of life,
thus no interpretation. He has too great a contempt for views of life,
which in his eyes are only fruits of idle speculation.
Description and interpretation can nevertheless be hard
to separate. He himself uses images and thoughts that pose
problems which in the long run cannot be overlooked. But it
must be said to his credit that what we hear from him through
the open door bears throughout the stamp of genuineness and of self-experience.
With intellectual honesty he has taken care
that neither handed-down opinions nor undigested
reflections come to inhibit the originality of the experiences.
And he keeps all hasty interpretation away, as far as
he is able. In this he has an advantage over Ludvig Feilberg,
who as the final interpretation of his observations had a whole
theological physics, to which few readers have, I believe, paid
much attention.

By his own declaration Anker Larsen has not been influenced
by Feilberg. Yet in many points there is a kinship.
Feilberg's ``values of possibility'' stand just as separated
from life's great work as Anker Larsen's ``eternal Nows''.
Common to both is above all the fresh, unprejudiced
% --- p. 153 ---
\setcounter{footnote}{0}%
\opage{153}gaze and the ability to give us immediate experiences. Here
it cannot be denied, however, that Feilberg is the master. He was
a psychological botanist without equal, unsurpassed in
grasping and rendering those glimpses in which so much that we as a rule
pass by indifferent and thoughtless can come into being for us. But
he will rejoice in his grave that a worthy
successor has come, who by his own way has been led to discover
valuable features of life's involuntary unfolding and
to protect their significance.

\piecerule
% --- p. 154 ---
\setcounter{footnote}{0}%

\opage{154}\essayhead{LUDVIG HOLSTEIN: ``DEN GRØNNE MARK''}

This little book comes as a refreshing breath from
nature and human life, invigorating after all the shallow
talk one has been able to hear here at home in the
recent past, a talk in which neither the sound and free
human life nor the work of research has come into its
due. The book comes from the one of our poets who has most
lived in nature, who has filled our imagination with the finest
and freshest pictures of nature, and who now proves that he can
look at the mental life of man and at human history with
the same clear eyes as at the nature around us. We are in
high degree in need of such a poet, for we often go by, learned
and unlearned, uncomprehending, not only the life of nature around
us, but also the inner life in the minds of men.

Ludvig Holstein takes his stand on purely human ground.
All religion has for him ultimately sprung from human
souls and their experiences. Here is the green field, where
all the plants of the life of the spirit, the religious like all others, grow.
Here the religion of life grows forth, with its deep root in human
life-feeling, in which all profound experiences
are laid down. It is the poet's calling to interpret life, not least
where it stirs deepest. And faith itself is the greatest of
all poets.

In this last sentence the author's personal
standpoint is felt. In itself it does not say enough, for a scoffer could
also utter it. But here it is uttered with seriousness and fine
understanding for the involuntary unfolding of religious
feelings and thoughts under the struggle of life, as well as with an open mind
and good understanding for what the history of religion can teach
us. --- I may perhaps remind the reader that I have often, the first time in
% --- p. 155 ---
\setcounter{footnote}{0}%
\opage{155}my book on Kierkegaard (1893), called religion the
poetry of life. But --- alas! --- what good is it for a philosopher
to speak about poetry! A poet is needed. And now he has come
as a poet should, not with attempts to stop up
the holes of theology with lyric poetry, but with the fresh green
of life and nature about his brow. Let us then see how he
arranges his thoughts. I cannot, unfortunately, present
them with the rich poetry in which his speech appears,
although here, for the first time, he writes in prose. ---

Holstein demands ``studies in the field''. If we go the way of observation,
we easily discover how many different views
there are among men, even among those who count
themselves as of the same confession. Everyone creates his paradise
according to his need, and there is good use for the many mansions
Jesus spoke of. But only when we leave the human world
and go out into free nature do we stand before the mystery
that holds the key to all mysteries, first and
foremost to the myth of the mother and the child, the basic myth
of all religion, springing from the heart of nature itself: ``Every
growing blade of grass is a mystery, indeed Mystery itself. Like the
shepherds we stand face to face with the mystery of coming-to-be of All-Nature.''
Oehlenschläger already saw this in
``Jesu Liv i den aarlige Natur'': ``Each year, when the mists
flee away, then the little child Jesus is born again.'' --- And out here,
Holstein continues, on the green field Mary
and the Child shall remain, the human mother, the mystery of coming-to-be
itself, incomprehensible and eternal, reality and illusion in One, the primal myth,
from whose womb all religions have sprung! --- No
dualism! In human life as in nature the
great process of birth, death and renewal makes itself known.

As Holstein here involuntarily passes from nature to
history by reason of the inward connection of the two, so
one also detects his historical sense, in that the myths of the ancient
Oriental religions about the mother of the gods and her son
and the mysteries attached to them come to stand as predecessors
of the Christian conceptions. His historical
sense shows itself also in the remark that the church as
% --- p. 156 ---
\setcounter{footnote}{0}%
\opage{156}church reached its zenith in the thirteenth and fourteenth centuries.
Protestantism is for him a half-measure when it wishes to appear
with authority. It is consistent only when it builds
on freedom of conscience. Out on the green field
we believe in the natural growth of life, borne by will and led
by honest thinking. We sing hymns --- but not all church
hymns can we use. We see clearly that time goes a long
way, and we ourselves do not know to what building we are gathering
stones. But to peep behind the curtain we cannot. ---

Holstein is not the only one of our poets who has
taken part in the discussion of the religious question. Before
him Helge Rode has made his contribution, not so much as
a poet (though more than he wishes to be), but as a critic
of the position of the sciences in relation to the church. Holstein
reproaches Rode with holding fast to the Christian dualism
that is due to the influence of Paul's way of thinking, but objects
especially to him that his standpoint is ``occult rationalism''.
This last expression aptly characterizes a
conception that wishes to prove the possibility of miracles by referring
to the inventions of recent times, surprising to lay souls.
The way of thought of the old rationalistic pastors was
the same as Helge Rode's. --- I regard it as
exceedingly significant that one poet has here countered
another poet.

Holstein stands somewhat unclear toward the so-called Biblical
criticism (which ought to be called Biblical research, by analogy with
research on Homer and Plato). Sometimes he speaks of ``Biblical criticism''
in general, sometimes of ``a certain Biblical criticism'',
which is rejected with particular emphasis, and sometimes
of ``objective Biblical criticism''. There was every reason to
keep to this last alone, which is the properly scientific one,
and which therefore, like all true science, has a clear
eye for the difficulties that still remain. This
kind of research will agree with Holstein's remark that,
as regards the origin of Christianity, ``it will hardly ever
be possible to unravel all the threads''.

Since the discussion is of Biblical research, I will remark that
% --- p. 157 ---
\setcounter{footnote}{0}%
\opage{157}many scholars will not agree with Holstein in the
weight he places on the claim that the development of Christianity, especially
in the Pauline and Johannine direction, is supposed to be understood
from the connection between the Old and the New
Testament. Undoubtedly certain Oriental influences
have made themselves felt here, which it would not be hard for
Holstein to admit, since he lays so much weight on
the ``primal myth'' of the mother and the son. ---

As to future development, Holstein first
points out that a crisis has set in here in this country,
because the time is past in which ``the high renaissance of Danish
spiritual life, unchecked by the dogmas of the church, could flourish
under its roof'', so that ``the church, in spite of its dogmas,
was the mediator of its time between Christianity and
spiritual life. Since then spiritual life has gone its own ways, and the
Lutheran church stands as an anachronism'' without connection
with spiritual life as a whole. Yet Holstein does not wish
the church separated from the state, since that would breed rule by pastors
and narrow orthodoxy. What matters is whether the church will set the spirit
above the dogma. Holstein's last words are: Let the bells
ring and the dogmas sleep!

With this last word the real state of affairs
among our people is probably also expressed. People do not want to be rid of the bells,
they like to hear them ring and are moved by them, as by the liturgical part
of the church service generally. But people do not
want to hear too much about the dogmas. History must show
how long this spiritual condition can last. Both
those who stand outside the national church and those who stand inside
it will follow with suspense its further historical development
here. And all of them will have reason to be grateful
for the reference to life's green field that one of
our finest poets has given us in the present book.

\piecerule

% --- p. 158 ---
\setcounter{footnote}{0}%

\opage{158}\essayhead{SHAKESPEARE'S ``STORMEN''}

For all who admire Shakespeare and live in his works,
a great event is at hand, as the Royal
Theater is now [1926] to perform the drama that sets the crown on
his life's work. For my part I have found in it
an incentive to immerse myself anew in this wonderful
work, and have felt a desire to write down something of how
I understand it. I am not a Shakespeare scholar, only a
Shakespeare reader, and I must say in advance how
much I owe to the excellent English interpreters of Shakespeare,
especially Edward Dowden and A. C. Bradley. It is among
my best memories from my visits to England
that in the autumn of 1904 at the University of Oxford
I heard Bradley's lecture on Shakespeare as a man
(``Shakespeare the man''), which was later printed in his ``Oxford
Lectures''.

I shall try to give the content of the drama in brief outline.

Prospero, Duke of Milan, was by treachery driven from
his land by his own brother, to whom he had left
the government of the country while he himself immersed himself in his studies.
Now he lives with his daughter on an island far out in
the sea, still absorbed in his investigations and occupied with instructing
his daughter, Miranda, who in her tender age was torn
away from all human intercourse but her father's.
Prospero has, like Faust, immersed himself in magic, but his
life of thought has during his stay on the island won such
power that he can command the elements without needing, like Faust,
so questionable a helper as Mephistopheles.
Yet he has two serving beings, whom he found on
the island. The one, Ariel, is an ethereal being, who was held bound by a witch,
% --- p. 159 ---
\setcounter{footnote}{0}%
\opage{159}who formerly ruled on the island, in
harsh torment until Prospero freed him. Ariel moves
with the speed of thought, acts swiftly and wisely in his master's
spirit and can give him good counsel. Indeed, though he is an ethereal
being, without human flesh and blood, he can tell
Prospero what magnanimous humanity commands; he knows
how he would act if he were a man. In the midst
of all his zeal in his benefactor's service he longs, however, to
be released from this service in order again to become One with the purest
elements of nature. The other servant, Caliban, is a
cross between beast and devil, a son of the witch. Prospero
has tried in vain to awaken a human mind in
him. He cannot forget that he is the son of the island's former
mistress. He continues to be bestial and vindictive
and is therefore treated severely, especially after he had
tried to rape Miranda. The human language Prospero
has taught him he uses mostly to curse his
master. The only human trait in him is his
joy in hearing music. He too longs for freedom,
but it is the unruliness of raw elements he yearns for.
He is forced to toil for Prospero's material necessities,
while Ariel works for the carrying out of the master's
ideal plans.

The magic into which Prospero has been initiated through his investigations
he now uses to get his enemies into his power and
thereby deliverance and happiness for himself and his daughter. In a storm
raised by him his treacherous brother,
Antonio, and the latter's helper, King Alonzo of Naples, are stranded on
the island. Ariel takes care that no one loses his life. And now
the shipwrecked wander about in various groups on the island. A sweet
music leads the King's son Ferdinand to Prospero's cave,
where he and Miranda are at once captivated by each other. One of the
most graceful love relationships Shakespeare has portrayed unfolds
now. It causes Prospero the greatest joy: ``As
glad as they I cannot become, and yet nothing could give me
greater joy!'' Yet he puts the Prince to a hard trial
through strict labor before he gives his consent.

% --- p. 160 ---
\setcounter{footnote}{0}%
\opage{160}When the most important group of the shipwrecked, Antonio,
Alonzo and their courtiers, are known to be in Ariel's power,
Prospero, in his great joy at the young couple's happiness,
has Ariel conjure up a series of scenes in which goddesses bring
their good wishes to the pact, and nymphs and harvesters
dance. And then comes one of the most characteristic scenes
in the play. In the midst of the conjured spectacle
Prospero suddenly starts up and bids the spirits vanish. Suddenly
the thought had shot through him that all was not yet
safe. He had, absorbed in the visions
that conjured themselves up, forgotten that the monster Caliban out in the woods
had met two men from the ship, who had saved a good
supply of wine, and that now all three were on their way to attack
Prospero and seize his daughter, at whose side the
one of the conspirators, the King of Naples' butler, in
whom Caliban saw a revelation of his erstwhile god, then
was to rule over the island. No wonder that Prospero was dismayed
and upset when he came to think of what was
brewing against him, while he had begun to celebrate the
settlement in which his magnanimity toward his enemies was to
bring peace and joy. His daughter had never seen him
so embittered. But he collects himself and calms himself and the others
with the thought that not only such visions as they had just
had must dissolve into nothing, but that it would go
so with everything great and glorious, with our very earth.
Our life is indeed only like one of those dream-visions that just now
came forth.

It now dawns on Prospero that his trouble in educating
Caliban has been in vain, and that he must proceed severely
before he lets magnanimity rule. Ariel brings the brutal
conspirators into his power, they get their punishment, and when
Prospero later sails away from the island, he leaves Caliban behind.
The monster has not improved, but he has seen
what folly he committed in seeing a god in a drunken
butler. --- In a theater in London one saw at the end the ship
with Prospero and his enemies, who had now received forgiveness,
sail away from the shore, while Caliban lying on
% --- p. 161 ---
\setcounter{footnote}{0}%
\opage{161}a cliff edge follows them with his eyes. He has evidently
not understood a bit of what has happened. ---

Gradually all the other shipwrecked are now gathered at
Prospero's cave, led by Ariel, now under sweet music,
``the best comforter for an unquiet mind'', now under threatening
sounds. The King and Antonio awake from the magic sleep
into which Ariel has put them, and Prospero receives them with
an exhortation to forget what has been caused and has happened, and to
set it right as well as possible. The union of the King's
son and Miranda is to seal the reconciliation. If there
has been any doubt in Prospero as to how he
should treat those who are so wholly in his power, then Ariel
appears as his good genius. Ariel thinks that if
his master had seen these enemies' sufferings, he would
pity them; Ariel himself would, if he had been
a human being! Prospero, who as a human being can feel their wretchedness
even better than Ariel, follows his exhortation.
His enemies' wicked conduct has wounded him in his
innermost being, but magnanimity gains power over bitterness and
strength of spirit over desire for revenge. His feeling of his power
over men as over nature, the consciousness that
his plans have succeeded and that the vindication won
by his own strength is near, gives him the superiority
to go ``the rarer way'', that which leads beyond
the urge to repay evil with evil. He has learned that
the evil in nature and in human life cannot vanish
from the world, but he has at the same time learned that it can be overcome
with patience, without one's needing to betray the
noblest that is contained in one's own mind. ---

Shakespeare was not alone in regarding magnanimity as
the supreme trait of character. In the thinking of the Renaissance
this idea recurs again and again. In a series
of philosophers the awakening research into nature and the widened
view of human conditions led to the insight that
selfishness and desire for revenge are a limitation of the mind's
power, in the end a weakness. This line of thought is later completed
in Spinoza's doctrine of the free man, who lives out of
% --- p. 162 ---
\setcounter{footnote}{0}%
\opage{162}himself, and whose power is able not merely to assert itself,
but to show its superiority by helping others. Generositas
as an expression of spiritual strength is a main thought in
the ethics of the Renaissance. It is not my opinion that Shakespeare
should have been influenced by the philosophers of his time, but through
his experiences and through the work in his art
he has come by his own way to the same result as the
thinkers of his time, and he has been able to express it in a
work that no one but he could create. It is faith in
life that leads beyond sorrow as beyond hatred. When King Alonzo,
despite the reconciliation, cannot forget what has happened, Prospero
stops him with the words: ``Let us not burden our
memory with a sorrow whose time is past!'' --- The line of thought
that had let Prospero think with calm of the fact that
not only the paradises of fancy, but all human work,
indeed the earth itself on which we live, is subject to dissolution,
leads him to recognize that also the sorrow that has been caused
him, and the wrong that has been done to him, are subject to this law.

The Shakespeare scholars are agreed that the spiritual loftiness
that is portrayed in ``Stormen'' and has peace and reconciliation as
its consequence is on the whole characteristic of the last period of
Shakespeare's work as a poet. The same mood prevails
in ``Cymbeline'' and in ``A Winter's Tale''. Despite all crime
and wrong, fidelity and love triumph. Imogen
and Hermione are natures like Prospero, although they command no
other magic than patient love.
All these dramas mark a peculiar period in
Shakespeare's development, which, by the way, is characterized not only
by the superior mildness, but also by the quite
special grace with which young women (Miranda,
Imogen, Perdita) are portrayed. What is peculiar to
``Stormen'' in relation to the other dramas from the same time
is that Shakespeare through Prospero directly and more distinctly
than in other dramas expresses the conception of life he has
won, and that he now announces that he intends to cease with his
% --- p. 163 ---
\setcounter{footnote}{0}%
\opage{163}poetic works. He takes leave of his hearers by
having Prospero break his magic staff, bury it deep in the earth,
and sink his magic book into the sea. The oppositions of life
have now faded away for him. He has conducted his
spiritual struggle with the experiences that fell
to his lot, whether they intervened in his own personal
life or not. He has fought as a man and as an
artist, and ``Stormen'' shows that he has won the victory in
both capacities. Now he returns to Stratford on
Avon, where with the practical sense that always seems to
have worked alongside his artistic urge, he had
prepared for himself a fine home. He who, according to the legend, had to
flee as a young man from his native town now owned one of the
most stately houses there, and as Ariel looked forward to his freedom
in the ether, Shakespeare has looked forward to walking again
along the banks of the Avon. An excellent Shakespeare scholar has also
noticed in Shakespeare's later production the fragrance of
the meadows by the Avon. (For my part, conversely,
a walk across these meadows has made the impression of Shakespeare's
little pictures of nature fresher and more vivid than
the mere reading could do.)

The Shakespeare whom we meet in the last period of his
work as a poet, and most markedly in ``Stormen'', must be seen
in connection with the one who earlier, about forty years
old, wrote Hamlet, Othello, Macbeth and Lear, those
mighty tragedies in which the oppositions of human life
come forth so sharply and ruthlessly, and in which downfall
often seems to be the last word. As Bradley has
shown, these tragedies deal for a substantial part with the
disappointments and adversity that open and ardent natures can
meet on their way. Such a nature was Hamlet before his father's
murder and his mother's hasty wedding, before his friends'
treachery and Ophelia's spying, such a one was also
Othello, before Iago devilishly awakened his doubt, and Lear, before
his daughters' ingratitude showed itself. Such a one was moreover
Macbeth before the witch sisters' prophecy and his wife's
% --- p. 164 ---
\setcounter{footnote}{0}%
\opage{164}incitement. And Bradley has aptly stressed that Shakespeare
himself is characterized by his contemporaries as an open, free
and fine nature. He must have had experiences of what
life can bring of misfortune and wickedness, even if it did not
directly intervene in his own fate. His imagination
was in any case seized by it, and he did not
rest until he had portrayed in great and strong strokes what he
saw. But he did not lose faith in the reality of the good.
He maintains it in the very fall of the great and noble, and he
portrays it as fidelity and patience in figures
such as Kent and Horatio. He stands in the midst of the strife and lets
the oppositions rebound against one another. The same spiritual
power that in ``Stormen'' carried him beyond life's oppositions
had made it possible for him in those great tragedies
to let these oppositions confront one another
in all their strength. It is fear and pity that come
to rule, not the transfigured and inward understanding
and the patient endurance that hold sway in ``Stormen''.
In ``King Lear'' too there is talk of downfall.
When the old King enters with Cordelia's corpse, Kent
cries: ``Is this the end of all that was foretold?'', --- Edgar:
``or an image of the horror that can be expected?'' --- and the Duke of Albany:
``Downfall and end?'' (V, 3, 263ff.). But
now, in ``Stormen'', it is with a moved mind, but not in
horror, that Prospero thinks of the downfall. A transfigured
mood now prevails in the poet's mind. We would not
rejoice in ``Stormen'' if we had not, by the same poet's
hand, lived through the tragic moods that in his early manhood
he felt the need to give expression to, and
there would not have been the great weight in the poet's
last proclamation if he had not been in
purgatory before he walked in paradise.

It is not possible without more to make use of the yield of others'
experience of life, so that one holds to the transfigured
form in which their conception of life can appear, after
they have tried life. We must all begin afresh. And however pure
% --- p. 165 ---
\setcounter{footnote}{0}%
\opage{165}and clear the atmosphere is into which we are led by Prospero and
Ariel, there will yet be those who, when they need
support in life's struggle, rather take refuge in ``King Lear'' than
in ``Stormen''. Even if the noble and great is consecrated to downfall,
it is yet uplifting to know that it has been and can
be in the world. And it can be strengthening to take a bout
with the thoughts that are awakened by the fate of Lear and Cordelia,
before one tries to live in the ethereal world of ``Stormen''. Great
it is to have given us a work like ``Stormen'', but
greater that this work was created only after it had been
lived, felt and thought in the world of strife and oppositions.
And this he has done whom Goethe called ``William, Stern
der höchsten Höhe''. He does not fail us, neither during
the struggle nor in defeat, nor when victory is won.

\piecerule

\end{document}
